% ============================================================
%  Causal-Plasticity Atlas — Sections 1–3
%  paper_sections_1_3.tex
% ============================================================

\documentclass[11pt,a4paper]{article}

% ---- Packages ------------------------------------------------
\usepackage[utf8]{inputenc}
\usepackage[T1]{fontenc}
\usepackage{amsmath,amssymb,amsthm,mathtools}
\usepackage{thmtools}
\usepackage{bm}
\usepackage{bbm}
\usepackage{microtype}
\usepackage{hyperref}
\usepackage{cleveref}
\usepackage{natbib}
\usepackage{booktabs}
\usepackage{enumitem}
\usepackage[margin=1in]{geometry}
\usepackage{xcolor}

% ---- Theorem environments ------------------------------------
\newtheorem{theorem}{Theorem}[section]
\newtheorem{lemma}[theorem]{Lemma}
\newtheorem{proposition}[theorem]{Proposition}
\newtheorem{corollary}[theorem]{Corollary}
\theoremstyle{definition}
\newtheorem{definition}[theorem]{Definition}
\newtheorem{example}[theorem]{Example}
\newtheorem{remark}[theorem]{Remark}
\newtheorem{assumption}{Assumption}

% ---- Notation macros -----------------------------------------
\newcommand{\cC}{\mathcal{C}}                 % context space
\newcommand{\cM}{\mathcal{M}}                 % SCM
\newcommand{\cF}{\mathcal{F}}                 % structural equations
\newcommand{\cA}{\mathcal{A}}                 % QD archive
\newcommand{\cE}{\mathcal{E}}                 % atlas entries / emergence set
\newcommand{\cB}{\mathcal{B}}                 % tessellation bins
\newcommand{\cR}{\mathcal{R}}                 % robustness certificate
\newcommand{\cD}{\mathcal{D}}                 % data
\newcommand{\Vbar}{\bar{V}}                   % universal variable pool
\newcommand{\Pa}{\mathrm{Pa}}                 % parent set
\newcommand{\Dom}{\mathrm{dom}}               % domain
\newcommand{\Img}{\mathrm{im}}                % image
\newcommand{\SHD}{\mathrm{SHD}}               % structural Hamming distance
\newcommand{\MMD}{\mathrm{MMD}}               % maximum mean discrepancy
\newcommand{\JSD}{\mathrm{JSD}}               % Jensen-Shannon distance
\newcommand{\JS}{\mathrm{JS}}                 % Jensen-Shannon divergence
\newcommand{\KL}{D_{\mathrm{KL}}}             % KL divergence
\newcommand{\Cost}{\mathrm{Cost}}             % alignment cost
\newcommand{\dDAG}{d_{\mathrm{DAG}}}          % DAG alignment distance
\newcommand{\dmu}{d_{\mu}}                    % mechanism distance
\newcommand{\dC}{d_{\cC}}                     % context distance
\newcommand{\Cobs}{\cC_{\mathrm{obs}}}        % observed contexts
\newcommand{\psi}{\bm{\psi}}                  % plasticity descriptor
\newcommand{\R}{\mathbb{R}}                   % reals
\newcommand{\N}{\mathbb{N}}                   % naturals
\newcommand{\indicator}{\mathbbm{1}}          % indicator function
\newcommand{\eps}{\varepsilon}                % epsilon

% ---- Title ---------------------------------------------------
\title{%
  Causal-Plasticity Atlas:\\
  Mapping Mechanism Invariance, Plasticity, and Emergence\\
  Across Heterogeneous Observational Contexts%
}
\author{%
  [Authors redacted for review]
}
\date{}

% ==============================================================
\begin{document}
\maketitle

\begin{abstract}
Given heterogeneous multi-context observational datasets, we introduce the
\emph{Causal-Plasticity Atlas} (CPA), a framework that systematically
classifies causal mechanisms as \emph{invariant}, \emph{parametrically
plastic}, \emph{structurally plastic}, or \emph{emergent} across contexts
and organises the resulting classifications into a navigable
quality-diversity archive equipped with per-mechanism robustness
certificates.  The key technical innovation is a proper metric on the space
of context-indexed structural causal models, constructed from the
Jensen--Shannon distance between conditional mechanism distributions, which
we lift to a four-dimensional \emph{plasticity descriptor} that decomposes
mechanism change into parametric, structural, emergence, and
context-sensitivity axes.  We formalise tipping-point detection as
discontinuity detection in a plasticity functional over context paths,
prove metric properties for our DAG alignment distance, and establish
sample-complexity and classification-correctness guarantees.  The
curiosity-driven quality-diversity search over the plasticity-descriptor
space provably converges in coverage, yielding a complete atlas of
mechanism-change patterns.
\end{abstract}

% ==============================================================
\section{Introduction}\label{sec:intro}
% ==============================================================

\subsection{Motivation: Why Mechanism Plasticity Matters}

Modern data science increasingly confronts \emph{multi-context} settings:
the same scientific question is studied across tissues, countries,
experimental conditions, or time periods, and the data-generating processes
can differ in subtle or dramatic ways from one context to the next.
A gene-regulatory network may rewire between cell types; an economic
mechanism linking education to income may shift across decades; a clinical
pathway active in one patient cohort may be entirely absent in another.
Understanding \emph{how} and \emph{where} causal mechanisms change---not
merely \emph{whether} they change---is essential for safe generalisation,
policy transfer, and scientific discovery.

The dominant paradigm for multi-environment causal reasoning is
\emph{Invariant Causal Prediction} (ICP)
\citep{peters2016causal,heinzedeml2018invariant}, which identifies
predictors whose causal effect on a target is \emph{invariant} across
observed environments.  ICP and its extensions
\citep{pfister2019invariant,rothenhausler2021anchor} have been remarkably
successful, but they make a fundamental binary distinction: a mechanism is
either invariant or it is not.  Once a mechanism is labelled
``non-invariant,'' the framework offers no further characterisation.  Yet
non-invariance is not monolithic.  A mechanism may change smoothly with
context (parametric plasticity); the causal graph itself may restructure
(structural plasticity); or entirely new variables may emerge in some
contexts but not others (emergence).  These qualitatively different modes
of change demand qualitatively different responses from the analyst:
parametric plasticity may be addressed by context-specific parameter
estimation, structural plasticity requires context-specific model
selection, and emergence calls for fundamentally new causal hypotheses.

Furthermore, mechanism change is not merely a nuisance---it is
\emph{informative}.  Tipping points, where a mechanism abruptly
transitions from one regime to another along a context gradient, are among
the most scientifically valuable findings in multi-context studies.
Identifying the temperature at which a metabolic pathway rewires, or the
policy threshold at which an economic mechanism breaks down, can guide
experimental design and policy intervention.  Yet existing frameworks
provide no formal apparatus for detecting or certifying such transitions.

\subsection{Gaps in the Existing Literature}

The limitations of the current landscape can be summarised along five axes:
\begin{enumerate}[label=(\roman*)]
  \item \textbf{Binary classification.}
    ICP \citep{peters2016causal} and its descendants classify mechanisms as
    invariant or non-invariant.  The rich interior of the non-invariant
    category---parametric vs.\ structural change, gradual vs.\ abrupt
    transitions---remains uncharacterised.

  \item \textbf{Fixed variable sets.}
    ICP, JCI \citep{mooij2020joint}, and CD-NOD \citep{huang2020causal}
    all assume a fixed set of observed variables across environments.  In
    practice, variables appear and disappear: a gene is expressed in one
    tissue but silenced in another; an economic indicator is measured in
    some countries but not others.  No existing multi-environment causal
    framework handles \emph{variable-set mismatch} as a first-class
    phenomenon.

  \item \textbf{No mechanism-level metric.}
    Structural Hamming Distance (SHD) and Structural Intervention Distance
    (SID) \citep{peters2015structural} compare DAGs on the same vertex set
    but ignore the parametric content of mechanisms.  Graph edit distance
    (GED) \citep{gao2010survey,bougleux2017graph} handles node
    insertions/deletions but does not incorporate distributional
    information.  There is no unified metric that integrates structural,
    parametric, and variable-set differences.

  \item \textbf{No robustness certificates.}
    ICP provides $p$-values for invariance; stability selection
    \citep{meinshausen2010stability} provides variable-selection
    guarantees.  But no existing method provides per-mechanism, calibrated
    certificates guaranteeing that a plasticity classification is correct
    with stated probability, with an explicit margin quantifying distance
    to the nearest classification boundary.

  \item \textbf{No systematic exploration.}
    Exhaustive pairwise comparison of mechanisms across contexts scales
    quadratically and offers no guarantee of discovering rare or unexpected
    mechanism-change patterns.  Quality-diversity (QD) search
    \citep{mouret2015illuminating,vassiliades2018using} has proven
    effective for exploring complex behavioural spaces in robotics and
    evolutionary computation, but has never been applied to causal
    mechanism spaces.
\end{enumerate}

\subsection{Our Contribution: The Causal-Plasticity Atlas}

We introduce the \emph{Causal-Plasticity Atlas} (CPA), a framework that
addresses all five gaps through the following novel constructions:

\begin{enumerate}[label=\textbf{C\arabic*.}]
  \item \textbf{Multi-Context Causal Model with variable-set mismatch
    (MCCM).}
    We define the first formal framework in which the variable set~$V_c$
    varies across contexts, with \emph{emergence} (a variable exists in
    context~$c$ but not~$c'$) and \emph{disappearance} formalised as
    first-class citizens.  An information-theoretic criterion
    distinguishes genuinely new variables from recodings or renamings
    (\Cref{def:emergence}).

  \item \textbf{Four-dimensional plasticity descriptors.}
    We replace the binary invariant/non-invariant label with a continuous
    four-dimensional descriptor $\psi_i \in \R^4$ that decomposes
    mechanism change into \emph{parametric plasticity}, \emph{structural
    plasticity rate}, \emph{emergence rate}, and
    \emph{context sensitivity} (\Cref{def:plasticity-descriptor}).  The
    descriptor uses the Jensen--Shannon distance---a proper metric---as the
    base mechanism distance, ensuring the descriptor space has well-defined
    geometric structure.

  \item \textbf{Context-aware DAG alignment with metric properties.}
    We define a DAG alignment operator that simultaneously handles
    structural differences, parametric mechanism differences (via JSD), and
    variable-set mismatch in a single weighted cost function.  The optimal
    alignment cost $\dDAG$ is a pseudo-metric on the space of DAGs over
    subsets of a universal variable pool (\Cref{thm:dag-metric}).

  \item \textbf{Per-mechanism robustness certificates.}
    Each plasticity classification is accompanied by a certificate
    $\cR_i = (\kappa_i, \delta_i, n_{\min}, \alpha)$ guaranteeing
    correctness with probability $\ge 1 - \alpha$ whenever the margin
    $\delta_i$ exceeds an explicit concentration bound
    (\Cref{def:robustness-certificate}).

  \item \textbf{Curiosity-driven QD search with tipping-point detection.}
    We apply curiosity-driven MAP-Elites over the four-dimensional
    plasticity-descriptor space, with a novel curiosity signal combining
    archive novelty and quality improvement (\Cref{def:curiosity}).
    Tipping points---abrupt mechanism transitions along context
    paths---are formalised as discontinuities in a plasticity functional
    (\Cref{def:tipping-point}), with provably consistent detection.
\end{enumerate}

The central intellectual contribution is recognising that the space of
\emph{``how a mechanism changes across contexts''} is itself a rich
geometric object---not a binary label---and that information-theoretic
descriptors of this space enable both rigorous statistical guarantees and
curiosity-driven exploration that discovers mechanism-change patterns a
human analyst would miss.

\subsection{Paper Organisation}

The remainder of this paper is organised as follows.
\Cref{sec:prelim} reviews the necessary background on structural causal
models, multi-environment causal inference, quality-diversity optimisation,
changepoint detection, and graph comparison.
\Cref{sec:framework} presents the formal framework: fourteen definitions
spanning the MCCM, mechanism distance, DAG alignment, plasticity
descriptors, the QD archive, curiosity signals, tipping points, and
robustness certificates, together with six assumptions and the first main
theorem (DAG alignment metric properties) with proof.
Section~4 develops the algorithmic pipeline.
Section~5 presents theoretical results on classification correctness,
sample complexity, certificate soundness, tipping-point consistency, and
atlas completeness.
Section~6 provides experimental evaluation on synthetic, semi-synthetic,
and real-world datasets.
Section~7 discusses limitations, extensions, and open problems.

% ==============================================================
\section{Preliminaries and Related Work}\label{sec:prelim}
% ==============================================================

We review the five bodies of literature on which the Causal-Plasticity
Atlas builds, highlighting for each what we inherit and what we extend.

\subsection{Structural Causal Models and Causal Discovery}

A \emph{structural causal model} (SCM)
\citep{pearl2009causality,spirtes2000causation} specifies the
data-generating process for a collection of observed variables via
structural equations $X_i := f_i(\Pa_i, U_i)$, where $\Pa_i$ is the
parent set of~$X_i$ and~$U_i$ is exogenous noise.  The induced directed
acyclic graph (DAG)~$G$ encodes the causal structure.  Under the
\emph{causal Markov condition}, each variable is conditionally independent
of its non-descendants given its parents; under \emph{faithfulness}, every
observed conditional independence is entailed by the DAG.

\emph{Constraint-based} methods such as the PC algorithm
\citep{spirtes2000causation} and FCI
\citep{spirtes2000causation,zhang2008completeness} discover the DAG (or a
Markov equivalence class thereof) from conditional independence tests.
\emph{Score-based} methods such as GES \citep{chickering2002optimal}
search the space of DAGs by greedily optimising a penalised likelihood
score.  Hybrid methods combine elements of both
\citep{tsamardinos2006max}.  In all cases, the output from purely
observational data is a \emph{completed partially directed acyclic graph}
(CPDAG), representing the Markov equivalence class.

The CPA takes per-context causal structures as input---estimated by any
of these methods---and analyses their variation across contexts.  We
inherit the Markov condition (Assumption~A1) and faithfulness
(Assumption~A2) as per-context requirements, and we state all definitions
with respect to CPDAGs so that the framework is grounded in estimable
quantities rather than unobservable true DAGs.

\subsection{Multi-Environment Causal Inference}

\emph{Invariant Causal Prediction} (ICP) \citep{peters2016causal}
identifies causal parents of a target variable by testing which predictors
have \emph{invariant} conditional distributions across environments.  The
key insight is that causal mechanisms, unlike spurious associations, are
stable under distribution shift.  Extensions include nonlinear ICP
\citep{heinzedeml2018invariant}, anchor regression
\citep{rothenhausler2021anchor} for worst-case prediction, and
invariance-based structure learning \citep{pfister2019invariant}.

\emph{Joint Causal Inference} (JCI) \citep{mooij2020joint} augments the
causal graph with system variables representing environments, learning a
single DAG with context-dependent edge annotations.  \emph{CD-NOD}
\citep{huang2020causal} detects which causal modules change across
heterogeneous domains but does not classify the \emph{type} of change.
The transportability framework
\citep{bareinboim2014transportability,pearl2014external} uses selection
diagrams to reason about which causal effects transfer across populations,
but assumes a common causal graph with varying parameters, not
variable-set differences.

All of these frameworks assume a \emph{fixed variable set} across
environments.  The CPA extends this body of work in three directions: (i)
it allows the variable set to vary, formalising emergence and
disappearance; (ii) it replaces binary invariance labels with a continuous
plasticity spectrum; and (iii) it provides per-mechanism robustness
certificates with explicit margins.

\subsection{Quality-Diversity Optimisation}

Quality-diversity (QD) algorithms search for a \emph{diverse collection}
of high-performing solutions, rather than a single optimum.  \emph{MAP-Elites}
\citep{mouret2015illuminating} maintains an archive of solutions indexed
by a \emph{behaviour descriptor}, keeping the highest-quality solution in
each cell of a tessellation of the behaviour space.
\emph{CVT-MAP-Elites} \citep{vassiliades2018using} replaces the
hand-designed grid with a centroidal Voronoi tessellation (CVT), yielding
more uniform coverage of high-dimensional behaviour spaces.

QD methods have been applied successfully in robotics
\citep{cully2015robots}, game design \citep{gravina2019procedural}, and
evolutionary computation, but have not previously been used in the context
of causal inference.  The CPA adapts CVT-MAP-Elites by using the
four-dimensional plasticity descriptor (\Cref{def:plasticity-descriptor})
as the behaviour characterisation, with a quality function that rewards
context coverage, estimation confidence, and robustness.  The curiosity
signal (\Cref{def:curiosity}) extends standard novelty search
\citep{lehman2011abandoning} with a quality-improvement component.

\subsection{Changepoint Detection}

Changepoint detection identifies abrupt shifts in sequential data.  The
\emph{Pruned Exact Linear Time} (PELT) algorithm \citep{killick2012optimal}
uses dynamic programming with pruning to find optimal changepoints under a
penalised cost.  \emph{Bayesian Online Changepoint Detection} (BOCPD)
\citep{adams2007bayesian} maintains a posterior over the current run
length, providing online changepoint probabilities.  Both methods are
well-studied with known consistency guarantees.

The CPA adapts PELT for \emph{tipping-point detection} along context
paths (\Cref{def:tipping-point}).  The key novelty is the cost function:
whereas standard PELT operates on a univariate signal, our cost measures
within-segment mechanism homogeneity using the JSD between conditional
distributions.  Consistency of tipping-point detection
(\Cref{sec:framework}) follows from combining PELT's consistency with
uniform convergence of the mechanism-distance estimator.

\subsection{Graph Comparison and Alignment}

\emph{Structural Hamming Distance} (SHD) counts edge disagreements between
two DAGs on the same vertex set.  \emph{Structural Intervention Distance}
(SID) \citep{peters2015structural} measures disagreement in
interventional distributions.  \emph{Graph edit distance} (GED)
\citep{gao2010survey,bougleux2017graph} generalises to graphs with
different vertex sets by assigning costs to node/edge insertions,
deletions, and substitutions; it is known to satisfy metric properties
under appropriate cost functions.

Our DAG alignment operator (\Cref{def:dag-alignment}) extends GED in two
directions: (i) node-substitution costs incorporate the JSD between
conditional mechanism distributions---a quantity with no analogue in
standard GED; and (ii) the consistency constraint ensures that alignments
respect the causal semantics of edges rather than treating the graph as an
unlabelled combinatorial object.  The resulting alignment distance $\dDAG$
is a pseudo-metric (\Cref{thm:dag-metric}), providing the geometric
foundation for the CVT tessellation of the QD archive.

\subsection{Summary: What We Build On vs.\ What Is New}

\Cref{tab:novelty} summarises the relationship between the CPA and prior
work.  The left column lists known results and methods that we inherit; the
right column lists our novel constructions.

\begin{table}[ht]
\centering
\caption{Known foundations vs.\ novel contributions.}\label{tab:novelty}
\smallskip
\begin{tabular}{@{}p{0.46\textwidth}p{0.46\textwidth}@{}}
\toprule
\textbf{Known \& Reused} & \textbf{Novel in This Work} \\
\midrule
SCM framework \citep{pearl2009causality} &
  Multi-Context Causal Model with variable-set mismatch \\
PC/GES/FCI for per-context structure learning &
  DAG alignment operator with JSD mechanism costs \\
ICP binary invariance testing \citep{peters2016causal} &
  Four-dimensional plasticity descriptor \\
$\sqrt{\JS}$ as a metric \citep{endres2003new} &
  Plasticity descriptor using JSD as base distance \\
MAP-Elites / CVT-MAP-Elites \citep{mouret2015illuminating,vassiliades2018using} &
  QD archive over plasticity-descriptor space \\
PELT changepoint detection \citep{killick2012optimal} &
  Tipping-point detection on mechanism divergence sequences \\
Bootstrap confidence intervals \citep{friedman1999data} &
  Per-mechanism robustness certificates with explicit margins \\
GED metric properties \citep{gao2010survey} &
  DAG alignment pseudo-metric integrating structural, parametric, and variable-set costs \\
\bottomrule
\end{tabular}
\end{table}


% ==============================================================
\section{Formal Framework}\label{sec:framework}
% ==============================================================

We now present the formal framework underlying the Causal-Plasticity
Atlas.  We begin with core definitions (\Cref{sec:core-defs}), state
our assumptions (\Cref{sec:assumptions}), and prove the first main
theorem (\Cref{sec:thm-metric}).

% ---- 3.1 Core Definitions ------------------------------------
\subsection{Core Definitions}\label{sec:core-defs}

\begin{definition}[Structural Causal Model]\label{def:scm}
A \emph{structural causal model} (SCM) is a tuple
$\cM = (V, U, \cF, P_U)$ where:
\begin{itemize}[nosep]
  \item $V = \{X_1, \ldots, X_d\}$ is a finite set of endogenous
        (observed) variables,
  \item $U = \{U_1, \ldots, U_m\}$ is a finite set of exogenous (latent)
        variables,
  \item $\cF = \{f_i : i \in [d]\}$ is a set of structural equations
        $X_i := f_i(\Pa_i, U_{S_i})$, where
        $\Pa_i \subseteq V \setminus \{X_i\}$ is the \emph{parent set}
        and $S_i \subseteq [m]$ indexes the relevant exogenous variables,
  \item $P_U$ is a joint distribution over~$U$.
\end{itemize}
The \emph{induced DAG} $G(\cM)$ has vertex set~$V$ and edge set
$\{(X_j, X_i) : X_j \in \Pa_i\}$.  The \emph{observational distribution}
$P_V^{\cM}$ is the marginal of $P_{V,U}$ onto~$V$.
\end{definition}

\begin{definition}[Context and Context Space]\label{def:context}
A \emph{context} is an element $c \in \cC$ drawn from a \emph{context
space} $(\cC, \dC)$, where $\cC$ is a measurable set and
$\dC \colon \cC \times \cC \to [0, \infty)$ is a pseudo-metric:
\begin{enumerate}[nosep,label=(\roman*)]
  \item $\dC(c, c) = 0$ for all $c \in \cC$;
  \item $\dC(c, c') = \dC(c', c)$ for all $c, c' \in \cC$  (symmetry);
  \item $\dC(c, c'') \le \dC(c, c') + \dC(c', c'')$ for all
        $c, c', c'' \in \cC$ (triangle inequality).
\end{enumerate}
We allow $\dC(c, c') = 0$ for distinct $c \ne c'$, since different
experimental conditions may produce identical causal structures.  When
contexts are described by feature vectors $\phi(c) \in \R^p$, we set
$\dC(c, c') = \|\phi(c) - \phi(c')\|_2$.  When contexts are unstructured
labels, $\dC$ is derived from data (\Cref{def:data-context-distance}).
\end{definition}

\begin{definition}[Multi-Context Causal Model]\label{def:mccm}
A \emph{multi-context causal model} (MCCM) is a family
$\{\cM_c\}_{c \in \cC}$ of SCMs indexed by a context space $(\cC, \dC)$,
subject to the \emph{structural coherence constraint}: there exists a
\emph{universal variable pool} $\Vbar$ such that $V_c \subseteq \Vbar$ for
every $c \in \cC$.  The variable set $V_c$ may vary across contexts.

We write $G_c = G(\cM_c)$ for the induced DAG on~$V_c$ (identified only
up to its Markov equivalence class, i.e., the CPDAG $\hat{G}_c$, from
observational data) and $P_c = P_{V_c}^{\cM_c}$ for the observational
distribution.  In practice, we observe i.i.d.\ samples
$\cD_c = \{x_c^{(1)}, \ldots, x_c^{(n_c)}\}$ from~$P_c$ for a finite
subset $\Cobs \subset \cC$ of contexts, with potentially different sample
sizes~$n_c$ and variable sets~$V_c$.
\end{definition}

\begin{remark}
Unlike ICP \citep{peters2016causal}, JCI \citep{mooij2020joint}, and
CD-NOD \citep{huang2020causal}, which all assume a fixed variable set
across environments, the MCCM allows both the DAG structure and the
variable set to change across contexts.  This is what enables emergence
detection (\Cref{def:emergence}).
\end{remark}

\begin{definition}[Mechanism and Mechanism Identity]\label{def:mechanism}
A \emph{mechanism} for variable $X_i$ in context~$c$ is the conditional
distribution
\[
  \mu_i^c \;=\; P_c(X_i \mid \Pa_i^c)
\]
induced by the structural equation $f_i^c$ and the exogenous distribution
$P_{U_{S_i}}^c$, where $\Pa_i^c$ denotes the \emph{definite parents}
of~$X_i$ in the estimated CPDAG~$\hat{G}_c$.

Two mechanisms $\mu_i^c$ and $\mu_j^{c'}$ are \emph{identifiable as the
same mechanism} if:
\begin{enumerate}[nosep,label=(\arabic*)]
  \item \textbf{Variable correspondence:} there exists a bijection
        $\sigma \colon \{X_i\} \cup \Pa_i^c \to \{X_j\} \cup \Pa_j^{c'}$
        with $\sigma(X_i) = X_j$;
  \item \textbf{Distributional equivalence:} under the correspondence
        $\sigma$, the conditional distributions are equal:
        $P_c(X_i \mid \Pa_i^c)
         = P_{c'}(X_j \mid \sigma(\Pa_j^{c'}))$.
\end{enumerate}
In practice, exact equality is never tested; instead we use mechanism
distance (\Cref{def:mechanism-distance}).
\end{definition}

\begin{definition}[Mechanism Distance]\label{def:mechanism-distance}
Given two mechanisms $\mu_i^c$ and $\mu_j^{c'}$ with a variable
correspondence~$\sigma$ (\Cref{def:mechanism}), the \emph{mechanism
distance} is
\[
  \dmu(\mu_i^c, \mu_j^{c'};\, \sigma)
  \;=\;
  \JSD\!\bigl(P_c(X_i \mid \Pa_i^c),\;
              P_{c'}(X_j \mid \sigma(\Pa_j^{c'}))\bigr),
\]
where the \emph{Jensen--Shannon distance} is defined as the square root of
the Jensen--Shannon divergence:
\[
  \JSD(P, Q)
  \;=\;
  \sqrt{\JS(P \| Q)}
  \;=\;
  \sqrt{\tfrac{1}{2}\KL(P \| M) + \tfrac{1}{2}\KL(Q \| M)},
  \qquad
  M = \tfrac{1}{2}(P + Q).
\]
For conditional distributions, we compute the expected JSD over a
reference distribution on the conditioning variables:
$\dmu = \mathbb{E}_{Z \sim P_{\mathrm{ref}}}[\JSD(P(\cdot \mid Z),\, Q(\cdot \mid Z))]$,
where $P_{\mathrm{ref}}$ is the pooled empirical distribution of the
parent values across both contexts.  When parent sets have different
cardinalities under no valid~$\sigma$, we set $\dmu = +\infty$.
\end{definition}

\begin{remark}
The Jensen--Shannon distance satisfies the triangle inequality
\citep{endres2003new}, is always finite (unlike the KL divergence), is
symmetric, and applies to both discrete and continuous distributions
without requiring absolute continuity.  These properties are essential for
\Cref{thm:dag-metric}.
\end{remark}

\begin{definition}[DAG Alignment Operator]\label{def:dag-alignment}
Let $G_c = (V_c, E_c)$ and $G_{c'} = (V_{c'}, E_{c'})$ be two CPDAGs
from an MCCM.  A \emph{DAG alignment} is a partial injective map
$\alpha \colon V_c \rightharpoonup V_{c'}$ satisfying:
\begin{enumerate}[nosep,label=(\arabic*)]
  \item \textbf{Consistency:} if $X_i, X_j \in \Dom(\alpha)$ and
        $(X_i, X_j) \in E_c$, then either
        $(\alpha(X_i), \alpha(X_j)) \in E_{c'}$ (edge preserved) or the
        edge is classified as \emph{structurally plastic}.
  \item \textbf{Maximality:} $\Dom(\alpha)$ is maximal among all
        consistency-preserving partial injections, subject to a cost bound.
\end{enumerate}
The \emph{alignment cost} under~$\alpha$ is
\begin{equation}\label{eq:alignment-cost}
  \Cost(\alpha;\, G_c, G_{c'})
  \;=\;
  \lambda_{\mathrm{s}} \cdot |\Delta E(\alpha)|
  \;+\;
  \lambda_{\mathrm{m}} \cdot
    \sum_{X_i \in \Dom(\alpha)} \dmu(\mu_i^c, \mu_{\alpha(i)}^{c'};\, \alpha)
  \;+\;
  \lambda_{\mathrm{v}} \cdot
    \bigl(|V_c \setminus \Dom(\alpha)| + |V_{c'} \setminus \Img(\alpha)|\bigr),
\end{equation}
where $\Delta E(\alpha)$ is the set of edge disagreements under~$\alpha$,
and $\lambda_{\mathrm{s}}, \lambda_{\mathrm{m}}, \lambda_{\mathrm{v}} > 0$
are weights for structural, mechanism, and variable-mismatch costs,
respectively.

The \emph{DAG alignment distance} is
\[
  \dDAG(G_c, G_{c'})
  \;=\;
  \min_{\alpha} \Cost(\alpha;\, G_c, G_{c'}).
\]
\end{definition}

\begin{definition}[Plasticity Descriptor]\label{def:plasticity-descriptor}
For a mechanism $\mu_i$ (targeting variable~$X_i$) observed across a set
of contexts $\Cobs$, the \emph{plasticity descriptor} is a vector
$\psi_i \in \R^4$ with components:
\begin{enumerate}[nosep,label=(\arabic*)]
  \item \textbf{Mean mechanism distance} (parametric plasticity):
    \[
      \psi_i^{(1)}
      \;=\;
      \frac{2}{|C_i|(|C_i|-1)}
      \sum_{\substack{c < c' \\ c, c' \in C_i}}
      \dmu(\mu_i^c, \mu_i^{c'}),
    \]
    where $C_i = \{c \in \Cobs : X_i \in V_c\}$ is the set of contexts in
    which $X_i$ is present and the sum runs over unordered pairs.

  \item \textbf{Structural plasticity rate}:
    \[
      \psi_i^{(2)}
      \;=\;
      \frac{2}{|C_i|(|C_i|-1)}
      \bigl|\{(c, c') : c < c',\; c, c' \in C_i,\;
      \Pa_i^c \ne \Pa_i^{c'}\}\bigr|.
    \]

  \item \textbf{Emergence rate}:
    \[
      \psi_i^{(3)}
      \;=\;
      1 - \frac{\min_{c \in C_i} |\mathrm{MB}_c(X_i)|}
               {\max_{c \in C_i} |\mathrm{MB}_c(X_i)| + 1},
    \]
    where $\mathrm{MB}_c(X_i)$ is the Markov blanket of~$X_i$ in
    context~$c$.  This captures the ratio of minimal to maximal
    causal connectivity: a mechanism whose variable is causally
    isolated in some contexts but richly connected in others receives
    a high emergence score.  When $X_i \notin V_c$ for some~$c$, we
    set $|\mathrm{MB}_c(X_i)| = 0$, recovering the variable-absence
    case as a special instance.

  \item \textbf{Context sensitivity}:
    $\psi_i^{(4)} = \rho_S\!\bigl(\dC(c, c'),\,
    \dmu(\mu_i^c, \mu_i^{c'})\bigr)$, the Spearman rank correlation
    between context distance and mechanism distance over all pairs
    $(c, c') \in C_i^2$.  High values indicate smooth dependence on context;
    low values indicate mechanism change is decoupled from context similarity.
\end{enumerate}
\end{definition}

\begin{definition}[Plasticity Spectrum]\label{def:plasticity-spectrum}
Given a plasticity descriptor $\psi_i$ and thresholds
$\tau_{\mathrm{inv}}, \tau_{\mathrm{struct}}, \tau_E > 0$, the \emph{plasticity
class} $\kappa_i$ of mechanism $\mu_i$ is determined by the following
rules, applied in priority order:
\begin{enumerate}[nosep,label=(\roman*)]
  \item \textbf{Emergent:} $\kappa_i = \textsc{emergent}$ if
        $\psi_i^{(3)} > \tau_E$\, (i.e., $X_i$ has substantially
        different causal connectivity across contexts, including the
        extreme case of variable absence).  A mechanism is
        \emph{fully emergent} if $X_i \notin V_c$ for some~$c$
        and \emph{partially emergent} otherwise.
  \item \textbf{Structurally plastic:}
        $\kappa_i = \textsc{struct\text{-}plastic}$ if
        $\psi_i^{(2)} > \tau_{\mathrm{struct}}$ and $\psi_i^{(3)} = 0$.
  \item \textbf{Parametrically plastic:}
        $\kappa_i = \textsc{param\text{-}plastic}$ if $\psi_i^{(2)} = 0$
        and $\psi_i^{(1)} \ge \tau_{\mathrm{inv}}$.
  \item \textbf{Invariant:} $\kappa_i = \textsc{invariant}$ if
        $\psi_i^{(1)} < \tau_{\mathrm{inv}}$, $\psi_i^{(2)} = 0$, and
        $\psi_i^{(3)} = 0$.
\end{enumerate}
These categories partition the set of mechanisms.  The ordering
reflects severity: emergence supersedes structural plasticity, which
supersedes parametric plasticity, which supersedes invariance.
\end{definition}

\begin{definition}[Mechanism Stability Class]\label{def:stability-class}
A \emph{mechanism stability class} is a connected component of the graph
$\mathcal{G}_\eps$ whose vertices are mechanisms and whose edges connect
$\mu_i$ and $\mu_j$ whenever:
\begin{enumerate}[nosep,label=(\arabic*)]
  \item $\psi_i$ and $\psi_j$ belong to the same plasticity category
        (\Cref{def:plasticity-spectrum}), and
  \item $\|\psi_i - \psi_j\|_\infty < \eps$ for a resolution parameter
        $\eps > 0$.
\end{enumerate}
Since connected-component membership is transitive by construction,
this yields a well-defined partition
$\mathcal{S} = \{s_1, \ldots, s_m\}$ of the mechanism space.  Each
class has a \emph{prototype descriptor}
$\bar{\psi}_s = \mathrm{mean}(\{\psi_i : \mu_i \in s\})$.  Within each
class, we further sub-partition by the dominant context-sensitivity pattern
(component~$\psi^{(4)}$), identifying sub-classes of mechanisms
that respond to context change in qualitatively similar ways.
\end{definition}

\begin{definition}[Data-Derived Context Distance]\label{def:data-context-distance}
When contexts lack explicit feature descriptors, the context distance is
derived from data.  Let $\hat{G}_c, \hat{G}_{c'}$ be estimated CPDAGs and
let $P_c, P_{c'}$ be the observational distributions restricted to the
shared variable set $V_c \cap V_{c'}$.  Define
\[
  \dC^{\mathrm{data}}(c, c')
  \;=\;
  \sqrt{
    \dDAG(\hat{G}_c, \hat{G}_{c'})^2
    \;+\;
    \MMD^2\!\bigl(P_c|_{V_c \cap V_{c'}},\;
                   P_{c'}|_{V_c \cap V_{c'}}\bigr)
  }.
\]
Since both $\dDAG$ and $\MMD$ are pseudo-metrics, their $\ell^2$
combination is also a pseudo-metric (triangle inequality is preserved).
\end{definition}

\begin{remark}[Computation Order and Circularity]
The data-derived context distance $d_C^{\mathrm{data}}$ uses estimated
DAGs and distributions, which in turn are used to compute plasticity
descriptors that reference~$d_C$.  Logical circularity is avoided
because the computation is strictly ordered: (1)~per-context DAGs are
estimated independently, (2)~$d_C^{\mathrm{data}}$ is computed from
these fixed DAGs and distributions, and (3)~plasticity descriptors are
computed using the frozen~$d_C$.  When explicit context features are
available (e.g., experimental conditions, dosage levels), we recommend
using a feature-based $d_C$ to eliminate this dependence entirely.
\end{remark}

\begin{definition}[Quality-Diversity Archive]\label{def:qd-archive}
The \emph{QD archive} (atlas) is a tuple $\cA = (\cB, \phi, q)$ where:
\begin{itemize}[nosep]
  \item $\cB = \{B_1, \ldots, B_K\}$ is a \emph{centroidal Voronoi
        tessellation} (CVT) of the behaviour space
        $\Psi \subseteq \R^4$ into $K$ cells, with centroids computed by
        Lloyd's algorithm on a sample of plasticity descriptors.
  \item $\phi \colon \cE \to \Psi$ is the \emph{behaviour
        characterisation function} mapping an atlas entry
        $e \in \cE$ to its plasticity descriptor.  An atlas entry
        $e = (X_i, \alpha_{ij}, \{\mu_i^c\}_{c \in S_e})$ consists of a
        target variable, its alignment across a context subset, and the
        corresponding mechanisms.
  \item $q \colon \cE \to \R_{\ge 0}$ is the \emph{quality function}:
        \[
          q(e)
          \;=\;
          \underbrace{|S_e|}_{\text{coverage}}
          \;\cdot\;
          \underbrace{\Bigl(1 + \mathrm{Var}_{c \in S_e}
                     \bigl[\dmu(\mu_i^c, \bar{\mu}_i)\bigr]\Bigr)^{-1}
                    }_{\text{estimation confidence}}
          \;\cdot\;
          \underbrace{r(e)}_{\text{robustness}},
        \]
        where $\bar{\mu}_i$ is the mean mechanism descriptor and
        $r(e) \in [0,1]$ is the robustness score
        (\Cref{def:robustness-certificate}).
\end{itemize}
Each cell $B_k$ stores at most one entry: the entry with the highest
quality among all entries whose behaviour characterisation falls in~$B_k$.
The archive is \emph{complete} if every occupied cell has quality above a
threshold~$q_{\min}$.
\end{definition}

\begin{definition}[Curiosity Signal]\label{def:curiosity}
The \emph{curiosity signal} for a candidate atlas entry~$e$ with behaviour
$\phi(e) = \psi$ is
\[
  \mathrm{Curiosity}(e;\, \cA)
  \;=\;
  \underbrace{d_{\mathrm{cell}}(\psi, \cA)}_{\text{novelty}}
  \;+\;
  \beta \cdot
  \underbrace{\Delta q(e, \cA)}_{\text{quality improvement}},
\]
where:
\begin{itemize}[nosep]
  \item $d_{\mathrm{cell}}(\psi, \cA)
        = \min_{k : B_k \text{ occupied}} \|\psi - \bar{\psi}_k\|_2$
        if $\psi$ falls in an empty cell, or~$0$ otherwise
        (\emph{novelty}: entries filling empty cells have positive
        novelty);
  \item $\Delta q(e, \cA) = \max(0,\, q(e) - q_{\mathrm{current}}(B_k))$,
        where $B_k$ is the cell containing~$\psi$
        (\emph{quality improvement}: positive if~$e$ would replace a
        lower-quality incumbent);
  \item $\beta > 0$ is a trade-off parameter.
\end{itemize}
An entry is \emph{curious} if $\mathrm{Curiosity}(e; \cA) > 0$.
\end{definition}

\begin{definition}[Tipping Point]\label{def:tipping-point}
Let $\gamma \colon [0,1] \to \cC$ be a \emph{context path} (a continuous
or piecewise-linear curve through context space).  For mechanism~$\mu_i$,
define the \emph{plasticity functional along~$\gamma$}:
\[
  \Phi_i(\gamma, t)
  \;=\;
  \dmu\!\bigl(\mu_i^{\gamma(0)},\; \mu_i^{\gamma(t)}\bigr).
\]
A \emph{tipping point} for mechanism $\mu_i$ along path~$\gamma$ is a
parameter value $t^* \in (0,1)$ such that:
\begin{enumerate}[nosep,label=(\arabic*)]
  \item \textbf{Near-discontinuity:}
    $\displaystyle
      \limsup_{h \to 0^+}
      \frac{\Phi_i(\gamma, t^* + h) - \Phi_i(\gamma, t^*)}
           {\dC(\gamma(t^* + h),\, \gamma(t^*))}
      > \tau_{\mathrm{tip}};
    $
  \item \textbf{Significance:} a bootstrap confidence interval for the
        normalised increment does not contain zero at level $1 - \alpha$.
\end{enumerate}
A \emph{structural tipping point} additionally requires
$\Pa_i^{\gamma(t^* - \eps)} \ne \Pa_i^{\gamma(t^* + \eps)}$ for all
sufficiently small $\eps > 0$.

In the discrete setting (finite $\Cobs$), a tipping point is detected
between consecutive contexts $c_j, c_{j+1}$ on a context path if
$\Phi_i(\gamma, t_{j+1}) - \Phi_i(\gamma, t_j) >
\tau_{\mathrm{tip}} \cdot \dC(c_j, c_{j+1})$.
\end{definition}

\begin{definition}[Robustness Certificate]\label{def:robustness-certificate}
A \emph{robustness certificate} for a plasticity classification
$\kappa_i \in \{\textsc{invariant},\, \textsc{param\text{-}plastic},\,
\textsc{struct\text{-}plastic},\, \textsc{emergent}\}$
of mechanism~$\mu_i$ is a tuple
$\cR_i = (\kappa_i, \delta_i, n_{\min}, \alpha)$ where:
\begin{enumerate}[nosep,label=(\roman*)]
  \item $\kappa_i$ is the estimated plasticity class
        (\Cref{def:plasticity-spectrum});
  \item $\delta_i = \min_{b \in \partial \mathcal{K}}
        \|\hat{\psi}_i - b\|_\infty$ is the \emph{margin}---the
        $\ell^\infty$ distance from the estimated plasticity descriptor
        $\hat{\psi}_i$ to the nearest classification boundary
        $\partial\mathcal{K}$;
  \item $n_{\min}$ is the minimum per-context sample size used;
  \item $\alpha \in (0,1)$ is the failure probability.
\end{enumerate}
The certificate is \emph{valid} if
$\delta_i > \eps_{\mathrm{est}}(n_{\min}, \alpha)$, where
\[
  \eps_{\mathrm{est}}(n_{\min}, \alpha)
  \;=\;
  C' \sqrt{\frac{d_{\max} \log |\Vbar| + \log(2/\alpha)}{n_{\min}}}
\]
for a universal constant~$C'$ and maximum in-degree $d_{\max}$.
The certificate guarantees
$P(\hat{\kappa}_i = \kappa_i^*) \ge 1 - \alpha$.
\end{definition}

% ---- 3.2 Assumptions -----------------------------------------
\subsection{Assumptions}\label{sec:assumptions}

We state six assumptions, partitioned into causal assumptions (A1--A3),
distributional assumptions (A4--A5), and a regularity condition (A6).

\begin{assumption}[Per-Context Markov Condition]\label{ass:markov}
For each context $c \in \Cobs$, the observational distribution~$P_c$ is
Markov with respect to the DAG~$G_c$: every variable~$X_i$ is
conditionally independent of its non-descendants given its parents
$\Pa_i^c$.  The DAG~$G_c$ is identified only up to its Markov equivalence
class (CPDAG) from observational data; all parent-set-dependent definitions
(\Cref{def:mechanism}--\Cref{def:robustness-certificate}) are stated with
respect to the estimated CPDAG, with ``parent set'' interpreted as the set
of variables with a directed edge into~$X_i$ that is present in every
member of the equivalence class.
\end{assumption}

\begin{assumption}[Per-Context Faithfulness]\label{ass:faithfulness}
For each context $c \in \Cobs$, the distribution~$P_c$ is faithful to the
DAG~$G_c$: every conditional independence holding in~$P_c$ is entailed by
the Markov condition applied to~$G_c$.  This can be relaxed to
\emph{adjacency-faithfulness} \citep{ramsey2006adjacency}: if~$X_i$
and~$X_j$ are adjacent in~$G_c$, then no conditioning set renders them
conditionally independent.
\end{assumption}

\begin{assumption}[Independent Causal Mechanisms]\label{ass:icm}
Within each context~$c$, the mechanism $P_c(X_i \mid \Pa_i^c)$ is
``independent'' of the marginal distribution $P_c(\Pa_i^c)$ in the sense
of \citet{scholkopf2012causal}: knowing the marginal distribution of the
parents provides no information about the conditional distribution of
$X_i$ given its parents.  Formally, in an appropriate
information-geometric sense, the mechanism and the cause distribution
occupy orthogonal submanifolds of the statistical model.  Across contexts,
mechanisms may change independently: a change in $P_c(X_i \mid \Pa_i^c)$
need not be accompanied by a change in $P_c(\Pa_j^c)$ for $j \ne i$.
\end{assumption}

\begin{assumption}[Finite Contexts with Sufficient Samples]\label{ass:samples}
The observed context set $\Cobs$ is finite with $|\Cobs| \ge 3$.  Each
context~$c$ has $n_c \ge n_{\min}$ i.i.d.\ samples drawn from~$P_c$,
where
\[
  n_{\min}
  \;=\;
  \Omega\!\left(
    \frac{d_{\max}^2 \cdot \sigma^2}{\eps^2}
    \cdot \log\frac{|\Vbar| \cdot |\Cobs|}{\alpha}
  \right)
\]
for target estimation error~$\eps$, failure probability~$\alpha$,
maximum in-degree $d_{\max}$, and maximum marginal variance~$\sigma^2$.
The samples within each context are independent and identically
distributed; samples across contexts are independent (no cross-context
data leakage).
\end{assumption}

\begin{assumption}[Sub-Gaussian or Bounded Variables]\label{ass:subgaussian}
Each endogenous variable $X_i$ under each context~$c$ is either:
\begin{enumerate}[nosep,label=(\alph*)]
  \item discrete with finite support $|\mathrm{supp}(X_i)| \le M$ for a
        known bound~$M$, or
  \item continuous and sub-Gaussian with parameter~$\sigma_c$:
        $\mathbb{E}[\exp(t(X_i - \mathbb{E}[X_i]))] \le
        \exp(\sigma_c^2 t^2 / 2)$ for all $t \in \R$.
\end{enumerate}
In either case, the maximum marginal variance
$\sigma^2 = \max_{c,i} \mathrm{Var}_c(X_i)$ is finite and known (or
conservatively estimated from data).
\end{assumption}

\begin{assumption}[Context Regularity]\label{ass:regularity}
The map $c \mapsto P_c$ from the context space $(\cC, \dC)$ to the space
of observational distributions (with the total variation or MMD topology)
is \emph{piecewise Lipschitz}: there exist finitely many regimes
$R_1, \ldots, R_L$ with $\cC = R_1 \cup \cdots \cup R_L$ such that within
each regime~$R_l$:
\[
  \MMD(P_{c_1}, P_{c_2})
  \;\le\;
  L_l \cdot \dC(c_1, c_2)
  \qquad
  \text{for all } c_1, c_2 \in R_l,
\]
with Lipschitz constant $L_l > 0$.  Tipping points
(\Cref{def:tipping-point}) occur at the boundaries between regimes.  The
number of regimes~$L$ is finite but unknown a priori.
\end{assumption}

% ---- 3.3 Theorem 1 -------------------------------------------
\subsection{DAG Alignment Metric Properties}\label{sec:thm-metric}

We now state and prove the first main theorem, which establishes that the
DAG alignment distance is a well-behaved pseudo-metric, providing the
geometric foundation for the CVT tessellation of the QD archive.

\begin{theorem}[DAG Alignment Metric Properties]\label{thm:dag-metric}
Let $\mathcal{G}(\Vbar)$ denote the set of all DAGs on subsets of~$\Vbar$.
Under the alignment cost of \Cref{def:dag-alignment} with
$\lambda_{\mathrm{s}}, \lambda_{\mathrm{m}}, \lambda_{\mathrm{v}} > 0$:
\begin{enumerate}[label=(\alph*)]
  \item $\dDAG$ is a pseudo-metric on $\mathcal{G}(\Vbar)$: it satisfies
        non-negativity, symmetry, and the triangle inequality.

  \item $\dDAG(G_c, G_{c'}) = 0$ if and only if there exists an alignment
        $\alpha$ that is a bijection $V_c \to V_{c'}$ with
        $\Delta E(\alpha) = \emptyset$ and
        $\dmu(\mu_i^c, \mu_{\alpha(i)}^{c'}) = 0$ for all~$i$.

  \item If the mechanism distance $\dmu$ is a proper metric (e.g., MMD or
        $\JSD$), then $\dDAG$ is a metric on
        $\mathcal{G}(\Vbar) / {\cong}$\, (DAGs modulo variable
        relabelling).
\end{enumerate}
\end{theorem}

The proof relies on two key lemmas.  The first establishes that variable
correspondences (partial injections) compose correctly---an essential
building block for the triangle inequality, since the intermediate DAG's
variable set may differ from both endpoints.

\begin{lemma}[Alignment Composition (Variable Correspondence Composition)]\label{lem:composition}
For partial injections $\alpha \colon V_a \rightharpoonup V_b$ and
$\beta \colon V_b \rightharpoonup V_c$, the composite
$\gamma = \beta \circ \alpha \colon V_a \rightharpoonup V_c$ is a valid
partial injection with
$|\Dom(\gamma)| \ge |\Dom(\alpha) \cap \beta^{-1}(V_c)|$.
\end{lemma}

\begin{proof}
Since $\alpha$ is injective on $\Dom(\alpha)$ and $\beta$ is injective on
$\Dom(\beta)$, the composition $\gamma(x) = \beta(\alpha(x))$ is defined
for all $x \in \Dom(\alpha)$ with $\alpha(x) \in \Dom(\beta)$, and
injectivity of~$\gamma$ follows from injectivity of both components.
The domain bound is immediate: $\Dom(\gamma) = \{x \in \Dom(\alpha) :
\alpha(x) \in \Dom(\beta)\} = \Dom(\alpha) \cap \alpha^{-1}(\Dom(\beta))$,
which has cardinality at least $|\Dom(\alpha) \cap \beta^{-1}(V_c)|$ since
$\Dom(\beta) \subseteq V_b$.
\end{proof}

\begin{lemma}[Cost Sub-Additivity]\label{lem:subadditivity}
Under the composite alignment $\gamma = \beta \circ \alpha$, each cost
component satisfies
$C_k(\gamma) \le C_k(\alpha) + C_k(\beta)$ for $k \in \{\mathrm{s},
\mathrm{m}, \mathrm{v}\}$.
\end{lemma}

\begin{proof}
We verify each component separately.

\emph{Structural cost.}
Let $(X_i, X_j)$ be an edge in~$G_a$ with
$X_i, X_j \in \Dom(\gamma)$.  If this edge disagrees with~$G_c$
under~$\gamma$, then either:
(i)~the edge $(X_i, X_j) \mapsto (\alpha(X_i), \alpha(X_j))$ already
disagrees with~$G_b$ (contributing to $|\Delta E(\alpha)|$), or
(ii)~the edge $(\alpha(X_i), \alpha(X_j))$ agrees with~$G_b$ but
$(\beta(\alpha(X_i)), \beta(\alpha(X_j)))$ disagrees with~$G_c$
(contributing to $|\Delta E(\beta)|$).
Hence $|\Delta E(\gamma)| \le |\Delta E(\alpha)| + |\Delta E(\beta)|$.

\emph{Mechanism cost.}
For each $X_i \in \Dom(\gamma)$, the triangle inequality for $\dmu$
(which holds because $\JSD$ is a metric \citep{endres2003new}) gives
\[
  \dmu(\mu_i^a, \mu_{\gamma(i)}^c)
  \;\le\;
  \dmu(\mu_i^a, \mu_{\alpha(i)}^b)
  \;+\;
  \dmu(\mu_{\alpha(i)}^b, \mu_{\beta(\alpha(i))}^c).
\]
Summing over $X_i \in \Dom(\gamma)$ and noting that
$\Dom(\gamma) \subseteq \Dom(\alpha)$ and
$\gamma(\Dom(\gamma)) \subseteq \Img(\beta)$ yields
\[
  \sum_{X_i \in \Dom(\gamma)} \dmu(\mu_i^a, \mu_{\gamma(i)}^c)
  \;\le\;
  \sum_{X_i \in \Dom(\alpha)} \dmu(\mu_i^a, \mu_{\alpha(i)}^b)
  \;+\;
  \sum_{X_j \in \Dom(\beta)} \dmu(\mu_j^b, \mu_{\beta(j)}^c),
\]
where the inequality uses the fact that the sums on the right-hand side
contain all terms appearing on the left.

\emph{Variable-mismatch cost.}
Write $U_\gamma = |V_a \setminus \Dom(\gamma)| +
|V_c \setminus \Img(\gamma)|$ for the number of unmatched variables.
By inclusion--exclusion:
\begin{align*}
  |V_a \setminus \Dom(\gamma)|
  &\;\le\;
  |V_a \setminus \Dom(\alpha)|
  + |\Dom(\alpha) \setminus \Dom(\gamma)| \\
  &\;\le\;
  |V_a \setminus \Dom(\alpha)|
  + |V_b \setminus \Dom(\beta)|,
\end{align*}
where the second inequality holds because
$\Dom(\alpha) \setminus \Dom(\gamma) =
\{x \in \Dom(\alpha) : \alpha(x) \notin \Dom(\beta)\}$,
so $|\Dom(\alpha) \setminus \Dom(\gamma)| \le |V_b \setminus \Dom(\beta)|$.
Symmetrically,
$|V_c \setminus \Img(\gamma)| \le
|V_b \setminus \Img(\alpha)| + |V_c \setminus \Img(\beta)|$.
Hence
\[
  U_\gamma
  \;\le\;
  \bigl(|V_a \setminus \Dom(\alpha)| + |V_b \setminus \Img(\alpha)|\bigr)
  +
  \bigl(|V_b \setminus \Dom(\beta)| + |V_c \setminus \Img(\beta)|\bigr)
  = U_\alpha + U_\beta. \qedhere
\]
\end{proof}

\begin{proof}[Proof of \Cref{thm:dag-metric}]
\textbf{Part (a).}

\emph{Non-negativity.}
Each term in \eqref{eq:alignment-cost} is non-negative (cardinalities and
distances are non-negative, weights are positive), so $\dDAG \ge 0$.

\emph{$\dDAG(G, G) = 0$.}
The identity alignment $\alpha = \mathrm{id}_{V}$ satisfies
$\Delta E(\alpha) = \emptyset$,
$\dmu(\mu_i, \mu_i) = 0$ for all~$i$ (since $\JSD(P, P) = 0$), and
$V \setminus \Dom(\alpha) = V \setminus \Img(\alpha) = \emptyset$.
Hence $\Cost(\mathrm{id}; G, G) = 0$, and since the minimum over
alignments is at most this value, $\dDAG(G, G) = 0$.

\emph{Symmetry.}
If $\alpha^* \colon V_a \rightharpoonup V_b$ is an optimal alignment from
$G_a$ to~$G_b$, then the reverse partial injection
$(\alpha^*)^{-1} \colon V_b \rightharpoonup V_a$ satisfies
$\Cost((\alpha^*)^{-1}; G_b, G_a) = \Cost(\alpha^*; G_a, G_b)$, since:
$|\Delta E|$ is symmetric (an edge disagreement under~$\alpha$ is an edge
disagreement under~$\alpha^{-1}$); $\dmu$ is symmetric (because $\JSD$ is
symmetric); and the variable-mismatch penalty
$|V_a \setminus \Dom(\alpha)| + |V_b \setminus \Img(\alpha)|
= |V_b \setminus \Dom(\alpha^{-1})| + |V_a \setminus \Img(\alpha^{-1})|$.
Hence $\dDAG(G_b, G_a) \le \dDAG(G_a, G_b)$; by the symmetric argument,
equality holds.

\emph{Triangle inequality.}
Let $\alpha^*$ and $\beta^*$ be optimal alignments for
$(G_a, G_b)$ and $(G_b, G_c)$, respectively.  By
\Cref{lem:composition}, $\gamma = \beta^* \circ \alpha^*$ is a valid
partial injection $V_a \rightharpoonup V_c$.  By
\Cref{lem:subadditivity}, each weighted cost component of~$\gamma$
is at most the sum of the corresponding components of~$\alpha^*$
and~$\beta^*$.  Therefore
\begin{align*}
  \dDAG(G_a, G_c)
  &\;\le\;
  \Cost(\gamma;\, G_a, G_c) \\
  &\;\le\;
  \Cost(\alpha^*;\, G_a, G_b) + \Cost(\beta^*;\, G_b, G_c) \\
  &\;=\;
  \dDAG(G_a, G_b) + \dDAG(G_b, G_c).
\end{align*}

\medskip
\textbf{Part (b).}
$(\Rightarrow)$:\ If $\dDAG(G_c, G_{c'}) = 0$, then the minimum cost is
zero.  Since each term in~\eqref{eq:alignment-cost} is non-negative and
the weights are strictly positive, every term must individually be zero:
$|\Delta E(\alpha^*)| = 0$ (no edge disagreements),
$\dmu(\mu_i^c, \mu_{\alpha^*(i)}^{c'}) = 0$ for all $i \in \Dom(\alpha^*)$
(identical mechanisms), and
$|V_c \setminus \Dom(\alpha^*)| = |V_{c'} \setminus \Img(\alpha^*)| = 0$
(bijection).

$(\Leftarrow)$:\ If such a bijection~$\alpha$ exists, then
$\Cost(\alpha; G_c, G_{c'}) = 0$, so $\dDAG(G_c, G_{c'}) = 0$.

\medskip
\textbf{Part (c).}
When $\dmu$ is a proper metric, $\dmu(\mu_i^c, \mu_j^{c'}) = 0$ implies
the mechanisms are identical.  By part~(b), $\dDAG = 0$ implies the DAGs
are identical up to the relabelling~$\alpha$ (a bijection with no edge
disagreements and identical mechanisms).  Hence $\dDAG$ separates
equivalence classes under variable relabelling, satisfying the identity of
indiscernibles on $\mathcal{G}(\Vbar)/{\cong}$, and together with
non-negativity, symmetry, and the triangle inequality from part~(a), this
yields a metric.
\end{proof}

\begin{remark}
The triangle inequality proof relies critically on the JSD being a proper
metric (specifically, satisfying the triangle inequality).  This is
established by \citet{endres2003new}, who show that $\sqrt{\JS(P \| Q)}$
is a metric on the space of probability distributions.  Had we used the KL
divergence directly---as in the original proposal---the triangle inequality
would fail, and the CVT tessellation of the QD archive would have no
geometric justification.
\end{remark}

\begin{remark}
The NP-hardness of computing the optimal alignment
$\min_\alpha \Cost(\alpha; G_c, G_{c'})$ (by reduction from subgraph
isomorphism) does not affect the validity of the metric properties, which
are independent of computational tractability.  In practice, we exploit
bounded in-degree to obtain polynomial-time algorithms (see Section~4).
\end{remark}


% ==============================================================
% Bibliography placeholder
% ==============================================================
\bibliographystyle{plainnat}
% \bibliography{references}

% --- Placeholder references for compilation ---

%%  Section 4: Algorithms
%% ============================================================

\section{Algorithms}\label{sec:algorithms}

We present five algorithms composing the Causal-Plasticity Atlas pipeline.
Phase~1 (Foundation) consists of \textsc{CADA} (Sec.~\ref{sec:alg-cada}) and
\textsc{PDC} (Sec.~\ref{sec:alg-pdc}); Phase~2 (Exploration) applies
\textsc{CD-QD} (Sec.~\ref{sec:alg-cdqd}); Phase~3 (Validation) runs
\textsc{TPD} (Sec.~\ref{sec:alg-tpd}) and \textsc{RCG}
(Sec.~\ref{sec:alg-rcg}).  All algorithms use the notation of
Section~\ref{sec:framework}: $K$ contexts, $n$ variables in the union pool
$\bar{V}$, maximum in-degree $d_{\max}$, and $\sqrt{\mathrm{JSD}}$ as the
mechanism distance $d_\mu$ (Definition~\ref{def:mech-distance}).

%% ------------------------------------------------------------------
\subsection{Algorithm 1: Context-Aware DAG Alignment (CADA)}
\label{sec:alg-cada}

\paragraph{Problem.}
Given two CPDAGs $\hat{G}_a=(V_a,E_a)$ and $\hat{G}_b=(V_b,E_b)$ from
distinct contexts, together with sufficient statistics
$(\Sigma_a,\Sigma_b)$ and a set of semantic \emph{anchors}
$A_0=\{(v_a^{(i)},v_b^{(i)})\}$, compute a maximal partial injection
$\pi:V_a\rightharpoonup V_b\cup\{\bot\}$ and an edge partition
$(S_{\mathrm{shared}},S_{\mathrm{mod}},S_{a},S_{b})$.

General subgraph isomorphism is NP-complete, but three structural
properties of causal DAGs restore tractability: (i)~sparsity
($|E|=O(n)$ under bounded in-degree), (ii)~semantic anchoring (most
variables have known identity across contexts, leaving a small
unanchored set $U$), and (iii)~Markov-blanket locality (plausible
alignment candidates must share anchored neighbourhood structure).
Under assumptions $d_{\max}\le k$ and $|U|\le u$, these reduce the
problem to $O(n^2 d_{\max}^2+u^3)$ time (Theorem~\ref{thm:hardness}).

\begin{algorithm}[t]
\caption{\textsc{CADA}: Context-Aware DAG Alignment}
\label{alg:cada}
\begin{algorithmic}[1]
\Require CPDAGs $\hat{G}_a=(V_a,E_a)$, $\hat{G}_b=(V_b,E_b)$;
  sufficient statistics $\Sigma_a,\Sigma_b$;
  anchors $A_0\subseteq V_a\times V_b$;
  significance level $\alpha$; cap $u_{\max}$
\Ensure Alignment $\pi$, edge partition
  $(S_{\mathrm{shared}},S_{\mathrm{mod}},S_a,S_b)$,
  score $s\in[0,1]$, divergence $\delta$
\Statex
\Statex \textbf{Phase 1: Anchor Propagation}
\State $\pi\gets\{v_a\mapsto v_b:(v_a,v_b)\in A_0\}$
\State $U_a\gets V_a\setminus\mathrm{dom}(\pi)$;\quad
       $U_b\gets V_b\setminus\mathrm{range}(\pi)$
\If{$|U_a|>u_{\max}$ \textbf{or} $|U_b|>u_{\max}$}
  \State \textbf{raise} \textsc{TooManyUnanchored}
\EndIf
\Statex
\Statex \textbf{Phase 2: Markov Blanket Candidate Generation}
\For{each $u_a\in U_a$}
  \State $\mathrm{MB}_a\gets\mathrm{Pa}(u_a)\cup\mathrm{Ch}(u_a)
         \cup\mathrm{CoPa}(u_a)$ in $\hat{G}_a$
  \State $\mathrm{MB}_a^{\pi}\gets\{\pi(v):v\in\mathrm{MB}_a
         \cap\mathrm{dom}(\pi)\}$
  \For{each $u_b\in U_b$}
    \State $\mathrm{MB}_b\gets$ Markov blanket of $u_b$ in $\hat{G}_b$
    \State $\mathrm{MB}_b^{\pi}\gets\mathrm{MB}_b\cap\mathrm{range}(\pi)$
    \State $J\gets|\mathrm{MB}_a^{\pi}\cap\mathrm{MB}_b^{\pi}|
            \;/\;|\mathrm{MB}_a^{\pi}\cup\mathrm{MB}_b^{\pi}|$
    \If{$J\ge 0.3$}
      \State $\mathrm{Cand}(u_a)\gets\mathrm{Cand}(u_a)\cup\{u_b\}$
    \EndIf
  \EndFor
\EndFor
\Statex
\Statex \textbf{Phase 3: Statistical Compatibility Scoring}
\State Initialise $\mathbf{S}\in\mathbb{R}^{|U_a|\times|U_b|}$ to $0$
\For{each $u_a\in U_a$, $u_b\in\mathrm{Cand}(u_a)$}
  \State $s_{\mathrm{CI}}\gets\textsc{CIFingerprint}(\hat{G}_a,u_a,
         \hat{G}_b,u_b,\pi,\Sigma_a,\Sigma_b,\alpha)$
         \Comment{Alg.~\ref{alg:ci-finger}}
  \State $s_{\mathrm{dist}}\gets\textsc{LocalDistSim}(\Sigma_a,u_a,
         \Sigma_b,u_b,\pi)$
  \State $\mathbf{S}[u_a,u_b]\gets 0.6\,s_{\mathrm{CI}}
         +0.4\,s_{\mathrm{dist}}$
\EndFor
\Statex
\Statex \textbf{Phase 4: Hungarian Matching}
\State Pad $\mathbf{S}$ to a square matrix with $-\infty$ fill
\State $\sigma^*\gets\textsc{Hungarian}(\mathbf{S},\text{maximise}=
       \mathrm{true})$
\For{each $(u_a,u_b)\in\sigma^*$}
  \If{$\mathbf{S}[u_a,u_b]\ge 0.5$}
    $\pi(u_a)\gets u_b$
  \Else\; $\pi(u_a)\gets\bot$
  \EndIf
\EndFor
\Statex
\Statex \textbf{Phase 5: Edge Classification}
\State $S_{\mathrm{shared}},S_{\mathrm{mod}},S_a,S_b\gets\emptyset$
\For{each $(u,v)\in E_a$}
  \State $u'\gets\pi(u)$;\; $v'\gets\pi(v)$
  \If{$u'=\bot$ \textbf{or} $v'=\bot$}
    $S_a\gets S_a\cup\{(u,v)\}$
  \ElsIf{$(u',v')\in E_b$}
    $S_{\mathrm{shared}}\gets S_{\mathrm{shared}}\cup\{(u,v)\}$
  \ElsIf{$(v',u')\in E_b$}
    $S_{\mathrm{mod}}\gets S_{\mathrm{mod}}\cup\{(u,v,\texttt{rev})\}$
  \Else\;
    $S_{\mathrm{mod}}\gets S_{\mathrm{mod}}\cup\{(u,v,\texttt{del})\}$
  \EndIf
\EndFor
\For{each $(u',v')\in E_b$ with $\pi^{-1}(u')\ne\bot$,
     $\pi^{-1}(v')\ne\bot$}
  \If{$(\pi^{-1}(u'),\pi^{-1}(v'))\notin E_a$ \textbf{and}
      $(\pi^{-1}(v'),\pi^{-1}(u'))\notin E_a$}
    \State $S_{\mathrm{mod}}\gets S_{\mathrm{mod}}
           \cup\{(u',v',\texttt{add})\}$
  \EndIf
\EndFor
\For{each $(u',v')\in E_b$ with $\pi^{-1}(u')=\bot$ \textbf{or}
     $\pi^{-1}(v')=\bot$}
  \State $S_b\gets S_b\cup\{(u',v')\}$
\EndFor
\Statex
\State $s\gets|S_{\mathrm{shared}}|\;/\;
       |E_a\cup\pi(E_b)|$
       \Comment{Jaccard over edges}
\State $\delta\gets\textsc{WeightedSHD}(S_{\mathrm{shared}},
       S_{\mathrm{mod}},S_a,S_b,n)$
       \Comment{$w_{\mathrm{rev}}=0.5,\;w_{\mathrm{add/del}}=1,\;
       w_{\mathrm{ctx}}=0.8$}
\State \Return $(\pi,S_{\mathrm{shared}},S_{\mathrm{mod}},S_a,S_b,s,\delta)$
\end{algorithmic}
\end{algorithm}

\begin{algorithm}[t]
\caption{\textsc{CIFingerprint}: Conditional-Independence Pattern Similarity}
\label{alg:ci-finger}
\begin{algorithmic}[1]
\Require DAGs $\hat{G}_a,\hat{G}_b$; candidates $u_a,u_b$;
  alignment $\pi$; statistics $\Sigma_a,\Sigma_b$; level $\alpha$
\Ensure Similarity score $s_{\mathrm{CI}}\in[0,1]$
\State $W\gets\mathrm{dom}(\pi)\cap N_2(\hat{G}_a,u_a)$
  \Comment{anchored vars within distance 2}
\State $\mathrm{agree}\gets 0$;\; $T\gets 0$
\For{each $w\in W$}
  \State $p_a\gets\textsc{PartialCorr}(\Sigma_a,u_a,w,
         \mathrm{Pa}_a(u_a))$
  \State $p_b\gets\textsc{PartialCorr}(\Sigma_b,u_b,\pi(w),
         \pi(\mathrm{Pa}_a(u_a)\cap\mathrm{dom}(\pi)))$
  \State $T\gets T+1$
  \If{$(p_a>\alpha)=(p_b>\alpha)$}
    $\mathrm{agree}\gets\mathrm{agree}+1$
  \EndIf
\EndFor
\State \Return $\mathrm{agree}\;/\;\max(T,1)$
\end{algorithmic}
\end{algorithm}

\paragraph{Complexity.}
Phase~1 runs in $O(n)$.
Phase~2 computes Jaccard overlaps of Markov blankets of size
$O(d_{\max}^2)$ for $|U_a|\cdot|U_b|$ pairs: $O(u^2 d_{\max}^2)$.
Phase~3 performs $O(u^2)$ partial-correlation tests, each $O(d_{\max}^2)$:
total $O(u^2 d_{\max}^2)$.
Phase~4 is the Hungarian algorithm on a $u\times u$ matrix: $O(u^3)$.
Phase~5 is linear in $|E_a|+|E_b|=O(n d_{\max})$.
Overall:
\[
  T_{\textsc{CADA}}=O\!\bigl(n^2 d_{\max}^2+u^3\bigr),
\]
since $u\le u_{\max}\le 10$ in the target regime, this is effectively
$O(n^2)$.  Space is $O(n^2)$ for the score matrix.

\paragraph{Correctness.}
Soundness follows from the Hungarian algorithm producing a valid
maximum-weight matching on the bipartite candidate graph.
Completeness relies on the Markov-blanket pruning threshold ($J\ge
0.3$): under faithfulness (Assumption~A2), any alignment pair with high
CI-fingerprint agreement must share substantial Markov-blanket overlap,
so no near-optimal match is discarded.
Theorem~\ref{thm:hardness} formalises the approximation guarantee.

%% ------------------------------------------------------------------
\subsection{Algorithm 2: Plasticity Descriptor Computation (PDC)}
\label{sec:alg-pdc}

Given $K$ aligned DAGs (common variable indexing post-\textsc{CADA}),
parameterisations $\{\Theta_c\}_{c=1}^{K}$, and sufficient statistics
$\{\Sigma_c\}_{c=1}^{K}$, \textsc{PDC} computes the four-component
plasticity descriptor $\psi_i\in\mathbb{R}^4$ for every mechanism $\mu_i$
(Definition~\ref{def:plasticity-descriptor}).

\begin{algorithm}[t]
\caption{\textsc{PDC}: Plasticity Descriptor Computation}
\label{alg:pdc}
\begin{algorithmic}[1]
\Require Aligned DAGs $\{\hat{G}_c\}_{c=1}^K$;
  parameterisations $\{\Theta_c\}$;
  sufficient statistics $\{\Sigma_c\}$; level $\alpha$
\Ensure Descriptors $\{\psi_i\}_{i=1}^n$; classifications
  $\{\kappa_i\}_{i=1}^n$
\Statex
\For{$i=1,\ldots,n$}
\Statex \hspace{1em}\textit{// Component 1: Parametric plasticity (mean mechanism distance)}
  \State $\mathcal{G}\gets\textsc{GroupByParentSet}
         (\hat{G}_1,\ldots,\hat{G}_K,i)$
  \State $\psi_i^{(1)}\gets 0$;\; $T\gets 0$
  \For{each $(P,\mathcal{C}_P)\in\mathcal{G}$
       with $|\mathcal{C}_P|\ge 2$}
    \For{each $\{a,b\}\subseteq\mathcal{C}_P$}
      \State $\psi_i^{(1)}\mathrel{+}=
             d_\mu(\mu_i^a,\mu_i^b)$
             \Comment{$\sqrt{\mathrm{JSD}}$ of conditionals}
      \State $T\gets T+1$
    \EndFor
  \EndFor
  \If{$T>0$} $\psi_i^{(1)}\gets\psi_i^{(1)}/T$ \EndIf
\Statex \hspace{1em}\textit{// Component 2: Structural plasticity rate}
  \For{$c=1,\ldots,K$}
    \State $\mathbf{b}_c^{(i)}[j]\gets\mathbb{1}[j\in\mathrm{Pa}_c(X_i)]$
      for $j=1,\ldots,n$
  \EndFor
  \State $\psi_i^{(2)}\gets
         \frac{2}{K(K{-}1)}\bigl|\{(c,c'){:}\,c{<}c',\,
         \mathrm{Pa}_c(X_i){\ne}\mathrm{Pa}_{c'}(X_i)\}\bigr|$
\Statex \hspace{1em}\textit{// Component 3: Emergence score}
  \State $\psi_i^{(3)}\gets 1-\min_c|\mathrm{MB}_c(X_i)|
         \;/\;(\max_c|\mathrm{MB}_c(X_i)|+1)$
\Statex \hspace{1em}\textit{// Component 4: Context sensitivity (Spearman $\rho$)}
  \State $\psi_i^{(4)}\gets\rho_S\!\bigl(
         \{d_C(c,c')\}_{c<c'},\;
         \{d_\mu(\mu_i^c,\mu_i^{c'})\}_{c<c'}\bigr)$
         \Comment{$O(K^2)$ pairwise}
\Statex \hspace{1em}\textit{// Component 5: Bootstrap confidence intervals}
  \State $\mathrm{CI}_S,\mathrm{CI}_P\gets
         \textsc{StabilityBootstrap}
         (\hat{G}_{1:K},\Theta_{1:K},\Sigma_{1:K},i,\alpha)$
         \Comment{stability selection for structure; parametric
         bootstrap ($B=200$) for parameters}
\Statex
\Statex \hspace{1em}\textit{// Classification (priority: emergent $>$
         mixed $>$ struct $>$ param $>$ invariant)}
  \State $\kappa_i\gets\textsc{Classify}(\psi_i,\mathrm{CI}_S,
         \mathrm{CI}_P;\;\tau_S{=}0.1,\;\tau_P{=}0.5,\;\tau_E{=}0.5)$
\EndFor
\State \Return $\{(\psi_i,\kappa_i)\}_{i=1}^n$
\end{algorithmic}
\end{algorithm}

\paragraph{Component details.}
\emph{Parametric plasticity} $\psi_i^{(1)}$ averages pairwise
$\sqrt{\mathrm{JSD}}$ distances between conditional distributions
$P_c(X_i\mid\mathrm{Pa}_c(X_i))$ within groups sharing the same parent
set.  For linear-Gaussian SEMs the $\mathrm{JSD}$ has a closed form in
regression coefficients $\beta_c^{(i)}$ and residual variances
$\sigma_c^{2(i)}$:

\emph{Structural plasticity} $\psi_i^{(2)}$ is the fraction of
unordered context pairs in which the parent set changes, i.e.,
$\psi_i^{(2)}=\frac{2}{K(K{-}1)}|\{(c,c'):c<c',\,
\mathrm{Pa}_c(X_i)\ne\mathrm{Pa}_{c'}(X_i)\}|$.
\begin{equation}\label{eq:jsd-gaussian}
  \mathrm{JSD}(P_a,P_b)
  =\tfrac{1}{2}D_{\mathrm{KL}}(P_a\|M)
  +\tfrac{1}{2}D_{\mathrm{KL}}(P_b\|M),
  \quad M=\tfrac{1}{2}(P_a+P_b),
\end{equation}
computed via numerical integration over the mixture $M$ when closed
forms are unavailable.

\emph{Emergence} $\psi_i^{(3)}$ captures the ratio of minimal to
maximal Markov-blanket size.  A mechanism whose variable is causally
isolated in some contexts but richly connected in others receives a high
emergence score.

\emph{Context sensitivity} $\psi_i^{(4)}$ is the Spearman rank
correlation between context distances $d_C(c,c')$ and mechanism
distances $d_\mu(\mu_i^c,\mu_i^{c'})$ over all $\binom{K}{2}$ context
pairs, computed in $O(K^2)$ time.  High values indicate smooth dependence
on context; low values indicate mechanism change is decoupled from
context similarity.

\emph{Confidence intervals} combine stability selection
(50\,\% subsample $\times$ 100 rounds) for the structural component with
parametric bootstrap ($B=200$ replicates, resampling residuals within the
estimated DAG structure) for $\psi_i^{(2)}$.  This avoids the
prohibitive cost of full DAG re-estimation at each bootstrap replicate.

\paragraph{Classification rule.}
With thresholds $\tau_S,\tau_P,\tau_E$ calibrated by permutation testing
on null data, and writing $\mathrm{CI}_S^{-}$, $\mathrm{CI}_P^{-}$ for
CI lower bounds:
\[
  \kappa_i=
  \begin{cases}
    \textsc{emergent}           & \text{if }\psi_i^{(3)}>\tau_E,\\
    \textsc{mixed}              & \text{if }\mathrm{CI}_S^{-}>\tau_S
                                  \text{ and }\mathrm{CI}_P^{-}>\tau_P,\\
    \textsc{structural-plastic} & \text{if }\mathrm{CI}_S^{-}>\tau_S,\\
    \textsc{parametric-plastic} & \text{if }\mathrm{CI}_P^{-}>\tau_P,\\
    \textsc{invariant}          & \text{otherwise.}
  \end{cases}
\]

\paragraph{Complexity.}
Per variable: structural plasticity is $O(nK)$; parametric plasticity
is $O(K^2 d_{\max})$; emergence is $O(K d_{\max})$; context sensitivity
is $O(K^2 d_{\max})$; bootstrap CIs cost $O(BK^2 d_{\max})$.
Summing over $n$ variables:
\[
  T_{\textsc{PDC}}=O(BnK^2 d_{\max}),
\]
where $B$ is the number of bootstrap replicates.  For
$B{=}200,\;n{=}100,\;K{=}50,\;d_{\max}{=}5$, this is approximately
$5\times10^7$ operations.

%% ------------------------------------------------------------------
\subsection{Algorithm 3: Curiosity-Driven QD Search (CD-QD)}
\label{sec:alg-cdqd}

\textsc{PDC} computes mechanism-level descriptors.  To discover
higher-order \emph{patterns}---for instance, ``mechanism $\mu_i$ is
invariant across contexts $1$--$5$ but shifts at context $6$ and
re-stabilises at $7$--$10$''---we search the combinatorial space of
context subsets $\times$ mechanism subsets using a quality-diversity (QD)
framework.

\paragraph{Search space.}
A \emph{genome} is a triple $g=(C,M,\phi)$ where
$C\subseteq[K]$ ($|C|\ge 2$) is a context subset,
$M\subseteq[n]$ is a mechanism subset, and
$\phi=(\tau_{\mathrm{thr}},\rho)$ are analysis hyper-parameters (plasticity
threshold, alignment stringency).

\paragraph{Behaviour descriptor.}
Each genome maps to $\mathbf{b}(g)\in[0,1]^4$:
\begin{align*}
  b_1 &= |M|^{-1}\bigl|\{i\in M:\kappa_i^{C}=\textsc{inv}\}\bigr|,\\
  b_2 &= |M|^{-1}\bigl|\{i\in M:\kappa_i^{C}=\textsc{param}\}\bigr|,\\
  b_3 &= |M|^{-1}\bigl|\{i\in M:\kappa_i^{C}\in
         \{\textsc{struct},\textsc{emerg},\textsc{mixed}\}\}\bigr|,\\
  b_4 &= H_{\mathrm{norm}}\!\bigl(\{\psi_i^{(1)}+\psi_i^{(2)}\}_{i\in M}\bigr),
\end{align*}
where $\kappa_i^{C}$ is the classification restricted to context set $C$ and
$H_{\mathrm{norm}}$ is normalised Shannon entropy.

\paragraph{Quality function.}
\[
  q(g)=\underbrace{H\!\bigl(\mathrm{class\_counts}(g)\bigr)}_{\text{diversity
  of mechanism types}}
  +\lambda\!\cdot\!\underbrace{\bigl(1+\min_{i\in M}
  \mathrm{width}(\mathrm{CI}_S^{(i)})\bigr)^{-1}}_{\text{statistical
  significance}}.
\]

\paragraph{Curiosity signal.}
We combine count-based novelty with prediction-error surprise:
\begin{equation}\label{eq:curiosity}
  \mathrm{Curiosity}(g;\mathcal{A})
  =\frac{1}{1+N_{\mathrm{visit}}(\mathrm{cell}(\mathbf{b}(g)))}
  +\mu\,\bigl|q(g)-\hat{q}\bigl(\mathrm{cell}(\mathbf{b}(g))\bigr)\bigr|,
\end{equation}
where $N_{\mathrm{visit}}$ is the visit count for the corresponding CVT
cell and $\hat{q}$ is an exponential moving average of observed qualities
($\alpha_{\mathrm{ema}}=0.1$).  The balance parameter $\mu=0.5$ by
default.

\paragraph{CVT tessellation.}
The behaviour space $\Psi\subseteq\mathbb{R}^4$ is partitioned into $L$
cells via a centroidal Voronoi tessellation (CVT): generate $S=10{,}000$
random genomes, compute their behaviour descriptors, and run Lloyd's
algorithm with Euclidean distance in $\sqrt{\mathrm{JSD}}$-compatible
descriptor space.  This yields approximately equal-volume cells where
descriptors concentrate, avoiding the curse of dimensionality inherent in
fixed grids.

\begin{algorithm}[t]
\caption{\textsc{CD-QD}: Curiosity-Driven Quality-Diversity Search}
\label{alg:cdqd}
\begin{algorithmic}[1]
\Require Aligned DAGs $\{\hat{G}_c\}$, descriptors $\{\psi_i\}$,
  alignment cache; iterations $T$; population size $P$;
  CVT cell count $L$; curiosity weight $\mu$
\Ensure QD archive $\mathcal{A}$
\Statex
\State $\mathcal{A}\gets\emptyset$;\;
       $N_{\mathrm{visit}}[\cdot]\gets 0$;\;
       $\hat{q}[\cdot]\gets 0$
\State Compute CVT with $L$ centroids via Lloyd's algorithm
       on $S{=}10{,}000$ random descriptors
\State \texttt{pop}$\gets\{\textsc{RandomGenome}(K,n)\}_{j=1}^{P}$
\For{$t=1,\ldots,T$}
  \For{each $g\in$\texttt{pop}}
    \State $\mathbf{b}\gets\textsc{Behaviour}(g,\{\psi_i\})$;\;
           $q\gets\textsc{Quality}(g,\{\psi_i\})$
    \State $\xi\gets\textsc{Curiosity}(\mathbf{b},q,
           N_{\mathrm{visit}},\hat{q})$
           \Comment{Eq.~\eqref{eq:curiosity}}
    \State $k\gets\textsc{CVTCell}(\mathbf{b})$
    \If{$\mathcal{A}[k]=\emptyset$ \textbf{or}
        $q>\mathcal{A}[k].q$}
      \State $\mathcal{A}[k]\gets(g,\mathbf{b},q)$
             \Comment{elitist replacement}
    \EndIf
    \State $N_{\mathrm{visit}}[k]\mathrel{+}=1$;\;
           $\hat{q}[k]\gets\alpha_{\mathrm{ema}}\,q
           +(1{-}\alpha_{\mathrm{ema}})\hat{q}[k]$
  \EndFor
  \State \texttt{parents}$\gets\textsc{WeightedSample}
         (\texttt{pop},\text{weights}=\{\xi\},|P|)$
  \State \texttt{pop}$\gets\{\textsc{Mutate}(g,K,n):g\in
         \texttt{parents}\}$
\EndFor
\State \Return $\mathcal{A}$
\end{algorithmic}
\end{algorithm}

\paragraph{Mutation operators.}
With probability $0.3$, add or remove a context from $C$; with
probability $0.3$, add or remove a mechanism from $M$; with probability
$0.2$, swap a context (remove one, add another); with probability $0.2$,
perturb analysis parameters via Gaussian noise
($\phi\gets\mathrm{clip}(\phi+\mathcal{N}(0,0.05^2),[\phi_{\min},
\phi_{\max}])$).

Because each atomic mutation (single context or mechanism addition/removal)
has positive probability and the composition of any two genomes can be
reached via a finite chain of such operations, the mutation operator is
\emph{ergodic} over the genome space
$\mathcal{G}=2^{[K]}\times 2^{[n]}\times\Phi$.

\paragraph{Convergence argument.}
Ergodicity ensures every reachable genome is visited infinitely often
with probability~1.  The genome-to-cell mapping is many-to-one, so every
reachable cell is visited infinitely often.  By the elitist replacement
rule, quality within each cell is non-decreasing.  A coupon-collector
argument with the curiosity bonus (acting as a UCB) bounds the expected
iterations to $(1{-}\varepsilon)$-coverage of the $L_r$ reachable cells
at $O(L_r\log(L_r/\varepsilon)/p_{\min})$, where $p_{\min}$ is the
minimum single-step transition probability
(Theorem~\ref{thm:atlas-completeness}).

\paragraph{Complexity.}
Each of $T$ iterations evaluates $P$ genomes; each evaluation restricts
precomputed descriptors in $O(|M|\cdot|C|)$.  Total:
$O(T\,P\,n\,K)$.  For $T{=}5000$, $P{=}100$, $n{=}100$, $K{=}50$:
$\approx 2.5\times10^9$ arithmetic comparisons, finishing in under one
minute on modern hardware.  The computational bottleneck is the
precomputation in Algorithms~\ref{alg:cada}--\ref{alg:pdc}, not the
search.

%% ------------------------------------------------------------------
\subsection{Algorithm 4: Tipping-Point Detection (TPD)}
\label{sec:alg-tpd}

When contexts admit a natural ordering (temporal, dose--response,
geographic gradient), we detect \emph{tipping points}: locations where
the mechanism divergence sequence exhibits discontinuous jumps rather
than smooth drift.

\paragraph{Divergence sequence.}
Let contexts be indexed $c_1,\ldots,c_K$ in order.  Define
\[
  d_k=\tfrac{1}{2}\,\delta_{\mathrm{struct}}(c_k,c_{k+1})
  +\tfrac{1}{2}\,\delta_{\mathrm{param}}(c_k,c_{k+1}),
  \quad k=1,\ldots,K{-}1,
\]
where $\delta_{\mathrm{struct}}$ is the structural divergence from
\textsc{CADA} and
$\delta_{\mathrm{param}}=|S_{\mathrm{shared}}|^{-1}
\sum_{e\in S_{\mathrm{shared}}}d_\mu(\mu_e^{c_k},\mu_e^{c_{k+1}})$.

\paragraph{Cost function.}
The segment cost for contexts $s,\ldots,t$ is
\begin{equation}\label{eq:tpd-cost}
  \mathcal{C}(s,t)=\sum_{k=s}^{t-1}d_k
  +\gamma\,\mathrm{Var}(d_s,\ldots,d_{t-1})\cdot(t-s),
\end{equation}
penalising both total divergence and within-segment heterogeneity; $\gamma>0$ is a regularisation parameter.

\begin{algorithm}[t]
\caption{\textsc{TPD}: Tipping-Point Detection via PELT}
\label{alg:tpd}
\begin{algorithmic}[1]
\Require Ordered DAGs $\{\hat{G}_{c_k}\}_{k=1}^K$;
  parameterisations $\{\Theta_{c_k}\}$;
  statistics $\{\Sigma_{c_k}\}$;
  penalty $\lambda=C_0\log K$;
  min.\ segment length $\ell_{\min}=2$;
  level $\alpha$; permutations $B_{\mathrm{perm}}=999$
\Ensure Validated tipping points with $p$-values,
  effect sizes, and affected mechanisms
\Statex
\Statex \textbf{Step 1: Mechanism Divergence Sequence}
\For{$k=1,\ldots,K{-}1$}
  \State $d_k\gets\textsc{MechDiv}(\hat{G}_{c_k},\hat{G}_{c_{k+1}},
         \Theta_{c_k},\Theta_{c_{k+1}},\Sigma_{c_k},\Sigma_{c_{k+1}})$
\EndFor
\Statex
\Statex \textbf{Step 2: PELT Dynamic Programming}
\State $F[0]\gets -\lambda$;\;
       $\mathrm{cp}[0]\gets\mathrm{nil}$;\;
       $R\gets\{0\}$
\For{$t=1,\ldots,K$}
  \State $(s^*,F[t])\gets
         \displaystyle\arg\min_{s\in R,\;t-s\ge\ell_{\min}}
         \bigl\{F[s]+\mathcal{C}(s{+}1,t)+\lambda\bigr\}$
  \State $\mathrm{cp}[t]\gets s^*$
  \State $R\gets\{s\in R:F[s]+\mathcal{C}(s{+}1,t)\le F[t]\}
         \cup\{t\}$
  \Comment{PELT prune}
\EndFor
\Statex
\Statex \textbf{Step 3: Backtrack}
\State $\hat{T}\gets\emptyset$;\; $t\gets K$
\While{$\mathrm{cp}[t]>0$}
  $\hat{T}\gets\hat{T}\cup\{\mathrm{cp}[t]\}$;\;
  $t\gets\mathrm{cp}[t]$
\EndWhile
\Statex
\Statex \textbf{Step 4: Permutation Validation}
\For{each $\hat{\tau}\in\hat{T}$}
  \State $L\gets\{c_{\hat{\tau}-\ell_{\min}+1},\ldots,c_{\hat{\tau}}\}$;\;
         $R\gets\{c_{\hat{\tau}+1},\ldots,c_{\hat{\tau}+\ell_{\min}}\}$
  \State $d_{\mathrm{obs}}\gets
         \textsc{MechDiv}(\text{pool}(L),\text{pool}(R))$
  \For{$b=1,\ldots,B_{\mathrm{perm}}$}
    \State $\sigma\gets$ random permutation of $L\cup R$
    \State $d_b\gets\textsc{MechDiv}
           (\text{pool}(\sigma_{1:|L|}),
           \text{pool}(\sigma_{|L|+1:}))$
  \EndFor
  \State $p_{\hat{\tau}}\gets
         (1+|\{b:d_b\ge d_{\mathrm{obs}}\}|)
         /(B_{\mathrm{perm}}+1)$
  \State $\mathrm{Cohen}_{\hat{\tau}}\gets
         (d_{\mathrm{obs}}-\bar{d}_{\mathrm{null}})
         /\hat{\sigma}_{\mathrm{null}}$
\EndFor
\State Apply Benjamini--Hochberg at level $\alpha$;
       retain significant $\hat{\tau}$
\Statex
\Statex \textbf{Step 5: Mechanism Attribution}
\For{each validated $\hat{\tau}$}
  \For{$i=1,\ldots,n$}
    \State Compare $\mathrm{Pa}_{c_{\hat{\tau}}}(X_i)$ vs.\
           $\mathrm{Pa}_{c_{\hat{\tau}+1}}(X_i)$ (structural);
           compute $d_\mu$ (parametric)
  \EndFor
  \State Record affected mechanisms and change types
\EndFor
\State \Return validated tipping points with metadata
\end{algorithmic}
\end{algorithm}

\paragraph{PELT pruning.}
The key to $O(K)$ expected-time segmentation is the pruning rule in
Step~2: a candidate $s$ is discarded if its cost-to-here plus the
segment cost to $t$ exceeds $F[t]$, since $s$ can never beat the
current optimum for any future endpoint.  Under the piecewise-Lipschitz
assumption~(A6), the number of non-pruned candidates remains bounded in
expectation~\citep{killick2012optimal}.

\paragraph{False-positive control.}
The permutation test in Step~4 has exact Type~I error $\alpha$ by
exchangeability under the null.  Benjamini--Hochberg correction across
multiple changepoints controls the FDR at level $\alpha$, avoiding the
conservatism of Bonferroni for problems with many candidate tipping
points.

\paragraph{Complexity.}
Divergence matrix: $O(K^2 n d_{\max})$ (or $O(K n d_{\max})$ if only
consecutive pairs are needed).  PELT: $O(K)$ expected, $O(K^2)$ worst
case.  Permutation validation: $O(B_{\mathrm{perm}}\,|\hat{T}|\,n\,d_{\max})$.
Overall:
\[
  T_{\textsc{TPD}}=O(K^2 n d_{\max}+B_{\mathrm{perm}}\,K\,n\,d_{\max}).
\]

%% ------------------------------------------------------------------
\subsection{Algorithm 5: Robustness Certificate Generation (RCG)}
\label{sec:alg-rcg}

For mechanisms classified as \textsc{invariant} by \textsc{PDC}, we
issue formal robustness certificates stating that the mechanism is
$(\tau,\beta)$-invariant across a context set $C$
(Definition~\ref{def:certificate}).

\paragraph{Certificate types.}
\begin{itemize}
\item \textsc{Strong Invariance}: parent set identical across $C$ and
  $\mathrm{UCB}_{\max d_\mu}<0.01$.
\item \textsc{Parametric Stability}: parent set identical and
  $\mathrm{UCB}_{\max d_\mu}\le\tau$.
\item \textsc{Structural Stability}: parent set differs by at most one
  edge, parametric distances bounded.
\item \textsc{Cannot Issue}: structural or parametric plasticity
  exceeds tolerance.
\end{itemize}

\begin{algorithm}[t]
\caption{\textsc{RCG}: Robustness Certificate Generation}
\label{alg:rcg}
\begin{algorithmic}[1]
\Require Aligned DAGs for context set $C$;
  parameterisations $\{\Theta_c\}_{c\in C}$;
  statistics $\{\Sigma_c\}_{c\in C}$;
  variable index $i$; tolerance $\tau$;
  confidence level $\beta$;
  stability rounds $R_{\mathrm{stab}}{=}100$;
  subsample ratio $\rho{=}0.5$;
  bootstrap replicates $B{=}1000$
\Ensure Certificate $\mathcal{R}_i=(\kappa_i,\delta_i,n_{\min},\alpha)$
  or \textsc{Cannot Issue}
\Statex
\Statex \textbf{Step 1: Structural Invariance via Stability Selection}
\State $P_c\gets\mathrm{Pa}_c(X_i)$ for each $c\in C$
\If{$|\mathrm{Unique}(\{P_c\})|>1$}
  \For{$r=1,\ldots,R_{\mathrm{stab}}$}
    \For{each $c\in C$}
      \State $\hat{G}_c^{(r)}\gets\textsc{EstimateDAG}
             (\text{subsample}(\Sigma_c,\rho))$
      \State $P_c^{(r)}\gets\mathrm{Pa}(\hat{G}_c^{(r)},X_i)$
    \EndFor
    \State $\mathrm{agree}^{(r)}\gets\mathbb{1}
           [|\mathrm{Unique}(\{P_c^{(r)}\}_{c\in C})|=1]$
  \EndFor
  \State $\hat{\pi}_{\mathrm{sel}}\gets
         R_{\mathrm{stab}}^{-1}\sum_r\mathrm{agree}^{(r)}$
         \Comment{selection probability}
  \If{$\hat{\pi}_{\mathrm{sel}}<0.5$}
    \State \Return \textsc{Cannot Issue}
           (structural difference robust)
  \EndIf
\EndIf
\Statex
\Statex \textbf{Step 2: Parametric Stability Check}
\State $P^*\gets\bigcap_{c\in C}P_c$
  \Comment{common parent set}
\If{$P^*=\emptyset$ \textbf{and}
    $|\mathrm{Unique}(\{P_c\})|>1$}
  \State \Return \textsc{Cannot Issue}
         (no common parents)
\EndIf
\State $d_{\max}^{\mathrm{obs}}\gets
       \max_{a,b\in C}d_\mu(\mu_i^a,\mu_i^b\mid P^*)$
\Statex
\Statex \textbf{Step 3: Bootstrap Upper Confidence Bound}
\For{$b=1,\ldots,B$}
  \For{each pair $\{a,c\}\subseteq C$}
    \State $\Theta_a^{(b)}\gets\textsc{BootParams}(\Theta_a,\Sigma_a)$;\;
           $\Theta_c^{(b)}\gets\textsc{BootParams}(\Theta_c,\Sigma_c)$
    \State $d_{ac}^{(b)}\gets d_\mu(\mu_i^{a(b)},\mu_i^{c(b)}\mid P^*)$
  \EndFor
  \State $D^{(b)}\gets\max_{a,c}d_{ac}^{(b)}$
\EndFor
\State $\mathrm{UCB}\gets\textsc{Quantile}(\{D^{(b)}\}_{b=1}^B,1{-}\beta)$
\Statex
\Statex \textbf{Step 4: Certificate Decision}
\If{$\mathrm{UCB}\le\tau$}
  \State $\kappa\gets\begin{cases}
         \textsc{Strong Invariance}&\text{if }\mathrm{UCB}<0.01,\\
         \textsc{Parametric Stability}&\text{otherwise}
         \end{cases}$
  \State $\delta_i\gets\tau-\mathrm{UCB}$
         \Comment{robustness margin}
  \State \Return $\mathcal{R}_i=(\kappa,\delta_i,
         \min_c N_c,\beta)$
\Else
  \State \Return \textsc{Cannot Issue}
         ($\mathrm{UCB}=\mathrm{UCB}>\tau$)
\EndIf
\end{algorithmic}
\end{algorithm}

\paragraph{Stability selection.}
Step~1 uses stability selection~\citep{meinshausen2010stability} with
50\,\% subsampling over 100 rounds.  The selection probability
$\hat{\pi}_{\mathrm{sel}}$ measures how often the parent set is
structurally identical across all contexts under data perturbation.  By
the Shah--B\"uhlmann bound~\citep{shah2013variable}, the probability of
falsely declaring structural invariance when the true parent sets differ
is controlled at level $\le 1/(2\hat{\pi}_{\mathrm{sel}}-1)$ per
mechanism.

\paragraph{Bootstrap UCB.}
In Step~3, the parametric bootstrap perturbs regression coefficients by
their standard errors and residual variances via $\chi^2$
resampling.  The $(1{-}\beta)$-quantile of the bootstrap
$\max d_\mu$ provides the upper confidence bound.  Consistency
of the parametric bootstrap~\citep{efron1979bootstrap} ensures
$\mathrm{UCB}\to Q_{1-\beta}(\max d_\mu^{\mathrm{true}})$ at rate
$O(N_{\min}^{-1/2})$.

\paragraph{Certificate guarantees.}
A certificate $\mathcal{R}_i=(\kappa_i,\delta_i,n_{\min},\beta)$
guarantees (Theorem~\ref{thm:certificate-soundness}):
\[
  \Pr\!\bigl(\psi_{P,\mathrm{true}}^{(i)}>\tau
  \mid\text{certificate issued}\bigr)\le\beta.
\]
The robustness margin $\delta_i=\tau-\mathrm{UCB}$ quantifies how much
additional parametric variation the certificate can absorb before
revocation.

\paragraph{Validity conditions.}
Every certificate records its assumptions: (i)~causal sufficiency,
(ii)~i.i.d.\ sampling within each context, (iii)~no unmeasured
confounders within the measured set, (iv)~Gaussian conditional
assumption (for the closed-form $\mathrm{JSD}$).  Violation of any
condition invalidates the certificate.  Small-sample warnings are issued
when $\min_c N_c<200$.

\paragraph{Complexity.}
Structural check: $O\!\bigl(R_{\mathrm{stab}}\,|C|\,
T_{\mathrm{DAG}}\bigr)$, where $T_{\mathrm{DAG}}$ is the cost of
re-estimating a DAG (linear in $|C|$ since each round independently
re-estimates each context).
Parametric check: $O\!\bigl(\binom{|C|}{2}\,B\,d_{\max}\bigr)$.
Per mechanism:
\[
  T_{\textsc{RCG}}
  =O\!\bigl(R_{\mathrm{stab}}\,|C|\,T_{\mathrm{DAG}}
  +|C|^2\,B\,d_{\max}\bigr).
\]
The structural bootstrap dominates; with $|C|{=}10$,
$R_{\mathrm{stab}}{=}100$, and $T_{\mathrm{DAG}}{=}O(n^2\log n)$ for
GES, this costs $\approx 4.5\times10^4\cdot T_{\mathrm{DAG}}$ per
mechanism.  Since the structural check is skipped when parent sets agree
(which is the common case for invariant mechanisms), the practical cost
is typically only the fast parametric bootstrap.


%% ============================================================
%%  Section 5: Theoretical Analysis
%% ============================================================

\section{Theoretical Analysis}\label{sec:theory}

This section establishes the formal guarantees underlying the
Causal-Plasticity Atlas.  All results use the notation of
Section~\ref{sec:framework} and rely on Assumptions~A1--A6.

\subsection{Classification Correctness}

\begin{theorem}[Classification Correctness]\label{thm:classification}
Let $\mu_i$ be a causal mechanism with true plasticity descriptor
$\psi_i=(\psi_S,\psi_P,\psi_E,\psi_{CS})\in\mathbb{R}^4$, and let
$\hat\psi_i$ be the estimated descriptor computed from $N_c\ge n_{\min}$
samples per context across $K$ contexts, where
\[
  n_{\min}\ge\frac{C\,d_{\max}^2\,\sigma^2}{\varepsilon^2}
  \log\!\Bigl(\frac{nK}{\alpha}\Bigr)
\]
for a universal constant $C>0$.  Let $\tau_S,\tau_P,\tau_E$ be the
classification thresholds.  If the \emph{separation condition} holds,
\[
  \min\bigl(|\psi_S-\tau_S|,\;|\psi_P-\tau_P|,\;|\psi_E-\tau_E|\bigr)
  >2\varepsilon,
\]
then
$\Pr\!\bigl[\textsc{Classify}(\hat\psi_i)
=\textsc{Classify}(\psi_i)\bigr]\ge 1-\alpha$.
\end{theorem}

\begin{proof}
We show that each component of $\hat\psi_i$ concentrates within
$\varepsilon$ of its population value, then apply the separation
condition.

\medskip\noindent\textbf{Step 1: Structural plasticity.}
The estimator $\hat\psi_S$ is the square root of the plug-in
multi-distribution JSD over binary parent indicators
$\mathbf{b}_1^{(i)},\ldots,\mathbf{b}_K^{(i)}$.  Each indicator
$\mathbf{b}_c^{(i)}[j]=\mathbb{1}[j\in\widehat{\mathrm{Pa}}_c(X_i)]$
is a Bernoulli random variable determined by $N_c$ samples via the
conditional-independence test in the causal discovery algorithm.  Under
Assumptions~A1--A2 (Markov + faithfulness) and A4 ($N_c\ge n_{\min}$),
the parent indicator is correctly estimated with probability
$\ge 1-\alpha/(nK)$ by the union bound over $n$ potential parents and $K$
contexts, using the standard Neyman--Pearson bound on partial-correlation
tests~\citep{kalisch2007estimating}.

Conditional on all indicators being correct, $\hat\psi_S=\psi_S$ exactly.
By McDiarmid's inequality, changing a single sample affects at most one
parent indicator in one context, which perturbs the JSD by at most
$O(1/(Kn))$.  Hence
$\Pr[|\hat\psi_S-\psi_S|>\varepsilon]\le
2\exp(-2\varepsilon^2 K^2 n^2/K)\le\alpha/(4n)$
for $N_c\ge n_{\min}$.

\medskip\noindent\textbf{Step 2: Parametric plasticity.}
Within each parent-set group, $\hat\psi_P$ is an average of pairwise
$\sqrt{\mathrm{JSD}}$ distances between estimated conditional
distributions.  For Gaussian linear SEMs, regression coefficients
$\hat\beta_c^{(i)}$ are sub-Gaussian with parameter
$\sigma^2 d_{\max}/N_c$~\citep{wainwright2019high}.  Since
$\sqrt{\mathrm{JSD}}$ is Lipschitz in the distribution parameters with
constant $L=O(1/\sigma)$~\citep{endres2003new}, the deviation of each
pairwise distance is bounded:
\[
  \bigl|\hat{d}_\mu(\mu_i^a,\mu_i^b)-d_\mu(\mu_i^a,\mu_i^b)\bigr|
  \le O\!\Bigl(\frac{d_{\max}\,\sigma}{\sqrt{N_{\min}}}\Bigr)
  \le\varepsilon
\]
with probability $\ge 1-\alpha/(4nK^2)$, by a union bound over the
$\binom{K}{2}$ pairs.  The average inherits the same rate.

\medskip\noindent\textbf{Step 3: Emergence.}
The emergence score $\hat\psi_E=1-\min_c|\widehat{\mathrm{MB}}_c(X_i)|
/(\max_c|\widehat{\mathrm{MB}}_c(X_i)|+1)$ depends on Markov-blanket
estimation.  Under the neighbourhood stability condition of
\citet{meinshausen2006high}, $\widehat{\mathrm{MB}}_c(X_i)
=\mathrm{MB}_c(X_i)$ with probability $\ge 1-\alpha/(4nK)$ when
$N_c=\Omega(d_{\max}\log(n)/\varepsilon^2)$.  Conditional on
correct Markov blankets, $\hat\psi_E=\psi_E$.

\medskip\noindent\textbf{Step 4: Context sensitivity.}
$\hat\psi_{CS}$ is the Spearman rank correlation between context
distances and mechanism distances over $O(K^2)$ pairs.  By the
asymptotic normality of the sample Spearman coefficient
\citep{hoeffding1948central}, $|\hat\rho_S-\rho_S|=O_P(K^{-1})$,
which converges faster than the parametric component (Step~2) and
therefore does not increase the overall sample requirement.

\medskip\noindent\textbf{Step 5: Union bound.}
By a union bound over the four components and all $n$ mechanisms:
\[
  \Pr\!\bigl[\|\hat\psi_i-\psi_i\|_\infty>\varepsilon
  \text{ for some }i\bigr]
  \le 4n\cdot\frac{\alpha}{4n}=\alpha.
\]
Under the separation condition, $\|\hat\psi_i-\psi_i\|_\infty\le
\varepsilon$ implies $|\hat\psi_S-\tau_S|>\varepsilon>0$ (and similarly
for $\tau_P,\tau_E$), so $\hat\psi_S$ and $\psi_S$ lie on the same side
of $\tau_S$.  The hierarchical classification rule
(Definition~\ref{def:plasticity-spectrum}) then yields the same label.
\end{proof}

%% ------------------------------------------------------------------
\subsection{Certificate Soundness (Parametric--Structural)}

\begin{theorem}[Certificate Soundness (Parametric--Structural)]\label{thm:certificate-soundness}
Let $\mathcal{R}_i=(\kappa_i,\delta_i,n_{\min},\beta)$ be a robustness
certificate issued by \textsc{RCG} (Algorithm~\ref{alg:rcg}) for
mechanism $\mu_i$, asserting $(\tau,\beta)$-invariance across context
set $C$.  Under Assumptions A1--A5 (Markov, faithfulness, ICM,
sufficient samples, sub-Gaussian variables), with $N_c\ge n_{\min}$ in
each context:
\[
  \Pr\!\bigl[\psi_{P}^{(i),\mathrm{true}}>\tau
  \;\big|\;\text{certificate issued}\bigr]\le\beta.
\]
\end{theorem}

\begin{proof}
The proof proceeds in three steps corresponding to the three
components of \textsc{RCG}.

\medskip\noindent\textbf{Step 1: Structural invariance guarantee.}
The stability-selection procedure (Step~1 of Algorithm~\ref{alg:rcg})
subsamples 50\,\% of data in each context and re-estimates the DAG
across $R_{\mathrm{stab}}=100$ rounds.  The selection probability
$\hat\pi_{\mathrm{sel}}$ is the fraction of rounds in which all
contexts yield identical parent sets for $X_i$.

By \citet{shah2013variable}, Theorem~1, the expected number of
falsely selected variables (i.e., variables whose parent sets appear
invariant when they truly differ) is bounded by
$q^2/(2\hat\pi_{\mathrm{sel}}-1)$ where $q$ is the number of candidate
parents.  Setting $\hat\pi_{\mathrm{sel}}\ge 0.5$ as our
acceptance threshold ensures the per-family-error-rate for structural
invariance is controlled.

Conditional on structural invariance being correctly declared, all
contexts share the same parent set $P^*$.

\medskip\noindent\textbf{Step 2: Bootstrap consistency.}
The parametric bootstrap (Step~3 of Algorithm~\ref{alg:rcg}) generates
$B$ replicates of $\max_{a,b\in C}d_\mu(\mu_i^a,\mu_i^b\mid P^*)$ by
perturbing regression coefficients according to their standard errors
and residual variances via $\chi^2_{N_c-|P^*|-1}/(N_c-|P^*|-1)$
scaling.

The parametric bootstrap is consistent for smooth functionals of
exponential-family parameters~\citep{efron1979bootstrap}.  The
$\sqrt{\mathrm{JSD}}$ distance between Gaussian conditionals is a smooth
function of $(\beta_c^{(i)},\sigma_c^{2(i)})$, and the maximum over
finitely many pairs is Lipschitz.  Therefore the bootstrap distribution
of $D^{(b)}=\max_{a,c}d_{ac}^{(b)}$ converges to the true sampling
distribution of $\max_{a,c}d_\mu(\hat\mu_i^a,\hat\mu_i^c)$ at rate
$O(N_{\min}^{-1/2})$ in the Kolmogorov distance.

\medskip\noindent\textbf{Step 3: Coverage guarantee.}
The UCB is defined as $\mathrm{UCB}=Q_{1-\beta}(\{D^{(b)}\}_{b=1}^B)$.
By bootstrap consistency:
\[
  \mathrm{UCB}\xrightarrow{N_{\min}\to\infty}
  Q_{1-\beta}\!\bigl(\max_{a,c}d_\mu^{\mathrm{true}}
  (\mu_i^a,\mu_i^c)\bigr).
\]
By definition of the quantile, the probability that the true maximum
exceeds the UCB is at most $\beta+o(1)$.  The certificate is issued
only when $\mathrm{UCB}\le\tau$, so
\[
  \Pr[\psi_P^{(i),\mathrm{true}}>\tau\mid\text{issued}]
  \le\Pr[\max d_\mu^{\mathrm{true}}>\mathrm{UCB}]
  \le\beta+o(1).
\]
For $N_{\min}$ satisfying Assumption~A4, the $o(1)$ term is negligible.
\end{proof}

\begin{remark}
Theorem~\ref{thm:certificate-soundness} covers only the parametric
plasticity component~$\psi_P$.  Analogous soundness guarantees for the
emergence certificate~($\psi_E$) and the context-sensitivity
certificate~($\psi_{CS}$) remain open problems: $\psi_E$ depends on
consistent Markov-blanket estimation across contexts, and $\psi_{CS}$
inherits the convergence properties of the Spearman rank correlation
estimator, both of which require additional regularity conditions beyond
Assumptions~A1--A5.
\end{remark}

%% ------------------------------------------------------------------
\subsection{Sample Complexity Bounds}

\begin{theorem}[Sample Complexity]\label{thm:sample-complexity}
To estimate the plasticity descriptor $\psi_i\in\mathbb{R}^4$ to within
$\varepsilon$ in $\|\cdot\|_\infty$ with probability $\ge 1-\alpha$, it
suffices to have
\[
  N_{\min}=\Omega\!\Bigl(\frac{d_{\max}^2\,\sigma^2}{\varepsilon^2}
  \log\frac{nK}{\alpha}\Bigr)
\]
samples per context.  This rate is minimax optimal up to logarithmic
factors for the Gaussian linear SEM class.
\end{theorem}

\begin{proof}[Proof sketch]
\textbf{Upper bound (achievability).}
We analyse each component separately and take the maximum sample
requirement.

\emph{$\psi_S$ (structural):}
The JSD over parent indicators requires each indicator
$\mathbb{1}[j\in\mathrm{Pa}_c(X_i)]$ to be correctly estimated.  By
Hoeffding's inequality on the partial-correlation test statistic, this
requires $N=\Omega(\varepsilon^{-2}\log n)$ per context per indicator.
Since $\sqrt{\mathrm{JSD}}$ is $1$-Lipschitz in the indicator means,
the rate transfers to $\psi_S$.

\emph{$\psi_P$ (parametric):}
Estimating $\beta_c^{(i)}\in\mathbb{R}^{d_{\max}}$ to precision
$\varepsilon$ requires
$N=\Omega(d_{\max}\sigma^2/\varepsilon^2)$ by standard
regression theory.  The Lipschitz constant of $\sqrt{\mathrm{JSD}}$
w.r.t.\ the regression parameters is $L=O(1/\sigma)$, so the descriptor
error is $O(\sigma\cdot L\cdot\sqrt{d_{\max}/N})=O(\sqrt{d_{\max}/N})$.
Setting this to $\varepsilon$ gives
$N=\Omega(d_{\max}/\varepsilon^2)$.

\emph{$\psi_E$ (emergence):}
Correct Markov-blanket identification requires
$N=\Omega(d_{\max}\log(n)/\varepsilon^2)$ by the coherence condition
of~\citet{meinshausen2006high}.

\emph{$\psi_{CS}$ (context sensitivity):}
The Spearman rank correlation over $\binom{K}{2}$ pairs inherits the
mechanism-distance estimation rate from $\psi_P$; additionally, the
sample Spearman coefficient converges at rate $O(K^{-1})$, which is
dominated by the parametric component for any fixed~$K$.

Taking the maximum over components and applying a union bound over $n$
mechanisms and $K$ contexts yields $N_{\min}$ as stated.

\medskip\noindent\textbf{Lower bound (minimax).}
Consider the class of Gaussian linear SEMs on $n$ variables with maximum
in-degree $d_{\max}$ and noise variance $\sigma^2$.  Two SEMs that
differ in a single edge weight by $\varepsilon$ in a single context
yield descriptors differing by $\Theta(\varepsilon)$ in $\psi_P$.
By Fano's inequality applied to the hypothesis-testing problem of
distinguishing these two SEMs from $N$ samples:
\[
  N\ge\frac{c\,d_{\max}\,\sigma^2}{\varepsilon^2}
  \bigl(1-o(1)\bigr),
\]
matching the upper bound up to logarithmic factors.
\end{proof}

%% ------------------------------------------------------------------
\subsection{Tipping-Point Consistency}

\begin{theorem}[Tipping-Point Detection Consistency]
\label{thm:tipping-point}
Let $\{G_k,\theta_k\}_{k=1}^K$ be a sequence of causal models satisfying
Assumption~A6 (piecewise Lipschitz) with finitely many true tipping
points $T^*=\{\tau_1^*,\ldots,\tau_m^*\}$.  Let $\hat{T}$ be the
output of \textsc{TPD} (Algorithm~\ref{alg:tpd}) with penalty
$\lambda=C_0\log K$.  Then:
\begin{enumerate}[label=(\roman*)]
\item \emph{Consistency}: As $K\to\infty$ and $N_{\min}\to\infty$,
  $\Pr[|\hat{T}|=|T^*|]\to 1$ and
  $\max_{\tau^*\in T^*}\min_{\hat\tau\in\hat{T}}
  |\hat\tau-\tau^*|\to 0$.
\item \emph{No false positives}: For any segment $[s,t]$ containing
  no true tipping point, the probability of detecting a spurious
  changepoint within $[s,t]$ is at most $\alpha$.
\item \emph{Localisation rate}:
  $|\hat\tau-\tau^*|=O_p(1/K)$ for each true tipping point $\tau^*$.
\end{enumerate}
\end{theorem}

\begin{proof}
\textbf{Part (i): Number and location consistency.}
Define the mechanism divergence sequence $d_k=\sqrt{\mathrm{JSD}}
(\text{mechanisms}_k,\text{mechanisms}_{k+1})$ for $k=1,\ldots,K{-}1$.
Under A6, $d_k$ is piecewise smooth with jumps of magnitude
$\Delta_\tau\ge\Delta_{\min}>0$ at each true tipping point.

Under sufficient samples ($N_{\min}\to\infty$), the estimation noise in
$d_k$ is $\hat{d}_k-d_k=O_p(N_{\min}^{-1/2})$, which vanishes relative
to $\Delta_{\min}$.  The PELT algorithm with BIC-type penalty
$\lambda=C_0\log K$ is known to consistently estimate the number and
locations of changepoints in piecewise-constant signals when the
signal-to-noise ratio diverges~\citep{killick2012optimal}.

More precisely, let $F_K(m)$ denote the minimum segmentation cost with
$m$ changepoints.  For $m<|T^*|$, a segment spanning a true tipping
point incurs cost at least $\Omega(\Delta_{\min}^2 K)$ above the
correctly segmented cost.  For $m>|T^*|$, the additional penalty
$\lambda$ exceeds the within-segment cost reduction, which is
$O(\sigma_{\mathrm{est}}^2)=o(\log K)$ under the Lipschitz condition.
Hence $\hat{m}=|T^*|$ with probability tending to~1.

Conditional on $|\hat{T}|=|T^*|$, the location of each estimated
changepoint minimises the local cost, which is achieved within
$O_p(1/K)$ of the true location by the convexity of the segment cost
around the true changepoint (since $\mathcal{C}(s,t)$ has a unique
minimum at the true segmentation boundary when the signal has a jump
discontinuity).

\medskip\noindent\textbf{Part (ii): False-positive control.}
The permutation validation in Step~4 tests the null hypothesis $H_0$:
the mechanism descriptors on the two sides of $\hat\tau$ are
exchangeable.  Under $H_0$, all $\binom{|L|+|R|}{|L|}$ allocations are
equally likely, so the permutation $p$-value has exact Type~I error
$\alpha$: $\Pr[p_{\hat\tau}\le\alpha\mid H_0]=\alpha$.

When multiple changepoints are tested, the Benjamini--Hochberg procedure
at level $\alpha$ controls the false discovery rate:
$\mathrm{FDR}=\mathbb{E}[V/\max(R,1)]\le\alpha$, where $V$ is the
number of false rejections and $R$ the total number of
rejections~\citep{benjamini1995controlling}.

\medskip\noindent\textbf{Part (iii): Localisation rate.}
At a true tipping point $\tau^*$, the divergence sequence has a jump of
size $\Delta_\tau$.  The CUSUM statistic
$S_t=\sum_{k=1}^{t}(d_k-\bar{d})$ has a peak within $O(1)$ indices of
$\tau^*$ in the discrete sequence, corresponding to $O(1/K)$ in the
normalised context parameter.  The estimation noise adds a perturbation
of $O_p(\sqrt{K}/N_{\min}^{1/2})$ to the CUSUM, which is $o(1)$ as
$N_{\min}/K\to\infty$.  Hence
$|\hat\tau-\tau^*|=O_p(1/K)$.
\end{proof}

%% ------------------------------------------------------------------
\subsection{Atlas Completeness}

\begin{theorem}[Atlas Completeness Under Ergodic Search]
\label{thm:atlas-completeness}
Let $\mathcal{A}_T$ be the QD archive after $T$ iterations of
\textsc{CD-QD} (Algorithm~\ref{alg:cdqd}) with CVT tessellation into
$L$ cells.  Assume the mutation operator is ergodic over the genome
space $\mathcal{G}=2^{[K]}\times 2^{[n]}\times\Phi$ with minimum
single-step transition probability $p_{\min}>0$.  Then:
\begin{enumerate}[label=(\roman*)]
\item \emph{Coverage convergence}:
  $\mathbb{E}[|\{k:\mathcal{A}_T[k]\ne\emptyset\}|]/L\to 1$ as
  $T\to\infty$ for all reachable cells.
\item \emph{Rate}: The expected iterations to $(1{-}\varepsilon)$-coverage
  of the $L_r\le L$ reachable cells is
  $O(L_r\log(L_r/\varepsilon)/p_{\min})$.
\item \emph{Quality monotonicity}: For each cell $k$,
  $q(\mathcal{A}_T[k])$ is non-decreasing in $T$.
\end{enumerate}
\end{theorem}

\begin{proof}[Proof sketch]
\textbf{(iii)} is immediate from the elitist replacement rule: a cell's
occupant is replaced only by a strictly higher-quality genome.

\textbf{(i):} The mutation operator's ergodicity means the induced
Markov chain on $\mathcal{G}$ is irreducible and positive recurrent
(since $\mathcal{G}$ is finite).  Every genome $g\in\mathcal{G}$ is
visited infinitely often with probability~1.  The behaviour mapping
$g\mapsto\textsc{CVTCell}(\mathbf{b}(g))$ is many-to-one, so every
reachable cell $k$ (i.e., $k\in\{$\textsc{CVTCell}$(\mathbf{b}(g)):
g\in\mathcal{G}\}$) is visited infinitely often.  Upon first visit, the
cell transitions from empty to occupied and remains occupied thereafter
by quality monotonicity.

\textbf{(ii):} This follows from a coupon-collector argument.  Let
$L_r$ be the number of reachable cells.  Suppose cell $k$ has
stationary visitation probability $\pi_k>0$ under the ergodic
chain.  The minimum visitation probability is $\pi_{\min}\ge
p_{\min}^{D}$ where $D$ is the mixing time.  The expected time to
visit all $L_r$ cells is the coupon-collector time with non-uniform
probabilities:
\[
  \mathbb{E}[T_{\mathrm{cover}}]
  =O\!\Bigl(\frac{L_r}{\pi_{\min}}
  \sum_{k=1}^{L_r}\frac{1}{k}\Bigr)
  =O\!\Bigl(\frac{L_r\log L_r}{\pi_{\min}}\Bigr).
\]
The curiosity bonus (inverse visit count) accelerates coverage by
preferentially selecting under-visited cells, analogous to a UCB
strategy.  By the analysis of~\citet{auer2002finite}, the
curiosity-weighted selection achieves regret $O(\sqrt{L_r T\log T})$,
yielding $(1{-}\varepsilon)$-coverage in
$O(L_r\log(L_r/\varepsilon)/p_{\min})$ iterations.
\end{proof}

%% ------------------------------------------------------------------
\subsection{Coverage Rate Bound}

\begin{theorem}[Coverage Rate Bound]\label{thm:coverage-rate}
Under the conditions of Theorem~\ref{thm:atlas-completeness}, with
curiosity weight $\mu>0$ and CVT cell count $L$, the expected fraction
of reachable cells occupied after $T$ iterations satisfies:
\[
  \mathbb{E}\!\left[\frac{|\mathcal{A}_T|}{L_r}\right]
  \ge 1-L_r\exp\!\left(-\frac{T\,p_{\min}}{2L_r}\right).
\]
In particular, $(1{-}\varepsilon)$-coverage is achieved for
$T\ge 2L_r\,p_{\min}^{-1}\log(L_r/\varepsilon)$.
\end{theorem}

\begin{proof}
For each reachable cell $k$, let $T_k$ be the first time cell $k$ is
occupied.  We show $\Pr[T_k>t]\le\exp(-t\,p_{\min}/(2L_r))$ via a
stochastic-dominance argument.

At each iteration, the probability that the selected genome maps to cell
$k$ is at least $p_{\min}/(2L_r)$.  The factor of $2$ accounts for the
worst-case dilution from the curiosity-weighted selection: even when cell
$k$ has the lowest curiosity score (fully visited), the ergodic
base measure ensures visitation probability $\ge p_{\min}$, and the
selection mechanism at most halves this probability (since the
curiosity-weighted distribution is a mixture of $\frac{1}{2}$ uniform $+$
$\frac{1}{2}$ curiosity-proportional, by construction of the weighted
sampling).

Therefore $T_k$ is stochastically dominated by a geometric random
variable with success probability $p_{\min}/(2L_r)$:
\[
  \Pr[T_k>t]\le\Bigl(1-\frac{p_{\min}}{2L_r}\Bigr)^t
  \le\exp\!\Bigl(-\frac{t\,p_{\min}}{2L_r}\Bigr).
\]

The expected number of unoccupied cells after $T$ iterations is
\[
  \mathbb{E}[L_r-|\mathcal{A}_T|]
  =\sum_{k=1}^{L_r}\Pr[T_k>T]
  \le L_r\exp\!\Bigl(-\frac{T\,p_{\min}}{2L_r}\Bigr).
\]
Dividing by $L_r$ yields the stated bound.  Setting the right-hand side
to $\varepsilon$ and solving for $T$ gives
$T\ge 2L_r\,p_{\min}^{-1}\log(L_r/\varepsilon)$.
\end{proof}

%% ------------------------------------------------------------------
\subsection{Hardness Results}

\begin{theorem}[Hardness and Tractability of DAG Alignment]
\label{thm:hardness}
\begin{enumerate}[label=(\alph*)]
\item \emph{NP-hardness}: Computing the optimal DAG alignment
  $d_{\mathrm{DAG}}(G_a,G_b)$ when no anchors are provided is
  NP-hard.
\item \emph{Fixed-parameter tractability}: When the maximum in-degree is
  bounded by $k$ and the number of unanchored variables by $u$, the
  optimal alignment can be computed in
  $O(n^2 d_{\max}^2+u^3)$ time.  For $k\le 5$, $u\le 10$, this is
  $O(n^2)$.
\item \emph{Approximation guarantee}: \textsc{CADA}
  (Algorithm~\ref{alg:cada}) computes a
  $(1{+}\varepsilon)$-approximate alignment in
  $O(n^2 d_{\max}^2+u^3)$ time, where $\varepsilon$ depends on the
  Markov-blanket pruning threshold.
\end{enumerate}
\end{theorem}

\begin{proof}[Reduction sketch]
\textbf{(a)} We reduce from subgraph isomorphism.  Given graphs $H$ and
$G$, construct DAGs $D_H$ and $D_G$ by orienting all edges according to
an arbitrary topological ordering of the vertices (e.g., by vertex
label).  Set $\lambda_{\mathrm{mech}}=0$ (ignore mechanism distances) and
$\lambda_{\mathrm{miss}}=\infty$ (require perfect matching of all
variables).  Then
\[
  d_{\mathrm{DAG}}(D_H,D_G)=0
  \quad\Longleftrightarrow\quad
  \text{$H$ is isomorphic to a subgraph of $G$},
\]
since the alignment $\pi$ becomes an injective vertex map that must
preserve all edges (as $\lambda_{\mathrm{struct}}$ penalises mismatches
and $\lambda_{\mathrm{miss}}=\infty$ penalises unmatched vertices).
Subgraph isomorphism is NP-complete~\citep{cook1971complexity}, so
computing $d_{\mathrm{DAG}}$ is NP-hard in the general (anchorless)
case.

\medskip\noindent\textbf{(b)} With $|U_a|\le u$ unanchored variables
and maximum in-degree $k$:
\begin{enumerate}
\item \emph{Anchor propagation} (Phase~1): $O(n)$.
\item \emph{Candidate generation} (Phase~2): each unanchored variable
  has a Markov blanket of size $O(k^2)$ (parents + children +
  co-parents under bounded in-degree).  The Jaccard overlap computation
  for each of the $u^2$ candidate pairs costs $O(k^2)$, giving
  $O(u^2 k^2)$.
\item \emph{CI scoring} (Phase~3): for each of the $O(u\cdot k^2)$
  surviving candidate pairs (after pruning by the overlap threshold),
  the CI-fingerprint computation involves $O(k^2)$ partial-correlation
  tests, each costing $O(k^2)$.  Total: $O(u\cdot k^4)$.
  Since $k$ is constant, this is $O(u)$.
\item \emph{Hungarian matching} (Phase~4): $O(u^3)$ on a $u\times u$
  matrix.
\item \emph{Edge classification} (Phase~5): $O(nk)$ linear in edges.
\end{enumerate}
Summing: $O(n+u^2 k^2+u\cdot k^4+u^3+nk)=O(n^2 k^2+u^3)$
(since $u\le n$ and the quadratic term from anchor propagation over all
$n$ variables dominates).

\medskip\noindent\textbf{(c)} The approximation guarantee follows from
the Markov-blanket pruning lemma:

\emph{Lemma} (MB Pruning Soundness). Under faithfulness (A2), if
$\pi^*$ is the optimal alignment and $(u_a,\pi^*(u_a))$ is an optimal
pair, then the Jaccard overlap of their anchored Markov blankets
satisfies $J(u_a,\pi^*(u_a))\ge\rho$ for a constant $\rho>0$ depending
on the minimum edge strength and graph sparsity.

\emph{Proof of Lemma.}
If $\pi^*(u_a)=u_b$ is the optimal match, then by definition of the
alignment cost, the CI fingerprints of $u_a$ and $u_b$ agree on most
anchored neighbours.  Under faithfulness, CI agreement implies Markov-blanket
membership agreement.  Since the anchored portion of the Markov blanket
is at least a constant fraction of the full blanket (by the bounded
in-degree assumption and the assumption that most variables are
anchored), the Jaccard overlap is $\ge\rho>0$.

Setting the pruning threshold to $0.3\le\rho$, no optimal pair is
discarded.  The Hungarian algorithm then finds the exact optimum on the
surviving candidates, yielding $\mathrm{cost}(\textsc{CADA})\le
\mathrm{cost}(\pi^*)$.  When the Jaccard pruning threshold~$J_{\min}$
is set below~$\rho$, all optimal pairs survive pruning and the
algorithm recovers the exact optimum.  In general, for threshold
$J_{\min}>\rho$, some near-optimal pairs may be pruned, and the
approximation ratio is $(1{+}\varepsilon)$ where
$\varepsilon\le\max_{(u_a,u_b)\,\text{pruned}}
[\mathbf{S}[u_a,\pi^*(u_a)]-\mathbf{S}[u_a,u_b]]
/\mathrm{cost}(\pi^*)$, which depends on the gap between the pruning
threshold and the minimum Jaccard overlap of optimal pairs.
\end{proof}

%% ------------------------------------------------------------------
\subsection{Structural Stability Under DAG Estimation Errors}

To complement the above results, we state a robustness guarantee showing
that the plasticity descriptor degrades gracefully when the input DAGs
contain estimation errors.

\begin{theorem}[Structural Stability]\label{thm:structural-stability}
Let $\hat{G}_c$ differ from the true DAG $G_c$ by
$\mathrm{SHD}(\hat{G}_c,G_c)\le s$ edges in each context $c$.  Let
$\hat\psi_i$ and $\psi_i$ be the descriptors from estimated and true
DAGs respectively.  Then:
\[
  \|\hat\psi_i-\psi_i\|_\infty
  \le f(s,n,d_{\max})
  =O\!\Bigl(\frac{s\,d_{\max}}{n}+\frac{s}{K}\Bigr),
\]
with component-wise bounds:
\begin{enumerate}[label=(\roman*)]
\item $|\hat\psi_S-\psi_S|\le 2s/(Kn)$,
\item $|\hat\psi_P-\psi_P|
  \le O(s\,d_{\max}\,\sigma\,N_{\min}^{-1/2})$,
\item $|\hat\psi_E-\psi_E|
  \le 2s/(\max_c|\mathrm{MB}_c(X_i)|+1)$.
\end{enumerate}
\end{theorem}

\begin{proof}
\textbf{(i) Structural plasticity.}
Each misestimated edge changes at most $2$ parent indicators (one
addition and one deletion) in at most one context.  The multi-JSD
over $K$ contexts with $n$ binary indicators has Lipschitz constant
$1/(Kn)$ per indicator flip (since changing one indicator in one context
shifts the empirical mean by $1/K$ and the binary entropy by
$O(1/K)$).  With $s$ misestimated edges, at most $2s$ indicators
are flipped, giving $|\hat\psi_S-\psi_S|\le 2s/(Kn)$.

\medskip\noindent\textbf{(ii) Parametric plasticity.}
A misestimated parent changes the conditioning set in one context.
By the omitted-variable bias formula, the regression coefficients
change by $O(\sigma\sqrt{d_{\max}/N_{\min}})$ per omitted or added
parent.  The $\sqrt{\mathrm{JSD}}$ distance has Lipschitz constant
$O(1/\sigma)$ w.r.t.\ regression parameters, so each misestimated
parent perturbs each pairwise distance by
$O(d_{\max}/\sqrt{N_{\min}})$.  With $s$ errors affecting at most $s$
parent sets, the average pairwise distance shifts by
$O(s\,d_{\max}\,\sigma/\sqrt{N_{\min}})$.

\medskip\noindent\textbf{(iii) Emergence.}
Each misestimated edge changes the Markov blanket size by at most $1$
in one context.  The emergence score
$\psi_E=1-\min_c|\mathrm{MB}_c|/(\max_c|\mathrm{MB}_c|+1)$ has
derivative $O(1/(\max|\mathrm{MB}|+1))$ w.r.t.\ a single Markov-blanket
size change.  With $s$ errors, each affecting at most one context,
$|\hat\psi_E-\psi_E|\le 2s/(\max_c|\mathrm{MB}_c|+1)$.

The $\|\cdot\|_\infty$ bound is the maximum of the three components.
\end{proof}

\paragraph{Practical implications.}
For a typical setting with $n=100$, $K=10$, $d_{\max}=5$, and $s=2$
misestimated edges per context, the perturbation bounds are:
$|\Delta\psi_S|\le 0.004$, $|\Delta\psi_P|\le O(10/\sqrt{N_{\min}})$
(negligible for $N_{\min}\ge 500$), and $|\Delta\psi_E|\le 0.1$.  These
are all well within the classification thresholds, confirming that
moderate DAG estimation errors do not invalidate the atlas.


% ==============================================================
% NEW: Advanced Theoretical Results (Sections 5A–5D)
% ==============================================================

\section{Measure-Theoretic Foundations}\label{sec:measure}

The theoretical results of Section~\ref{sec:theory} establish the
correctness and complexity of the CPA pipeline.  In this section and the
three that follow, we deepen the mathematical foundations by proving
structural properties of the mechanism-distance space, establishing
minimax optimality of our sample complexity bounds, revealing
information-geometric structure in the plasticity-descriptor space,
and strengthening our certificate guarantees to uniform (family-wise)
validity.

Throughout Sections~\ref{sec:measure}--\ref{sec:uniform-certs} we
work in the \emph{Gaussian linear SEM} setting unless otherwise
stated: each mechanism $\mu_i^c$ is a conditional Gaussian
$X_i \mid \mathrm{Pa}_i^c \sim \mathcal{N}(\beta_c^{(i)\top}
\mathrm{Pa}_i^c,\; \sigma_c^{2(i)})$.  We note explicitly where
results extend to broader model classes and at what cost.

\subsection{Mechanism-Distance Kernel Properties}\label{sec:kernel}

The mechanism distance $d_\mu$ (Definition~\ref{def:mechanism-distance})
is built from the Jensen--Shannon divergence.  We now show that $d_\mu^2$
induces a \emph{conditionally negative definite} (CND) kernel, enabling
kernel embeddings of mechanism spaces into reproducing kernel Hilbert
spaces.

\begin{definition}[Conditionally Negative Definite Kernel]\label{def:cnd}
A symmetric function $\kappa : \mathcal{X} \times \mathcal{X} \to \R$
is \emph{conditionally negative definite} (CND) if $\kappa(x,x) \geq 0$
for all $x$ and, for every $n \geq 2$, every $x_1,\ldots,x_n \in
\mathcal{X}$, and every $\alpha_1,\ldots,\alpha_n \in \R$ with
$\sum_i \alpha_i = 0$,
\[
  \sum_{i,j=1}^n \alpha_i \alpha_j \,\kappa(x_i,x_j) \;\leq\; 0.
\]
By Schoenberg's theorem, $\kappa$ is CND if and only if
$e^{-t\kappa}$ is positive definite for all $t > 0$.
\end{definition}

\begin{theorem}[Mechanism Distance Kernel]\label{thm:mechanism-kernel}
Let $\mathcal{P}$ be a set of probability measures on $\R^d$, each
absolutely continuous with respect to a common reference measure~$\lambda$.
Define $\kappa(P,Q) = \mathrm{JSD}(P \| Q)$, where $\mathrm{JSD}$ is
the Jensen--Shannon divergence.  Then:
\begin{enumerate}[label=(\alph*)]
  \item $\kappa$ is conditionally negative definite on $\mathcal{P}$.
  \item The Gaussian kernel
    $k_h(P,Q) = \exp\bigl(-\mathrm{JSD}(P\|Q)/h^2\bigr)$ is positive
    definite on $\mathcal{P}$ for every bandwidth $h > 0$.
  \item For conditional mechanisms
    $d_\mu(\mu_i^c,\mu_i^{c'})^2 =
    \mathbb{E}_{Z\sim P_{\mathrm{ref}}}
    [\mathrm{JSD}(\mu_i^c(\cdot\mid Z)\|\mu_i^{c'}(\cdot\mid Z))]$,
    the function $d_\mu^2$ is CND on the space of conditional mechanisms.
\end{enumerate}
\end{theorem}

\begin{proof}
\textbf{Part (a).}
We use the Endres--Schindelin decomposition \citep{endres2003new}.
Write $M = \frac{1}{2}(P+Q)$ and recall
$\mathrm{JSD}(P\|Q) = H(M) - \frac{1}{2}H(P) - \frac{1}{2}H(Q)$,
where $H$ denotes differential entropy.  Define the Hellinger-type
embedding $\phi: P \mapsto \sqrt{dP/d\lambda} \in L^2(\lambda)$.
The squared Hellinger distance satisfies
$H^2(P,Q) = \frac{1}{2}\int(\sqrt{p}-\sqrt{q})^2\,d\lambda
= 1 - A(P,Q)$
where $A(P,Q) = \int\sqrt{pq}\,d\lambda$ is the Bhattacharyya affinity.

By \citet{topsoe2000some}, the Jensen--Shannon divergence admits
the variational representation
\begin{equation}\label{eq:jsd-variational}
  \mathrm{JSD}(P\|Q) = \inf_{R}\,\bigl[\tfrac{1}{2}\KL(P\|R)
  + \tfrac{1}{2}\KL(Q\|R)\bigr],
\end{equation}
where the infimum is achieved at $R=M$.  For any fixed~$R$, the function
$(P,Q) \mapsto \frac{1}{2}\KL(P\|R) + \frac{1}{2}\KL(Q\|R)$ decomposes
as $f(P) + f(Q)$ with $f(P) = \frac{1}{2}\KL(P\|R)$.  Functions of the
form $\kappa(P,Q) = f(P) + f(Q)$ are trivially CND (since
$\sum_{i,j}\alpha_i\alpha_j[f(x_i)+f(x_j)] = 2(\sum_i\alpha_i)
\sum_j\alpha_j f(x_j) = 0$ when $\sum_i\alpha_i=0$).

The infimum of CND functions is not CND in general, so we proceed
differently.  Instead, we use the integral representation due to
\citet{briët2009properties}: for absolutely continuous measures,
\begin{equation}\label{eq:jsd-integral}
  \mathrm{JSD}(P\|Q) = \int_0^1 \frac{1}{t(1-t)}
  \,D_t(P\|Q)\,dt,
\end{equation}
where $D_t(P\|Q) = \int p^t q^{1-t}\,d\lambda - 1$ is related to the
R\'enyi divergence of order $t$.  For each fixed $t \in (0,1)$, the
function $(P,Q) \mapsto 1 - \int p^t q^{1-t}\,d\lambda$ equals
$1 - \langle \phi_t(P), \phi_{1-t}(Q)\rangle_{L^2}$ where
$\phi_s(P) = p^s$.  The negative inner product $-\langle\phi_t(P),
\phi_{1-t}(Q)\rangle$ is CND whenever the embedding maps are
consistent (i.e., $\phi_t(P) = \phi_{1-t}(P)$ at $t=1/2$, which
holds: both equal $\sqrt{p}$).  More directly,
\citet{sejdinovic2013equivalence} prove that $-\int p^{1/2}q^{1/2}
d\lambda = -(1-H^2(P,Q)/2)$ is CND, and the integrand at each~$t$
has the same structure under the $t$-power embedding.

Since the integrand in~\eqref{eq:jsd-integral} is CND for each~$t$
and the integral over $t$ with positive weight $1/(t(1-t))$ preserves
the CND property (CND functions form a convex cone closed under
non-negative integration), we conclude that $\mathrm{JSD}$ is CND.

\textbf{Part (b).}
Immediate from Schoenberg's theorem \citep{schoenberg1938metric}:
if $\kappa$ is CND then $e^{-t\kappa}$ is positive definite for all
$t>0$.  Setting $t = 1/h^2$ gives the result.

\textbf{Part (c).}
Write $d_\mu^2(\mu,\mu') =
\mathbb{E}_{Z}[\mathrm{JSD}(\mu(\cdot|Z)\|\mu'(\cdot|Z))]$.
For each fixed $Z=z$, part~(a) gives that
$\kappa_z(\mu,\mu') = \mathrm{JSD}(\mu(\cdot|z)\|\mu'(\cdot|z))$
is CND.  The expectation $d_\mu^2 = \mathbb{E}_Z[\kappa_Z]$ is a
non-negative integral of CND functions, hence CND (the CND cone
is closed under integration with respect to any probability measure).
\end{proof}

\begin{remark}[Kernel Methods for Mechanism Comparison]
Theorem~\ref{thm:mechanism-kernel} enables the use of kernel methods
(kernel PCA, kernel $k$-means, MMD tests) directly in mechanism space.
For instance, one can test $H_0: \mu_i^c = \mu_i^{c'}$ using an
MMD test with the kernel $k_h$, obtaining a consistent nonparametric
test of mechanism identity without specifying the parametric form of
mechanisms.  The bandwidth~$h$ can be selected by the median heuristic
$h = \mathrm{median}\{d_\mu(\mu_i^c,\mu_i^{c'}) : c \neq c'\}$.
\end{remark}

\subsection{Convergence Rate for Plasticity Estimation}\label{sec:convergence-rate}

We now establish the rate at which the empirical plasticity descriptor
$\hat\psi_n$ converges to the population descriptor~$\psi^*$ as the
per-context sample size~$n$ grows.

\begin{theorem}[Plasticity Estimator Convergence Rate]\label{thm:convergence-rate}
Under Assumptions~A1--A5, in the Gaussian linear SEM setting with
$K$ contexts, $|\bar V|$ variables, maximum in-degree $d_{\max}$,
and residual variance bounded by $\sigma^2$, the empirical plasticity
descriptor~$\hat\psi_n$ computed by Algorithm~PDC satisfies
\begin{equation}\label{eq:convergence-bound}
  \mathbb{E}\bigl[\|\hat\psi_n - \psi^*\|_\infty\bigr]
  \;\leq\;
  \frac{C_1\, d_{\max}\, \sigma \sqrt{\log(|\bar V| K)}}{\sqrt{n}}
  \;+\; \frac{1}{K-1},
\end{equation}
where $C_1 = 8\sqrt{2}$ is an explicit constant.  The first term is
the estimation error (vanishing with~$n$) and the second is a
discretisation bias (vanishing with~$K$).  In particular, for
$n \geq C_1^2 d_{\max}^2 \sigma^2 \log(|\bar V|K)/\eps^2$ and
$K \geq 1+1/\eps$, we have
$\mathbb{E}[\|\hat\psi_n - \psi^*\|_\infty] \leq 2\eps$.
\end{theorem}

\begin{proof}
We bound the sup-norm error by taking the maximum over the four
descriptor components.

\textbf{Component $\psi^{(1)}$ (parametric plasticity).}
For Gaussian conditionals,
$\mathrm{JSD}(\mu_i^c \| \mu_i^{c'})$ has the closed form
\begin{equation}\label{eq:gaussian-jsd}
  \mathrm{JSD}(P\|Q) = \frac{1}{2}\log\frac{\sigma_M^2}
  {\sigma_P\sigma_Q} + \frac{(\mu_P-\mu_Q)^2}{8\sigma_M^2}
  + \frac{\sigma_P^2+\sigma_Q^2}{4\sigma_M^2} - \frac{1}{2},
\end{equation}
where $\sigma_M^2 = \frac{1}{2}(\sigma_P^2+\sigma_Q^2)
+ \frac{1}{4}(\mu_P-\mu_Q)^2$ for the univariate case
(with $\mu_P = \beta_c^{(i)\top}z$, $\sigma_P = \sigma_c^{(i)}$).
The empirical version $\widehat{\mathrm{JSD}}$ replaces $(\hat\beta,
\hat\sigma^2)$ from OLS on $n$ samples.  By the delta method and
sub-Gaussian concentration of OLS estimators under~A5:
\[
  |\hat\beta_c^{(i)} - \beta_c^{(i)*}| \leq
  \frac{2\sigma}{\sqrt{n}} \sqrt{d_{\max}\log(2d_{\max}/\delta)}
\]
with probability $\geq 1-\delta$.  Similarly,
$|\hat\sigma_c^{2(i)} - \sigma_c^{2(i)*}| \leq
4\sigma^2\sqrt{\log(2/\delta)/n}$.

The JSD is Lipschitz in $(\beta,\sigma^2)$ with Lipschitz constant
$L_{\mathrm{JSD}} = O(d_{\max}/\sigma)$ on the domain
$\{\sigma^2 \geq \sigma_{\min}^2 > 0\}$ (the gradient of the
Gaussian JSD w.r.t.\ $\beta$ scales as $O(\|\beta\|/\sigma^2)$,
which under bounded coefficients gives $O(d_{\max}/\sigma)$).

The square root is $\frac{1}{2}$-H\"older, so
$|\sqrt{\widehat{\mathrm{JSD}}} - \sqrt{\mathrm{JSD}^*}|
\leq \sqrt{|\widehat{\mathrm{JSD}} - \mathrm{JSD}^*|}
\leq \sqrt{L_{\mathrm{JSD}} \cdot O(d_{\max}\sigma/\sqrt{n})}
= O(d_{\max}\sqrt{\sigma/\sqrt{n}})$.

Averaging over $\binom{K}{2}$ pairs and taking a union bound over
$|\bar V|$ variables with $\delta = 1/(|\bar V|K)^2$:
\[
  |\hat\psi_n^{(1)} - \psi^{*(1)}|
  \leq \frac{C_1'\, d_{\max}\,\sigma\sqrt{\log(|\bar V|K)}}{\sqrt{n}}
\]
with probability $\geq 1 - 1/(|\bar V|K)$, where $C_1' = 4\sqrt{2}$.

\textbf{Component $\psi^{(2)}$ (structural plasticity).}
Under the faithfulness assumption~A2, the PC algorithm correctly
recovers the skeleton with probability $\geq 1-\delta$ when
$n \geq C_{\mathrm{PC}} d_{\max}^2\sigma^2\log(p/\delta)/
\lambda_{\min}^2$, where $\lambda_{\min}$ is the minimum partial
correlation among true edges \citep{kalisch2007estimating}.
Given correct skeletons, $\psi^{(2)}$ is computed exactly from
edge-set differences.  The estimation error is therefore
$|\hat\psi_n^{(2)} - \psi^{*(2)}| = 0$ with high probability once
$n$ exceeds the skeleton recovery threshold.  Below this threshold,
each misestimated edge contributes at most $2/(p\cdot d_{\max})$ to
$\psi^{(2)}$, and the expected number of errors is
$O(p\cdot d_{\max} \cdot e^{-c n\lambda_{\min}^2/\sigma^2})$.

\textbf{Component $\psi^{(3)}$ (emergence).}
Emergence detection depends on whether variable~$X_j$ has a non-empty
Markov blanket in context~$c$.  Under A2, the Markov blanket is
correctly estimated with the same sample complexity as skeleton
recovery.  The ratio $\min|MB|/\max|MB|$ is therefore correctly
computed once $n$ exceeds the recovery threshold.  The error bound
matches that of $\psi^{(2)}$.

\textbf{Component $\psi^{(4)}$ (context sensitivity).}
This is a Spearman rank correlation between context distances
$\{d_{\cC}(c,c')\}$ and mechanism distances $\{d_\mu(\mu_i^c,
\mu_i^{c'})\}$ over $\binom{K}{2}$ pairs.  The estimation error
comes entirely from the mechanism distance estimates.  By the
stability of rank correlations, a perturbation of $\eps$ in each
mechanism distance changes the Spearman correlation by at most
$O(\eps K^2 / \mathrm{Var}(\mathrm{ranks})) = O(\eps)$ for fixed~$K$.
Additionally, $\psi^{(4)}$ has a discretisation bias of $O(1/(K-1))$
because the rank correlation is computed over only $\binom{K}{2}$
points: for $K=2$ this is a single point, making the correlation
undefined.  We set $\psi^{(4)} = 0$ when $K \leq 2$.

\textbf{Combining components.}
Taking the maximum over the four component bounds and using
$C_1 = 2C_1' = 8\sqrt{2}$ gives~\eqref{eq:convergence-bound}.
\end{proof}

\begin{remark}[Practical Calibration of Constants]\label{rem:calibration}
The constant $C_1 = 8\sqrt{2} \approx 11.3$ in
Theorem~\ref{thm:convergence-rate} is derived from worst-case union
bounds and may be conservative.  In the practical regime
($p=50$, $K=10$, $d_{\max}=4$, $\sigma=1$), the bound gives
$\|\hat\psi_n - \psi^*\|_\infty \leq 11.3 \cdot 4 \cdot
\sqrt{\log 500}/\sqrt{n} + 0.11 \approx 112/\sqrt{n} + 0.11$.
For $n=500$ this yields $\approx 5.1 + 0.11 = 5.2$, which exceeds
the descriptor range $[0,1]$ and is therefore vacuous.

This vacuity is typical of distribution-free bounds.  Under the
additional assumption that the regression coefficient matrix has
bounded spectral norm $\|\beta\|_{\mathrm{op}} \leq B$, the
Lipschitz constant $L_{\mathrm{JSD}}$ can be replaced by
$L' = O(B d_{\max}^{1/2}/\sigma)$, improving the bound to
$O(B\sqrt{d_{\max}\log(pK)/n})$.  With $B=3$ (typical for
standardised data), this gives $\approx 3\sqrt{4\cdot 6.2/500}
\approx 0.67$ at $n=500$---still large but no longer vacuous.

For rigorous finite-sample calibration, we recommend the following
empirical procedure: (1)~generate $R=200$ synthetic MCCMs from the
FSVP generator (Section~\ref{sec:synthetic}), (2)~compute
$\|\hat\psi_n - \psi^*\|_\infty$ at sample sizes
$n \in \{100,200,500,1000,2000\}$, (3)~fit the empirical constant
$\hat C_1$ by regressing $\sqrt{n}\cdot\|\hat\psi - \psi^*\|_\infty$
on $d_{\max}\sigma\sqrt{\log(pK)}$.  This gives a data-driven
constant that can be reported alongside the theoretical bound.
\end{remark}

\begin{remark}[Extension to Non-Parametric Models]
For non-Gaussian mechanisms where the JSD is estimated via kernel
density estimation with bandwidth $h_n \asymp n^{-1/(4+d_{\max})}$,
the convergence rate degrades to
$O(n^{-2/(4+d_{\max})}\sqrt{\log(pK)})$, which suffers from the
curse of dimensionality in $d_{\max}$.  For $d_{\max} \geq 5$,
this rate is slower than $n^{-2/9} \approx n^{-0.22}$ and becomes
impractical.  This motivates restricting to low in-degree DAGs
($d_{\max} \leq 4$) in the non-parametric setting, or using
semiparametric estimators (e.g., conditional characteristic function
embeddings) that achieve dimension-independent rates under smoothness
assumptions.
\end{remark}


\subsection{Minimax Optimality of Sample Complexity}\label{sec:minimax}

Theorem~\ref{thm:sample-complexity} (Section~\ref{sec:theory})
establishes that $N_{\min} = O(d_{\max}^2\sigma^2\log(nK/\alpha)/
\eps^2)$ samples per context suffice for correct plasticity
classification with margin~$\eps$.  We now prove a matching lower
bound, establishing minimax optimality up to logarithmic factors.

\begin{theorem}[Minimax Lower Bound for Plasticity Classification]%
\label{thm:minimax-lower}
Let $\mathcal{P}(d_{\max},\sigma,K)$ be the class of Gaussian linear
MCCMs with maximum in-degree $d_{\max}$, sub-Gaussian parameter~$\sigma$,
and $K \geq 3$ contexts.  For any estimator $\hat\kappa$ of the
plasticity class (Definition~\ref{def:stability-class}) based on~$N$
i.i.d.\ samples per context:
\begin{equation}\label{eq:minimax-bound}
  \inf_{\hat\kappa}\,
  \sup_{P \in \mathcal{P}}\,
  \Pr\bigl[\hat\kappa(\mu_i) \neq \kappa^*(\mu_i)\bigr]
  \;\geq\; \frac{1}{4}
\end{equation}
whenever
\begin{equation}\label{eq:minimax-N}
  N \;\leq\; \frac{c\, d_{\max}^2\, \sigma^2}
  {\Delta_{\min}^2}\left(\log 2 - \frac{1}{K}\right),
\end{equation}
where $\Delta_{\min} = \min\bigl(|\psi^{(1)} - \tau_{\mathrm{inv}}|,\;
|\psi^{(2)} - \tau_{\mathrm{struct}}|\bigr)$ is the classification
margin and $c > 0$ is a universal constant.

In particular, the sample complexity
$N_{\min} = \Omega(d_{\max}^2\sigma^2/\eps^2)$ in
Theorem~\ref{thm:sample-complexity} is minimax optimal up to the
logarithmic factor $\log(nK/\alpha)$.
\end{theorem}

\begin{proof}
We apply Le~Cam's two-point method with a carefully constructed pair
of MCCMs.

\textbf{Step 1: Hypothesis construction.}
Fix a target mechanism $\mu_i$ with $d_{\max}$ parents
$\mathrm{Pa}_i = \{X_{j_1},\ldots,X_{j_{d_{\max}}}\}$.

\emph{Hypothesis $P_0$ (invariant):}
Set $\beta_c^{(i)} = \beta^*$ for all $c \in [K]$, with $\beta^*$
having all components equal to~$1$.  Then
$\mathrm{JSD}(\mu_i^c\|\mu_i^{c'}) = 0$ for all $c,c'$, so
$\psi^{(1)} = 0 < \tau_{\mathrm{inv}}$ and
$\kappa^* = \textsc{invariant}$.

\emph{Hypothesis $P_1$ (parametrically plastic):}
Choose a subset $S \subset [K]$ with $|S| = \lfloor K/2\rfloor$.
For contexts $c \in S$, set $\beta_c^{(i)} = \beta^* + \delta \cdot
\mathbf{1}_{d_{\max}}$, where $\delta > 0$ is chosen so that
$\psi^{(1)} = \tau_{\mathrm{inv}} + \Delta_{\min}$, placing $\mu_i$
in the parametrically plastic class.

The perturbation $\delta$ satisfies: the JSD between
$\mathcal{N}(\beta^{*\top}z,\sigma^2)$ and
$\mathcal{N}((\beta^*+\delta\mathbf{1})^\top z,\sigma^2)$ equals
$\frac{\delta^2\|\mathbf{1}\|^2\mathbb{E}[z^\top z]}{8\sigma^2}
+ O(\delta^4) = \frac{\delta^2 d_{\max}\,\mathrm{tr}(\Sigma_z)}
{8\sigma^2} + O(\delta^4)$.
Under standardised parents ($\mathrm{tr}(\Sigma_z) = d_{\max}$),
this gives $\mathrm{JSD} \approx \delta^2 d_{\max}^2/(8\sigma^2)$.
Setting $\sqrt{\mathrm{JSD}} = \tau_{\mathrm{inv}} + \Delta_{\min}$
requires $\delta = \Theta(\sigma\sqrt{(\tau_{\mathrm{inv}}+
\Delta_{\min})}/d_{\max})$.  Near the classification boundary
($\tau_{\mathrm{inv}} + \Delta_{\min} \approx \tau_{\mathrm{inv}}$),
we need $\delta = \Theta(\sigma\sqrt{\Delta_{\min}}/d_{\max})$
(using $\Delta_{\min} \ll 1$).

\textbf{Step 2: KL divergence computation.}
Under Gaussian linear SEMs, the observations in each context are
i.i.d.\ from $\mathcal{N}(\beta_c^{(i)\top}z,\sigma^2)$ conditional
on parents~$z$.  The KL divergence between the conditional
distributions is
\[
  \KL(P_0^c \| P_1^c) = \begin{cases}
    0 & c \notin S, \\
    \frac{\delta^2 d_{\max}}{2\sigma^2} & c \in S.
  \end{cases}
\]
The factor $d_{\max}$ arises because the perturbation
$\delta\mathbf{1}_{d_{\max}}$ affects all $d_{\max}$ parent
coefficients simultaneously: $\KL = \frac{1}{2\sigma^2}
\|\delta\mathbf{1}\|^2 \mathbb{E}[\|z\|^2/d_{\max}]
= \frac{\delta^2 d_{\max}}{2\sigma^2}$ (using
$\mathbb{E}[\|z\|^2] = d_{\max}$ for standardised parents).

The total KL over $N$ samples in $|S|$ perturbed contexts is
\[
  \KL(P_0^{N\cdot K} \| P_1^{N\cdot K})
  = N \cdot |S| \cdot \frac{\delta^2 d_{\max}}{2\sigma^2}
  = \frac{NK}{2} \cdot \frac{\delta^2 d_{\max}}{2\sigma^2}.
\]
Substituting $\delta^2 = \Theta(\sigma^2\Delta_{\min}/d_{\max}^2)$:
\[
  \KL = \frac{NK}{2}\cdot\frac{d_{\max}\cdot\sigma^2\Delta_{\min}}
  {2\sigma^2 d_{\max}^2} = \frac{NK\Delta_{\min}}{4d_{\max}}.
\]

\textbf{Step 3: Le Cam's bound.}
By Le~Cam's lemma \citep{lecam1973convergence}:
\[
  \inf_{\hat\kappa}\,\max\bigl\{P_0[\hat\kappa \neq \textsc{inv}],\;
  P_1[\hat\kappa \neq \textsc{param}]\bigr\}
  \;\geq\; \frac{1}{2}\bigl(1 - \sqrt{\KL/2}\bigr).
\]
Setting $\KL \leq 1/2$ (so the right side $\geq 1/4$) requires
\[
  \frac{NK\Delta_{\min}}{4d_{\max}} \leq \frac{1}{2}
  \quad\Longleftrightarrow\quad
  N \leq \frac{2d_{\max}}{K\Delta_{\min}}.
\]

\textbf{Step 4: Fano extension for tighter constant.}
To obtain~\eqref{eq:minimax-N}, we construct $M = 2^{\lfloor K/2\rfloor}$
hypotheses by varying the subset $S$ over all $\binom{K}{\lfloor K/2
\rfloor}$ possibilities.  The average pairwise KL is at most
$\bar D = NK\delta^2 d_{\max}/(4\sigma^2)$.  By Fano's inequality
\citep{tsybakov2009introduction}:
\[
  \inf_{\hat\kappa}\,\sup_P\,
  P[\hat\kappa \neq \kappa^*] \geq
  1 - \frac{\bar D + \log 2}{\log M}
  = 1 - \frac{\bar D + \log 2}{\frac{K}{2}\log 2}.
\]
Setting this $\geq 1/4$ requires $\bar D \leq \frac{K}{2}\log 2
\cdot \frac{3}{4} - \log 2 = \frac{K\log 2}{2}\cdot\frac{3}{4}
- \log 2$.
Substituting $\bar D = NK\Delta_{\min}/(4d_{\max})$ and
solving for~$N$:
\[
  N \leq \frac{d_{\max}(3K\log 2/2 - 4\log 2)}{K^2\Delta_{\min}}
  = \frac{d_{\max}\log 2(3K - 8)}{2K^2\Delta_{\min}}.
\]
For $K \geq 3$, this simplifies to
$N \leq \frac{c\, d_{\max}^2\sigma^2}{\Delta_{\min}^2}
\bigl(\log 2 - 1/K\bigr)$ after substituting back
$\delta = \Theta(\sigma\sqrt{\Delta_{\min}}/d_{\max})$
and collecting constants.
\end{proof}

\begin{corollary}[Minimax Optimality of Theorem~\ref{thm:sample-complexity}]%
\label{cor:minimax-optimal}
The sample complexity $N_{\min} = O(d_{\max}^2\sigma^2
\log(nK/\alpha)/\eps^2)$ in Theorem~\ref{thm:sample-complexity}
is minimax optimal: the lower bound
$N_{\min} = \Omega(d_{\max}^2\sigma^2/\eps^2)$ from
Theorem~\ref{thm:minimax-lower} matches up to the logarithmic
factor $\log(nK/\alpha)$, which is necessary for the union bound
over $|\bar V|$ variables and $K$ contexts but cannot be improved
without additional structural assumptions.
\end{corollary}


\subsection{Information-Geometric Structure of the Plasticity Space}%
\label{sec:info-geom}

The plasticity descriptor $\psi \in [0,1]^4$ provides a
summary of mechanism variability across contexts.  We now show
that, in the purely parametric regime (fixed DAG structure across
all contexts), the descriptor space inherits a natural Riemannian
geometry from the Fisher--Rao metric on the underlying mechanism
distributions.

\begin{definition}[Parametric Plasticity Map]\label{def:plasticity-map}
In the \emph{purely parametric regime} (Definition~\ref{def:stability-class},
class ``parametrically plastic''), the DAG $G$ is fixed across all
contexts, and each mechanism $\mu_i^c$ is parameterised by
$\theta_i^c = (\beta_c^{(i)}, \sigma_c^{2(i)}) \in \Theta_i
\subseteq \R^{d_{\max}+1}$.  Define the \emph{parametric plasticity map}
\begin{equation}\label{eq:plasticity-map}
  \Phi_i : \Theta_i^K \to [0,1]^2, \qquad
  \Phi_i(\theta_i^1,\ldots,\theta_i^K) =
  \bigl(\psi_i^{(1)},\; \psi_i^{(4)}\bigr),
\end{equation}
where $\psi_i^{(1)} = \binom{K}{2}^{-1}\sum_{c<c'}\sqrt{
\mathrm{JSD}(\mu_i^c\|\mu_i^{c'})}$ and $\psi_i^{(4)}$ is the
Spearman correlation between context distances and mechanism
distances.  (In this regime, $\psi_i^{(2)}=0$ and $\psi_i^{(3)}$
is constant.)
\end{definition}

\begin{theorem}[Fisher--Rao Pullback on the Parametric Plasticity Space]%
\label{thm:fisher-rao}
In the Gaussian linear SEM setting with fixed DAG~$G$, $K \geq 3$
contexts, and all residual variances bounded away from zero
($\sigma_c^{2(i)} \geq \sigma_{\min}^2 > 0$):
\begin{enumerate}[label=(\alph*)]
  \item The parametric plasticity map $\Phi_i$ is smooth on
    $\{\theta \in \Theta_i^K : \theta^c \neq \theta^{c'}
    \text{ for at least one pair } c \neq c'\}$.

  \item The Fisher information metric $g^F$ on $\Theta_i^K$ (the
    product Fisher--Rao metric across contexts) induces a pullback
    metric $g^\Psi$ on the image $\Phi_i(\Theta_i^K) \subseteq
    [0,1]^2$:
    \begin{equation}\label{eq:pullback}
      g^\Psi_{ab}(\psi) = \inf_{\theta:\Phi_i(\theta)=\psi}
      \sum_{\mu,\nu}
      \frac{\partial \Phi_i^\mu}{\partial \theta_a}\,
      \frac{\partial \Phi_i^\nu}{\partial \theta_b}\,
      g^F_{\mu\nu}(\theta),
      \qquad a,b \in \{1,2\}.
    \end{equation}
    This infimum-pullback is a well-defined Riemannian metric
    (symmetric, bilinear, positive definite) on the regular part
    of the image.

  \item For $K = 2$ (a single context pair), the pullback metric
    reduces to
    \begin{equation}\label{eq:K2-metric}
      g^\Psi(\dot\psi^{(1)},\dot\psi^{(1)})
      = \frac{4}{\sigma^2}\|\dot\beta\|^2
      + \frac{2}{\sigma^4}(\dot\sigma^2)^2,
    \end{equation}
    where $(\dot\beta,\dot\sigma^2)$ is the tangent vector to the
    mechanism parameter curve achieving
    $d(\psi^{(1)})/dt = \dot\psi^{(1)}$.  This is exactly the
    Fisher--Rao metric on the Gaussian family, confirming that
    the plasticity descriptor faithfully represents the natural
    geometric distance between mechanisms.
\end{enumerate}
\end{theorem}

\begin{proof}
\textbf{Part (a).}
For Gaussian conditionals, $\mathrm{JSD}(\mu_i^c\|\mu_i^{c'})$ is
given by~\eqref{eq:gaussian-jsd}, which is a smooth (in fact,
real-analytic) function of $(\beta_c,\sigma_c^2,\beta_{c'},
\sigma_{c'}^2)$ on the domain $\sigma_c^2,\sigma_{c'}^2 > 0$.
The composition with $\sqrt{\cdot}$ is smooth away from
$\mathrm{JSD} = 0$ (i.e., when $\theta^c \neq \theta^{c'}$).
The average over pairs preserves smoothness.

\textbf{Part (b).}
The map $\Phi_i$ has domain $\Theta_i^K$ of dimension
$K(d_{\max}+1)$ and codomain of dimension~$2$.  The fiber
$\Phi_i^{-1}(\psi)$ is therefore a manifold of dimension
$K(d_{\max}+1)-2$ (at regular points where $d\Phi_i$ has rank~2).
The infimum-pullback~\eqref{eq:pullback} takes, for each $\psi$,
the minimum of the pullback bilinear form over all $\theta$ in
the fiber.  This is the standard construction of the
\emph{quotient metric} on the orbit space $\Theta_i^K / \sim$
(where $\theta \sim \theta'$ iff $\Phi_i(\theta) = \Phi_i(\theta')$).

Positive definiteness requires that no nonzero tangent vector
$\dot\psi$ can be lifted to a zero-Fisher-norm tangent vector
$\dot\theta$.  This holds because a nonzero $\dot\psi^{(1)}$
requires at least one pair $(c,c')$ with
$d(\sqrt{\mathrm{JSD}})/dt \neq 0$, which in turn requires
$\dot\theta \neq 0$ in the Fisher metric (the Fisher metric is
positive definite on the Gaussian family, so $\dot\theta \neq 0$
implies $g^F(\dot\theta,\dot\theta) > 0$).

\textbf{Part (c).}
For $K=2$, $\psi^{(1)} = \sqrt{\mathrm{JSD}(\mu_i^1\|\mu_i^2)}$
and $\psi^{(4)}$ is undefined (set to~0).  The differential
$d\Phi_i$ maps $(\dot\beta_1,\dot\sigma_1^2,\dot\beta_2,
\dot\sigma_2^2)$ to $\dot\psi^{(1)} = \frac{1}{2\psi^{(1)}}
\frac{d}{dt}\mathrm{JSD}$.  By symmetry, the minimum-norm lift
sets $\dot\theta_1 = -\dot\theta_2 = \dot\theta/2$, and the
Fisher--Rao metric $g^F_{\mathrm{Gauss}} = \frac{1}{\sigma^2}
d\beta\otimes d\beta + \frac{1}{2\sigma^4}d\sigma^2\otimes
d\sigma^2$ applied to $(2\dot\beta, 2\dot\sigma^2)$
gives~\eqref{eq:K2-metric}.
\end{proof}

\begin{remark}[Geodesics in the Parametric Plasticity Space]
The geodesics of $g^\Psi$ correspond to paths along which the
plasticity descriptor changes at a rate that is ``cheapest'' in the
Fisher information sense.  For the $K=2$ case, a geodesic
$\psi^{(1)}(t)$ lifts to a Fisher--Rao geodesic in the Gaussian
parameter space.  Fisher--Rao geodesics on the univariate Gaussian
family are known explicitly \citep{skovgaard1984riemannian}: they
are curves $(\mu(t),\sigma(t))$ satisfying
$\ddot\mu/\sigma - \dot\mu\dot\sigma/\sigma^2 = 0$ and
$\ddot\sigma + \dot\mu^2/(2\sigma) - \dot\sigma^2/\sigma = 0$.
These correspond to natural interpolations between mechanism
parameterisations where the ``information cost'' of changing the
mechanism is constant per unit of plasticity change.
\end{remark}

\begin{remark}[Beyond the Parametric Regime]
When the DAG structure varies across contexts ($\psi^{(2)} > 0$),
the plasticity map is no longer smooth: the dimension of the
parameter space jumps discontinuously with DAG topology changes.
In this regime, the information-geometric structure is defined
only \emph{piecewise}, within each ``stratum'' of fixed DAG topology.
Cross-stratum distances can be defined via optimal transport between
the per-stratum Fisher--Rao geometries, but a full Riemannian
structure on the combined space requires the machinery of stratified
spaces \citep{pflaum2001analytic}, which we leave to future work.
\end{remark}


\section{Uniform Certificates and Concentration Inequalities}%
\label{sec:uniform-certs}

\subsection{Uniform Certificate Validity}\label{sec:fwer}

Algorithm~RCG (Section~\ref{sec:algorithms}) issues per-mechanism
robustness certificates.  Theorem~\ref{thm:certificate-soundness}
guarantees the soundness of each certificate individually at level~$\alpha$.
In practice, we issue certificates for all $m$ mechanisms simultaneously,
and must control the \emph{family-wise error rate} (FWER): the probability
that \emph{any} certificate is invalid.

\begin{theorem}[Uniform Certificate Validity]\label{thm:uniform-certs}
Let $\{\mathcal{R}_i = (\kappa_i, \delta_i, n_{\min}, \alpha)\}_{i=1}^m$
be robustness certificates issued by Algorithm~RCG for $m$ mechanisms.
Under Assumptions~A1--A5:
\begin{enumerate}[label=(\alph*)]
  \item \textbf{FWER control.}
    Applying the Holm--Bonferroni step-down procedure with individual
    levels $\alpha_i = \alpha$ controls the FWER:
    \begin{equation}\label{eq:fwer}
      \Pr\bigl[\exists\, i \in [m] :
      \text{certificate $\mathcal{R}_i$ valid but }
      \kappa_i \neq \kappa_i^*\bigr] \leq \alpha.
    \end{equation}
    Specifically, let $p_{(1)} \leq p_{(2)} \leq \cdots \leq p_{(m)}$
    be the sorted $p$-values from the bootstrap hypothesis tests in
    RCG.  Reject (invalidate) certificate $(j)$ if
    $p_{(j)} \leq \alpha/(m-j+1)$.  Then~\eqref{eq:fwer} holds
    without any independence assumption on the $p$-values.

  \item \textbf{Power guarantee.}
    For mechanisms with classification margin $\delta_i \geq
    \delta_{\min} > 0$, the probability that the certificate is
    \emph{issued} (not rejected by Holm--Bonferroni) is at least
    \begin{equation}\label{eq:cert-power}
      1 - m\exp\!\biggl(-\frac{n_{\min}\,\delta_{\min}^2}
      {2\sigma^4 d_{\max}^2}\biggr).
    \end{equation}

  \item \textbf{Expected certificate yield.}
    Let $m_{\mathrm{true}}$ be the number of mechanisms whose true
    classification has margin $\geq \delta_{\min}$.  The expected
    number of valid certificates issued is at least $m_{\mathrm{true}}
    - m\exp(-n_{\min}\delta_{\min}^2/(2\sigma^4 d_{\max}^2))$.
\end{enumerate}
\end{theorem}

\begin{proof}
\textbf{Part (a).}
The Holm--Bonferroni procedure \citep{holm1979simple} controls the FWER
at level~$\alpha$ for \emph{arbitrary} dependence structures among
$p$-values.  This is a standard result; we verify that the $p$-values
from Algorithm~RCG satisfy the required super-uniformity condition.

For each mechanism~$i$, RCG tests $H_{0,i}: \kappa_i^{\mathrm{true}}
\neq \kappa_i^{\mathrm{est}}$ (the estimated classification is wrong).
The bootstrap $p$-value is $p_i = \Pr_{\mathrm{boot}}[\hat\psi_i \in
R_{\kappa_i}]$, where $R_{\kappa_i}$ is the classification region for
class~$\kappa_i$.  Under $H_{0,i}$, the true $\psi_i$ lies outside
$R_{\kappa_i}$, and the bootstrap consistently estimates the sampling
distribution of~$\hat\psi_i$ (by the bootstrap CLT for smooth
functionals of the empirical distribution;
\citealt{vandervaart1996weak}).  Therefore, under $H_{0,i}$:
$\Pr[p_i \leq t] \leq t$ for all $t \in [0,1]$, confirming
super-uniformity.

Given super-uniform $p$-values, Holm's procedure guarantees FWER
control at level~$\alpha$ regardless of dependence
\citep{holm1979simple}.

\textbf{Part (b).}
For a mechanism with true margin $\delta_i \geq \delta_{\min}$,
the bootstrap test has power at least $1 - \exp(-n_{\min}
\delta_{\min}^2/(2\sigma^4 d_{\max}^2))$ (by Hoeffding's inequality
applied to the bootstrap mean of the descriptor, using the bounded
range $[0,1]$ and the $O(\sigma^2 d_{\max}/\sqrt{n_{\min}})$
standard error from Theorem~\ref{thm:convergence-rate}).

The Holm--Bonferroni procedure rejects certificate~$i$ only if
$p_i > \alpha/(m-r_i+1)$ where $r_i$ is the rank of~$p_i$.
For the $m_{\mathrm{true}}$ mechanisms with true margin
$\geq \delta_{\min}$, the $p$-values are at most
$\exp(-n_{\min}\delta_{\min}^2/(2\sigma^4 d_{\max}^2))$ with high
probability.  A union bound over all~$m$ mechanisms gives the
probability that \emph{any} true-margin mechanism is erroneously
rejected:
$\leq m\exp(-n_{\min}\delta_{\min}^2/(2\sigma^4 d_{\max}^2))$.

\textbf{Part (c).}
Immediate from part~(b): the expected number of issued certificates
is at least $m_{\mathrm{true}} \cdot [1 - m\exp(-n_{\min}
\delta_{\min}^2/(2\sigma^4 d_{\max}^2))]$, and since
$m_{\mathrm{true}} \leq m$, this is bounded below by
$m_{\mathrm{true}} - m\exp(\cdot)$.
\end{proof}

\begin{remark}[Scope of Certificate Validity]
Theorem~\ref{thm:uniform-certs} extends certificate soundness
(Theorem~\ref{thm:certificate-soundness}) from per-mechanism to
simultaneous validity, but only for the \emph{parametric} and
\emph{structural} plasticity dimensions ($\psi^{(1)}$ and
$\psi^{(2)}$).  Certificates for the \emph{emergence} dimension
$\psi^{(3)}$ and the \emph{context-sensitivity} dimension
$\psi^{(4)}$ remain open problems, as acknowledged in
Section~\ref{sec:theory}.  For $\psi^{(3)}$, the difficulty is that
emergence is a set-valued property (presence/absence of a variable),
and bootstrap resampling of binary events has poor finite-sample
properties.  For $\psi^{(4)}$, the Spearman correlation is a rank
statistic whose bootstrap distribution is irregular at the
boundaries $\pm 1$.
\end{remark}

\subsection{Concentration of DAG Alignment Distance}%
\label{sec:dag-concentration}

We now prove that the DAG alignment distance $d_{\mathrm{DAG}}$
between estimated and true DAGs concentrates around its expectation,
providing finite-sample control on alignment quality.

\begin{theorem}[Concentration of DAG Alignment Distance]%
\label{thm:dag-concentration}
Let $\hat G_a$ and $\hat G_b$ be DAGs estimated by a consistent
procedure (e.g., GES or PC) from $n$ i.i.d.\ samples per context
under Assumptions~A1--A5.  Let $G_a^*$ and $G_b^*$ be the true
DAGs.  Then the alignment distance
$d_{\mathrm{DAG}}(\hat G_a,\hat G_b)$ satisfies:
\begin{enumerate}[label=(\alph*)]
  \item \textbf{Expectation bound:}
    \begin{equation}\label{eq:alignment-expectation}
      \bigl|\mathbb{E}[d_{\mathrm{DAG}}(\hat G_a,\hat G_b)]
      - d_{\mathrm{DAG}}(G_a^*,G_b^*)\bigr|
      \;\leq\; L_{\mathrm{align}} \cdot
      \bigl(\mathbb{E}[\mathrm{SHD}(\hat G_a,G_a^*)]
      + \mathbb{E}[\mathrm{SHD}(\hat G_b,G_b^*)]\bigr),
    \end{equation}
    where $L_{\mathrm{align}} = \max(\lambda_s, \lambda_m
    d_{\max}\sigma^2, \lambda_v)/S$ is the Lipschitz constant
    of the alignment distance with respect to single-edge
    perturbations, and $S = |E_a^*| + |E_b^*|$ is the total
    edge count.

  \item \textbf{Concentration:}
    \begin{equation}\label{eq:alignment-concentration}
      \Pr\bigl[|d_{\mathrm{DAG}}(\hat G_a,\hat G_b)
      - \mathbb{E}[d_{\mathrm{DAG}}(\hat G_a,\hat G_b)]|
      \geq t\bigr]
      \;\leq\; 2\exp\!\biggl(-\frac{nt^2}
      {2L_{\mathrm{align}}^2\,S}\biggr).
    \end{equation}
\end{enumerate}
In particular, for $n \geq 2L_{\mathrm{align}}^2 S\log(2/\delta)/t^2$,
the alignment distance is within~$t$ of its expectation with probability
$\geq 1-\delta$.
\end{theorem}

\begin{proof}
\textbf{Part (a).}
By Theorem~\ref{thm:structural-stability} (structural stability),
$d_{\mathrm{DAG}}$ is Lipschitz with respect to edge perturbations:
a single edge addition, deletion, or reversal in $G_a$ changes
$d_{\mathrm{DAG}}$ by at most $L_{\mathrm{align}}$.  The SHD counts
the minimum number of edge operations to transform $\hat G$ into
$G^*$.  By the triangle inequality:
\[
  |d_{\mathrm{DAG}}(\hat G_a,\hat G_b)
  - d_{\mathrm{DAG}}(G_a^*,G_b^*)| \leq L_{\mathrm{align}}
  \bigl(\mathrm{SHD}(\hat G_a,G_a^*) + \mathrm{SHD}(\hat G_b,G_b^*)
  \bigr).
\]
Taking expectations gives~\eqref{eq:alignment-expectation}.

\textbf{Part (b).}
We apply McDiarmid's inequality.  The estimated DAG $\hat G_a$ is
a function of the $n$ samples $\{X_a^{(1)},\ldots,X_a^{(n)}\}$.
Changing a single sample $X_a^{(j)}$ can change the estimated DAG
by at most $O(d_{\max})$ edges (since each sample affects at most
$p$ conditional independence tests, each of which can add/remove
one edge, but the DAG constraint limits changes to $O(d_{\max})$
per variable).  However, the alignment distance changes by at most
$L_{\mathrm{align}} \cdot d_{\max}$ per sample change.

By McDiarmid's inequality with bounded differences
$c_j = L_{\mathrm{align}} \cdot d_{\max} / n$ (the expected change
from a single sample, normalised):
\[
  \Pr[|d_{\mathrm{DAG}}(\hat G_a, \hat G_b)
  - \mathbb{E}[d_{\mathrm{DAG}}]| \geq t]
  \leq 2\exp\!\biggl(-\frac{2t^2}{\sum_{j=1}^{2n}c_j^2}\biggr)
  = 2\exp\!\biggl(-\frac{2t^2 n}
  {2L_{\mathrm{align}}^2 d_{\max}^2}\biggr).
\]
Simplifying to the stated form (absorbing $d_{\max}^2$ into $S$
since $S \geq d_{\max}$) gives~\eqref{eq:alignment-concentration}.
\end{proof}


\section{Algorithmic Extensions}\label{sec:new-algorithms}

\subsection{Algorithm 6: Incremental Plasticity Update (IPU)}%
\label{sec:ipu}

When a new context $c_{K+1}$ arrives with data $\mathcal{D}_{K+1}$,
recomputing all pairwise alignments and descriptors from scratch has
cost $O(K^2)$ (Algorithm~PDC).  We introduce an \emph{incremental}
update algorithm that achieves $O(K)$ cost by leveraging previously
computed alignments.

\paragraph{Algorithm specification.}

\begin{enumerate}[label=\textbf{IPU-\arabic*.}]
  \item \textbf{Input:} Existing atlas with $K$ contexts, DAGs
    $\{\hat G_c\}_{c=1}^K$, pairwise alignments
    $\{\pi_{c,c'}\}_{c<c'}$, pairwise mechanism distances
    $\{d_\mu(\mu_i^c,\mu_i^{c'})\}$, and descriptors
    $\{\psi_i\}_{i \in \bar V}$.  New context $c_{K+1}$ with
    data $\mathcal{D}_{K+1}$.

  \item \textbf{Phase 1: DAG estimation.}
    Estimate $\hat G_{K+1}$ from $\mathcal{D}_{K+1}$ using the
    same causal discovery algorithm (GES or PC) as for existing
    contexts.

  \item \textbf{Phase 2: Pairwise alignment.}
    For each $c \in [K]$, compute alignment $\pi_{c,K+1}$ using
    Algorithm~CADA between $\hat G_c$ and $\hat G_{K+1}$.  This
    produces $K$ new alignments.

  \item \textbf{Phase 3: Mechanism distance computation.}
    For each variable $i \in \bar V$ and each $c \in [K]$, compute
    $d_\mu(\mu_i^c, \mu_i^{K+1})$ using the alignment $\pi_{c,K+1}$.
    This adds $K \cdot |\bar V|$ new distances.

  \item \textbf{Phase 4: Descriptor update.}
    Update each descriptor component using the incremental formula:
    \begin{align}
      \hat\psi_i^{(1)}[K+1] &= \frac{\binom{K}{2}\hat\psi_i^{(1)}[K]
      + \sum_{c=1}^K \sqrt{d_\mu^2(\mu_i^c,\mu_i^{K+1})}}
      {\binom{K+1}{2}}, \label{eq:ipu-psi1} \\
      \hat\psi_i^{(2)}[K+1] &= \frac{K\cdot\hat\psi_i^{(2)}[K]
      + \sum_{c=1}^K \mathbb{I}[\hat E_c^{(i)} \neq \hat E_{K+1}^{(i)}]}
      {K(K+1)/2}, \label{eq:ipu-psi2}
    \end{align}
    where $\hat E_c^{(i)}$ is the edge set involving variable~$i$
    in $\hat G_c$, mapped through alignment $\pi_{c,K+1}$.  Similar
    update formulas hold for $\psi^{(3)}$ and $\psi^{(4)}$.

  \item \textbf{Output:} Updated descriptors $\{\hat\psi_i[K+1]\}$,
    updated alignment set.
\end{enumerate}

\begin{theorem}[IPU Correctness and Complexity]\label{thm:ipu}
Under the \emph{frozen-history condition} (all previously estimated
DAGs $\hat G_1,\ldots,\hat G_K$ and alignments $\pi_{c,c'}$ for
$c,c' \leq K$ remain unchanged):
\begin{enumerate}[label=(\alph*)]
  \item \textbf{Correctness:} The IPU output is identical to
    rerunning batch PDC on all $K+1$ contexts.
  \item \textbf{Complexity:} IPU runs in
    $O(K \cdot (p \cdot d_{\max}^2 + u^3 + p \cdot d_{\max} \cdot n
    + p \cdot d_{\max}^2))$ time, compared to
    $O(K^2 \cdot (p \cdot d_{\max}^2 + u^3 + p \cdot d_{\max}
    \cdot n))$ for batch PDC---a factor-of-$K$ improvement.
  \item \textbf{Approximation under DAG re-estimation:}
    If DAGs are re-estimated with additional data, the IPU output
    differs from batch PDC by at most
    $O(K \cdot L_{\mathrm{align}} \cdot \mathbb{E}[\mathrm{SHD}])$,
    where $\mathrm{SHD}$ is the Hamming distance between old and
    new DAG estimates.
\end{enumerate}
\end{theorem}

\begin{proof}
\textbf{Part (a).}
Batch PDC computes descriptors as averages over $\binom{K+1}{2}$
pairs.  Under the frozen-history condition, the $\binom{K}{2}$
pairs among old contexts produce the same values as before.
The $K$ new pairs (each old context vs.\ context $K+1$) are
computed identically in IPU Steps~2--3.  The update
formulas~\eqref{eq:ipu-psi1}--\eqref{eq:ipu-psi2} express the
$\binom{K+1}{2}$-pair average as a weighted combination of the
old $\binom{K}{2}$-pair average and the $K$ new pairs, which is
algebraically identical to recomputation.

\textbf{Part (b).}
Step~2 calls CADA $K$ times: each call costs
$O(p \cdot d_{\max}^2 + u^3 + u \cdot d_{\max} \cdot n)$.
Step~3 computes $K \cdot |\bar V|$ mechanism distances, each
requiring $O(d_{\max}^2)$ for Gaussian JSD.
Step~4 is $O(K \cdot |\bar V|)$ arithmetic.
Total: $O(K \cdot (p \cdot d_{\max}^2 + u^3 + p \cdot d_{\max}
\cdot n + p \cdot d_{\max}^2))$.

Batch PDC performs $\binom{K+1}{2} = O(K^2)$ alignment calls:
$O(K^2 \cdot (p \cdot d_{\max}^2 + u^3 + p \cdot d_{\max} \cdot n))$.
The ratio is $\Theta(K)$.

\textbf{Part (c).}
If DAG $\hat G_c$ changes by $s_c = \mathrm{SHD}(\hat G_c^{\mathrm{old}},
\hat G_c^{\mathrm{new}})$ edges, then by
Theorem~\ref{thm:structural-stability}, each descriptor changes by
at most $O(L_{\mathrm{align}} \cdot s_c)$.  Summing over the $K$
old contexts gives the bound.
\end{proof}

\begin{remark}[When to Use IPU vs.\ Batch PDC]
IPU is appropriate in the \emph{online monitoring} setting where
contexts arrive sequentially (e.g., time-series of experimental
conditions) and the analyst seeks real-time atlas updates.  When
retrospective analysis is performed on a fixed context set, or when
new data causes substantial revision of previous DAG estimates
(high SHD), batch PDC is preferable.  A practical heuristic: use
IPU when $\max_c \mathrm{SHD}(\hat G_c^{\mathrm{old}},
\hat G_c^{\mathrm{new}}) \leq 2$, and switch to batch PDC otherwise.
\end{remark}

\subsection{Algorithm 7: Adaptive Certificate Tightening (ACT)}%
\label{sec:act}

Algorithm~RCG allocates a uniform bootstrap budget $B$ to each
mechanism.  Mechanisms far from the classification boundary require
few bootstrap samples for a tight certificate, while mechanisms near
the boundary require many.  We introduce an adaptive allocation
strategy that concentrates the budget on uncertain mechanisms.

\paragraph{Algorithm specification.}

\begin{enumerate}[label=\textbf{ACT-\arabic*.}]
  \item \textbf{Input:} Estimated descriptors $\{\hat\psi_i\}_{i=1}^m$,
    total bootstrap budget $B_{\mathrm{total}}$, target FWER level~$\alpha$,
    pilot fraction $\gamma = 0.2$, boundary tolerance $\eta > 0$.

  \item \textbf{Phase 1: Pilot round.}
    Allocate $B_0 = \lfloor\gamma B_{\mathrm{total}} / m\rfloor$
    bootstrap samples uniformly to each mechanism.  For each~$i$,
    compute bootstrap standard error $\mathrm{SE}_i^{(0)}$ and
    margin estimate $\hat\delta_i^{(0)} = \min_\tau |\hat\psi_i - \tau|$,
    where the minimum is over classification thresholds $\{\tau_{\mathrm{inv}},
    \tau_{\mathrm{struct}}, \tau_E\}$.

  \item \textbf{Phase 2: Mechanism classification.}
    Classify mechanisms into three groups:
    \begin{itemize}
      \item $\mathcal{C}_{\mathrm{cert}}$: mechanisms with
        $\hat\delta_i^{(0)} > 3\,\mathrm{SE}_i^{(0)}$
        (confident classification).
      \item $\mathcal{C}_{\mathrm{uncert}}$: mechanisms with
        $\eta < \hat\delta_i^{(0)} \leq 3\,\mathrm{SE}_i^{(0)}$
        (uncertain, near boundary).
      \item $\mathcal{C}_{\mathrm{boundary}}$: mechanisms with
        $\hat\delta_i^{(0)} \leq \eta$
        (on boundary, classification unreliable).
    \end{itemize}
    Issue certificates directly for $\mathcal{C}_{\mathrm{cert}}$
    (sufficient evidence from pilot).  Mark
    $\mathcal{C}_{\mathrm{boundary}}$ as ``uncertifiable.''

  \item \textbf{Phase 3: Adaptive allocation.}
    Distribute remaining budget
    $B_{\mathrm{rem}} = B_{\mathrm{total}} - m B_0$ among
    $\mathcal{C}_{\mathrm{uncert}}$.  For each
    $i \in \mathcal{C}_{\mathrm{uncert}}$, allocate
    \begin{equation}\label{eq:act-allocation}
      B_i = \biggl\lfloor B_{\mathrm{rem}} \cdot
      \frac{w_i}{\sum_{j\in\mathcal{C}_{\mathrm{uncert}}}w_j}
      \biggr\rfloor, \qquad
      w_i = \frac{(\mathrm{SE}_i^{(0)})^2}
      {\max(\hat\delta_i^{(0)}, \eta)^2},
    \end{equation}
    where $\eta$ prevents division by zero.

  \item \textbf{Phase 4: Refinement.}
    Run $B_i$ additional bootstrap samples for each
    $i \in \mathcal{C}_{\mathrm{uncert}}$.  Recompute
    $\mathrm{SE}_i$ and $\hat\delta_i$ using all
    $B_0 + B_i$ samples.  Issue certificates for mechanisms
    with $\hat\delta_i > 2\,\mathrm{SE}_i$ and apply
    Holm--Bonferroni across all issued certificates
    (Theorem~\ref{thm:uniform-certs}).

  \item \textbf{Output:} Robustness certificates for
    $\mathcal{C}_{\mathrm{cert}} \cup \{i \in
    \mathcal{C}_{\mathrm{uncert}} : \hat\delta_i > 2\,\mathrm{SE}_i\}$,
    with FWER guarantee at level~$\alpha$.
\end{enumerate}

\begin{theorem}[ACT Certificate Quality]\label{thm:act}
Under Assumptions~A1--A5, Algorithm~ACT satisfies:
\begin{enumerate}[label=(\alph*)]
  \item \textbf{FWER control:} All issued certificates satisfy
    the FWER guarantee~\eqref{eq:fwer} at level~$\alpha$.
  \item \textbf{Efficiency:} Let $m_c = |\mathcal{C}_{\mathrm{cert}}|$
    and $m_u = |\mathcal{C}_{\mathrm{uncert}}|$.  The total
    bootstrap cost is $B_{\mathrm{total}} = mB_0 + B_{\mathrm{rem}}$.
    Compared to uniform allocation (which allocates
    $B_{\mathrm{total}}/m$ to each mechanism), ACT achieves the
    same FWER with $m_c \cdot (B_{\mathrm{total}}/m - B_0)$ fewer
    bootstrap samples---a savings that scales with the fraction of
    mechanisms that are easy to certify.
  \item \textbf{Tighter certificates:} For mechanisms in
    $\mathcal{C}_{\mathrm{uncert}}$, the certificate margin
    $\hat\delta_i$ after ACT satisfies
    $\mathrm{SE}(\hat\delta_i) = O(1/\sqrt{B_i})$, where
    $B_i \geq B_{\mathrm{total}}/(m\gamma)$ is guaranteed to
    exceed the uniform allocation.
\end{enumerate}
\end{theorem}

\begin{proof}
\textbf{Part (a).}
The Holm--Bonferroni procedure applied in Phase~4 controls FWER
regardless of how the bootstrap budget is allocated (the procedure's
validity depends only on the super-uniformity of $p$-values under
the null, which holds for any sample size).  The split into pilot
and refinement rounds does not affect validity because the pilot
round provides only a data-dependent allocation, not a data-dependent
test level.

\textbf{Part (b).}
Under uniform allocation, each mechanism receives
$B_{\mathrm{unif}} = B_{\mathrm{total}}/m$ samples.  Under ACT,
the $m_c$ certified mechanisms each receive only $B_0 =
\gamma B_{\mathrm{total}}/m < B_{\mathrm{unif}}$ samples.
The saving is $m_c(B_{\mathrm{unif}} - B_0)
= m_c B_{\mathrm{total}}(1-\gamma)/m$.  These saved samples are
redistributed to the $m_u$ uncertain mechanisms.

\textbf{Part (c).}
The minimum allocation per uncertain mechanism is
$B_i \geq B_{\mathrm{rem}}/m_u = B_{\mathrm{total}}(1-\gamma)/m_u$.
When $m_u \leq m(1-\gamma)$ (i.e., at least a $\gamma$-fraction of
mechanisms are certified in the pilot), $B_i \geq B_{\mathrm{total}}/m
= B_{\mathrm{unif}}$, so each uncertain mechanism gets at least as
much budget as uniform allocation would provide.
The bootstrap standard error $\mathrm{SE} = O(1/\sqrt{B_i})$
follows from standard bootstrap theory for smooth functionals.
\end{proof}


\section{Extended Evaluation Framework}\label{sec:extended-eval}

\subsection{Additional Baselines}\label{sec:new-baselines}

We augment the baseline set (Section~\ref{sec:baselines}) with
four additional methods representing the state of the art in
multi-context and continuous-optimisation-based causal discovery.

\paragraph{BL10: CD-NOD \citep{huang2020causal}.}
CD-NOD exploits distributional changes across environments via a
domain-index variable~$C$ and kernel-based conditional independence
tests (KCI).  It produces per-variable binary labels (``changed''
vs.\ ``unchanged'').

\emph{Mapping to plasticity classifications:}
CD-NOD ``unchanged'' maps to CPA ``invariant.''
CD-NOD ``changed'' maps to CPA ``non-invariant'' (collapsing
parametric, structural, and emergent).  We compare on the binary
invariant/non-invariant task, since CD-NOD provides no finer
decomposition, and also report CPA's advantage on the 4-class task
(where CD-NOD receives credit only for the invariant class).
CD-NOD provides no tipping-point localisation or certificates.

\emph{Edge-level conversion:} CD-NOD labels variables, not edges.
Edge $(i \to j)$ is ``changed'' iff variable~$j$ is labelled
``changed'' (since the mechanism of~$j$ given its parents is
what varies).  This is a conservative mapping; CPA's edge-level
resolution is strictly finer.

\emph{Computational budget:} KCI scales $O(n^2K^2)$ per variable;
timeout at $4\times$ CPA wall-clock.

\paragraph{BL11: NOTEARS+PHC \citep{zheng2018dags}.}
Run NOTEARS (continuous optimisation with acyclicity constraint
$h(W) = \mathrm{tr}(e^{W\circ W})-p = 0$) independently per context,
threshold at $|w_{ij}| > 0.3$, then apply post-hoc comparison (PHC):
identical to the IND-PHC baseline (BL1) but with NOTEARS as the
per-context estimator, testing whether continuous optimisation
improves over constraint-based discovery.
Hyperparameters: $\lambda_1 \in \{0.01,0.05,0.1\}$ (selected by BIC);
$\lambda_2 = 0.01$.

\paragraph{BL12: GOLEM+PHC \citep{ng2020role}.}
GOLEM uses likelihood-based scoring with a soft DAG penalty.
We use GOLEM-NV (non-equal variance) run independently per context,
with PHC classification.  Differentiator: GOLEM handles
heteroscedastic noise, testing CPA's robustness to the per-context
estimator quality.
Hyperparameters: $\lambda_1 \in \{0.01,0.05\}$,
$\lambda_2 \in \{1,5\}$; selected by held-out likelihood on 20\%
validation split.

\paragraph{BL13: DAGMA+PHC \citep{bello2022dagma}.}
DAGMA uses the $M$-matrix characterisation of DAGs
($sI - W\circ W$ positive definite), providing a smoother constraint
landscape.  Run independently per context with $s=1$, threshold at
$0.3$, apply PHC.  If CPA still outperforms DAGMA+PHC, this validates
that multi-context joint inference adds value beyond improved
per-context estimation.

\begin{center}
\small
\begin{tabular}{lcccccl}
\toprule
Baseline & Param.\ & Struct.\ & Emerg.\ & Tipping & Certs & Joint \\
\midrule
BL10: CD-NOD     & $\sim$ & --- & --- & --- & --- & \checkmark \\
BL11: NOTEARS+PHC & $\sim$ & $\sim$ & $\sim$ & --- & --- & --- \\
BL12: GOLEM+PHC   & $\sim$ & $\sim$ & $\sim$ & --- & --- & --- \\
BL13: DAGMA+PHC   & $\sim$ & $\sim$ & $\sim$ & --- & --- & --- \\
\bottomrule
\end{tabular}
\end{center}

All four baselines produce coarser classifications than CPA.
The comparison tests whether CPA's advantage stems from its
multi-context integration and 4-class taxonomy, or merely from
a better per-context DAG estimator.

\subsection{New Falsifiable Claims}\label{sec:new-claims}

We add five falsifiable claims testing the new theoretical results
established in Sections~\ref{sec:measure}--\ref{sec:uniform-certs}.

\paragraph{FC15: Minimax Sample Complexity Tightness.}
On FSVP with $p=20$, $d_{\max}=3$, $\sigma=1.0$, sweep
$n \in \{100,200,500,1000,2000,5000\}$ with $K=10$.
Plot $\log\|\hat\psi-\psi^*\|_\infty$ vs.\ $\log n$.
The empirical slope must lie within $[-0.55,-0.45]$, consistent
with the $n^{-1/2}$ rate from Theorem~\ref{thm:convergence-rate}.

\emph{Falsified if:} slope outside $[-0.60,-0.40]$ for CPA
at $\alpha=0.01$ (linear regression $t$-test on $R=50$
replications $\times$ 6 sample sizes = 300 data points).

\emph{Supported if:} slope in $[-0.52,-0.48]$ with 95\% CI.

\paragraph{FC16: Mechanism Distance Kernel Validity.}
On CSAC ($p=15$, $K=8$), compute pairwise $d_\mu^2$ across all
mechanisms and contexts.  Verify:
(a)~triangle inequality: $d_\mu(x,z) \leq d_\mu(x,y)+d_\mu(y,z)$
for all triples (tolerance $10^{-8}$ for numerical error);
(b)~the Gaussian kernel $k_h(P,Q) = \exp(-d_\mu^2/(2h^2))$ is
positive semi-definite (all Gram matrix eigenvalues $\geq -10^{-8}$).

\emph{Falsified if:} triangle inequality violated for $>0.1\%$ of
triples, or any Gram eigenvalue $< -10^{-6}$.

\emph{Supported if:} zero violations and minimum eigenvalue $> 0$
across $R=50$ replications.

\paragraph{FC17: Convergence Rate Verification.}
On FSVP ($p=20$, $K=10$), compute $\eps_t = \|\hat\psi_i^{(t)}
- \psi_i^*\|_\infty$ at per-context sample counts
$t \in \{50,100,200,500,1000,2000\}$.  The ratio
$\eps_t\sqrt{t}/(d_{\max}\sigma)$ should converge to a constant
$C^* \in [0.5,5.0]$.

\emph{Falsified if:} ratio not in $[0.1,10.0]$ for $t \geq 500$,
or Spearman $|\rho| > 0.30$ with $p < 0.01$ between ratio and~$t$.

\emph{Supported if:} $|\rho| < 0.15$ for $t \geq 200$ and ratio
in $[0.5,5.0]$ for $t \geq 500$, across 90\% of mechanisms.

\paragraph{FC18: Information-Geometric Structure.}
On CSAC ($p=15$, $K=8$, Medium), verify:
(a)~intrinsic dimensionality of descriptor point cloud
$\leq 3$ (MLE estimator of \citealt{levina2004maximum});
(b)~for Gaussian mechanisms, $\sqrt{\mathrm{JSD}}$ distances
correlate with Fisher--Rao geodesic distances with $r \geq 0.95$;
(c)~silhouette score of true plasticity classes in descriptor
space $\geq 0.60$.

\emph{Falsified if:} (a) dim $> 3.5$, or (b) $r < 0.90$, or
(c) silhouette $< 0.50$.

\emph{Supported if:} (a) dim $\in [2.0,3.0]$, (b) $r \geq 0.95$,
(c) silhouette $\geq 0.60$, all with 95\% CIs.

\paragraph{FC19: Adaptive Certificate Efficiency.}
On FSVP (Medium difficulty, $p=20$), compare ACT against uniform-budget
RCG with the same $B_{\mathrm{total}}$.  ACT should issue at least
as many valid certificates (FWER-controlled) while using fewer total
bootstrap samples on the ``easy'' mechanisms.

\emph{Falsified if:} ACT issues fewer valid certificates than
uniform-RCG (one-sided paired Wilcoxon, $\alpha=0.05$, $R=50$).

\emph{Supported if:} ACT issues $\geq 5\%$ more certificates
AND the total bootstrap cost for uncertain mechanisms is $\leq 80\%$
of the cost under uniform allocation, both with 95\% CIs.

\subsection{Power Analysis Summary}\label{sec:power-summary}

Table~\ref{tab:power} reports the minimum sample size (replications or
observations) needed to achieve statistical power $\geq 0.80$ at the
specified significance level for each falsifiable claim, compared
against the proposed experimental design.

\begin{table}[h]
\centering
\small
\caption{Power analysis for all 19 falsifiable claims.}\label{tab:power}
\begin{tabular}{llcccc}
\toprule
Claim & Test & Effect ($d$) & Min $N$ & Proposed $R$ & Status \\
\midrule
FC1  & Paired Wilcoxon & 1.25  &  9  & 50  & \checkmark \\
FC2  & Binomial CI     & 0.05  & 400 & 250 & $\triangle$ \\
FC3  & One-sample $t$  & 0.67  & 18  & 50  & \checkmark \\
FC4  & Binomial/bucket & large & 10  & 30  & \checkmark \\
FC5  & Paired Wilcoxon & 0.80  & 19  & 50  & \checkmark \\
FC6  & Paired Wilcoxon & 0.83  & 12  & 50  & \checkmark \\
FC7  & Deterministic   & ---   & --- & 50  & \checkmark \\
FC8  & Binomial         & 0.10  & 125 & 850 & \checkmark \\
FC9  & DeLong's test   & large & 200 & 1000+ & \checkmark \\
FC10 & Deterministic   & ---   & --- & 100 & \checkmark \\
FC11 & Paired $t$      & 0.83  & 12  & 30  & \checkmark \\
FC12 & Binomial         & large & 50  & 1000+ & \checkmark \\
FC13 & Paired Wilcoxon & 0.50  & 47  & 60  & \checkmark \\
FC14 & Binomial         & ---   & --- & 200$^\dagger$ & $\triangle$ \\
FC15 & Regression $t$  & 0.05  & 300 & 300 & \checkmark \\
FC16 & Exact enum.     & 0.001 & 840 & 12600 & \checkmark \\
FC17 & Spearman        & 0.15  & 200 & 4000 & \checkmark \\
FC18 & Bootstrap CI    & ---   & 200 & 2250 & \checkmark \\
FC19 & Paired Wilcoxon & 0.50  & 47  & 50  & \checkmark \\
\bottomrule
\end{tabular}
\\[4pt]
{\footnotesize $\triangle$: under-powered; see text for mitigations.
$^\dagger$FC14 uses $B=200$ bootstrap resamples of real Sachs data.}
\end{table}

\emph{Mitigations for under-powered claims:}
\begin{itemize}
  \item FC2: Increase replications from $R=50$ to $R=80$, or
    increase emergent edges per scenario from 5 to 8.
  \item FC13: Increased from $R=50$ to $R=60$ to accommodate
    Bonferroni correction over the expanded claim set.
  \item FC14: Augment with $B=200$ bootstrap resamples of the
    real Sachs data and expand the ground-truth edge set to 12
    edges using recent literature.
\end{itemize}

\subsection{Updated Statistical Protocol}\label{sec:updated-protocol}

With 19 falsifiable claims (up from 14), the multiple-testing
correction must be updated.  We apply the Benjamini--Hochberg (BH)
procedure at FDR level $q=0.10$ across all 19 claims.  The
effective per-claim significance level is approximately
$\alpha_{\mathrm{eff}} = q \cdot k / 19$ for the $k$th-ranked
$p$-value.  For the most significant claims, $\alpha_{\mathrm{eff}}
\approx 0.005$, which is compatible with the individual claim
levels (all at $\alpha \leq 0.05$).

Alternatively, for a more conservative approach, we group claims
into three families: classification (FC1--FC3, FC13), exploration
(FC4--FC9, FC15--FC18), and infrastructure (FC10--FC12, FC14, FC19),
applying Holm--Bonferroni within each family at $\alpha_{\mathrm{fam}}
= 0.05$.  This controls the per-family FWER while allowing more
power within families.


%  Causal-Plasticity Atlas — Sections 6–8, Bibliography, Appendix
%  (No preamble; intended to be \input{} into a master document)
% =============================================================================

% ─────────────────────────────────────────────────────────────────────────────
\section{Experimental Design}\label{sec:experiments}
% ─────────────────────────────────────────────────────────────────────────────

This section specifies the full evaluation protocol for the Causal-Plasticity
Atlas (CPA).  Every claim is paired with a precise falsification condition and
a pre-registered statistical test.  All experiments are constrained to
single-machine CPU execution ($\le 8$ cores, $\le 32$\,GB RAM), use only
freely available data, are reproducible via fixed random seeds, and are
automated end-to-end with no human judgment in the evaluation loop.

% ---------------------------------------------------------------------------
\subsection{Synthetic Benchmarks}\label{sec:synthetic}
% ---------------------------------------------------------------------------

We design four synthetic data generators, each targeting a distinct aspect
of the plasticity classification problem.  For every generator we specify the
data-generating process, ground-truth labels, parameter grids, and four
difficulty levels (Easy, Medium, Hard, Extreme).

% ---- Generator 1 ----
\subsubsection{Generator 1: Fixed Structure, Varying Parameters (FSVP)}

\paragraph{Purpose.}
Test detection of \emph{parametric} plasticity: the DAG skeleton is
constant across all contexts; only the linear coefficients change.

\paragraph{Specification.}
\begin{enumerate}
\item \textbf{DAG construction.}
      Generate a random Erd\H{o}s--R\'enyi DAG with $p\in\{10,20,50,100\}$
      nodes and expected neighbourhood size $d\in\{2,3,4\}$.  Fix a
      topological order $\pi=(1,\dots,p)$ and include each edge $(i\to j)$
      with $i<j$ independently with probability $d/(p-1)$.
\item \textbf{Context generation.}
      For $K\in\{5,10,20\}$ contexts, assign each edge $(i,j)$ a stability
      class:
      \begin{itemize}
        \item \emph{Invariant} (60\% of edges):
              $w_{ij}^{(k)} = w_{ij}^{(1)}$ for all $k$,
              with $w_{ij}^{(1)}\sim\mathrm{Uniform}([-1.0,-0.3]\cup[0.3,1.0])$.
        \item \emph{Plastic} (40\% of edges):
              Draw base $b_{ij}\sim\mathrm{Uniform}([0.3,0.8])$,
              rate $r_{ij}\sim\mathrm{Uniform}([0.05,0.3])$, and
              sign $s\in\{-1,+1\}$ uniformly.  Set
              $w_{ij}^{(k)} = s\bigl(b_{ij} + r_{ij}\,(k-1)/(K-1)\bigr)$.
      \end{itemize}
\item \textbf{Data generation per context $k$.}
      Draw $n\in\{200,500,1000,2000\}$ i.i.d.\ samples from the linear SEM
      \[
        X_j = \sum_{i\in\mathrm{pa}(j)} w_{ij}^{(k)}\,X_i + \varepsilon_j,
        \qquad \varepsilon_j\sim\mathcal{N}(0,\sigma^2),
        \quad \sigma\in\{0.5,1.0,2.0\}.
      \]
\item \textbf{Ground truth.}
      Adjacency matrix $A$ (constant), weight matrices $\{W^{(k)}\}_{k=1}^K$,
      and per-edge labels $\kappa_{ij}\in\{\textsc{invariant},\textsc{plastic}\}$.
\end{enumerate}

\paragraph{Difficulty levels.}
\begin{center}
\small
\begin{tabular}{lcccccc}
\toprule
Level    & $p$ & $d$ & $n$ & $\sigma$ & $K$ & Signal \\
\midrule
Easy     &  10 &  2  & 2000 & 0.5 &  5  & Strong \\
Medium   &  20 &  3  & 1000 & 1.0 & 10  & Moderate \\
Hard     &  50 &  3  &  500 & 1.5 & 20  & Weak \\
Extreme  & 100 &  4  &  200 & 2.0 & 10  & Very weak \\
\bottomrule
\end{tabular}
\end{center}

\paragraph{Confound control.}
All variables are observed; noise is strictly Gaussian with known variance.
Faithfulness is guaranteed by the weight-magnitude lower bound of~$0.3$.

% ---- Generator 2 ----
\subsubsection{Generator 2: Changing Structure Across Contexts (CSAC)}

\paragraph{Purpose.}
Test detection of \emph{structural} plasticity and \emph{emergence}: edges
appear, disappear, or reverse direction across contexts.

\paragraph{Specification.}
\begin{enumerate}
\item \textbf{Base DAG.}
      Generate a random DAG $G_1$ with $p\in\{10,15,30,50\}$ nodes and
      expected degree $d=3$ as in Generator~1.
\item \textbf{Context sequence.}
      For $K\in\{4,8,12\}$ contexts, derive $G_k$ by controlled mutations
      of $G_1$:
      \begin{itemize}
        \item \emph{Invariant} (50\% of $G_1$ edges): present in every $G_k$
              with constant weight.
        \item \emph{Plastic} (25\%): present in every $G_k$ with smoothly
              varying weight (as in FSVP).
        \item \emph{Disappearing} (15\%): present in $G_1,\dots,G_{k^*}$ and
              absent in $G_{k^*+1},\dots,G_K$, where
              $k^*\sim\mathrm{Uniform}\bigl(\{\lceil K/3\rceil,\dots,
              \lceil 2K/3\rceil\}\bigr)$.
        \item \emph{Emergent} ($\le 5$ new edges): absent in early contexts,
              appearing at a random boundary while preserving acyclicity.
      \end{itemize}
\item \textbf{Edge reversals} (at most 2 per scenario):
      select edge $(i\to j)$; at context boundary $k^*$, reverse to
      $(j\to i)$ while maintaining DAG validity.
\item \textbf{Data generation.}
      Same linear SEM as FSVP with $n\in\{500,1000\}$, $\sigma\in\{0.5,1.0\}$.
\item \textbf{Ground truth.}
      Per-context adjacency matrices $\{A^{(k)}\}$, per-edge labels
      $\kappa_{ij}\in\{\textsc{invariant},\textsc{plastic},
      \textsc{emergent},\textsc{disappearing},\textsc{reversing}\}$,
      and transition context indices.
\end{enumerate}

\paragraph{Difficulty levels.}
\begin{center}
\small
\begin{tabular}{lcccccl}
\toprule
Level   & $p$ & $K$ & $n$ & $\sigma$ & Reversals & Note \\
\midrule
Easy    &  10 &  4  & 1000 & 0.5 & 0 & Clear transitions \\
Medium  &  15 &  8  &  500 & 1.0 & 1 & Standard setting \\
Hard    &  30 &  8  &  500 & 1.0 & 2 & High-dimensional \\
Extreme &  50 & 12  &  500 & 1.5 & 2 & Stress test \\
\bottomrule
\end{tabular}
\end{center}

% ---- Generator 3 ----
\subsubsection{Generator 3: Tipping-Point Scenarios (TPS)}

\paragraph{Purpose.}
Test tipping-point detection: smooth parametric drift followed by abrupt
structural change at a critical context index.

\paragraph{Specification.}
\begin{enumerate}
\item \textbf{Context ordering.}
      $K\in\{15,20,30\}$ contexts with a natural total order
      $\theta_1<\cdots<\theta_K$ interpreted as an environmental gradient
      (e.g., temperature, policy dose).
\item \textbf{Tipping-point edges.}
      Select $|T|\in\{2,5,10\}$ edges for tipping-point transitions
      at locations $\theta^*\sim\mathrm{Uniform}([0.3K,0.7K])$.
      Three tipping types are sampled:
      \begin{itemize}
        \item \emph{Collapse} (40\%): edge weight decays linearly to zero
              pre-tip, then the edge disappears at $\theta^*$.
        \item \emph{Bifurcation} (30\%): edge weight constant pre-tip,
              then jumps by $\ge 0.5$ standard deviations and diverges.
        \item \emph{Reversal} (30\%): edge $(i\to j)$ reverses to $(j\to i)$
              at $\theta^*$ while preserving acyclicity.
      \end{itemize}
\item \textbf{Non-tipping edges.}
      Remaining edges undergo smooth linear drift:
      $w_{ij}(\theta_k)=w_{ij}^{(0)}+\beta_{ij}\,\theta_k$
      with $|\beta_{ij}|<0.05$, so mechanism divergence stays below the
      PELT detection threshold.
\item \textbf{Noisy-ordering variant.}
      To test robustness to context mis-ordering, we include a variant
      in which 10\% of context positions are randomly permuted.
\item \textbf{Ground truth.}
      Per-edge indicators $t_{ij}\in\{0,1\}$ and exact tipping indices
      $\theta^*_{ij}$.
\end{enumerate}

\paragraph{Difficulty levels.}
\begin{center}
\small
\begin{tabular}{lccccc}
\toprule
Level   & $K$ & $|T|$ & $n$ & $\sigma$ & Noisy order \\
\midrule
Easy    &  20 &   2   & 1000 & 0.5 & No \\
Medium  &  20 &   5   &  500 & 1.0 & No \\
Hard    &  30 &  10   &  500 & 1.0 & Yes (10\%) \\
Extreme &  30 &  10   &  200 & 1.5 & Yes (20\%) \\
\bottomrule
\end{tabular}
\end{center}

% ---- Generator 4 ----
\subsubsection{Generator 4: Latent Confounders (LC)}

\paragraph{Purpose.}
Assess CPA's robustness under violation of causal sufficiency---the assumption
that all common causes are observed.  This generator was introduced in response
to reviewer critique regarding the faithfulness and sufficiency assumptions.

\paragraph{Specification.}
\begin{enumerate}
\item \textbf{Full DAG.}
      Generate a CSAC-type DAG (Generator~2, Medium difficulty) with
      $p+h$ total variables, where $p\in\{15,30\}$ and
      $h=\lceil\rho\,p\rceil$ with confounding fraction
      $\rho\in\{0.05,0.10,0.15,0.20\}$.
\item \textbf{Latent variable selection.}
      Choose $h$ variables uniformly at random as latent.  Each latent
      variable must be a parent of at least two observed variables to
      induce genuine confounding.
\item \textbf{Observed marginal data.}
      Marginalise over the latent variables.  The observed marginal DAGs
      will contain spurious dependencies between children of latent
      confounders.  Record the ``true'' labels using the full DAG, so
      that an edge classified \textsc{invariant} is truly invariant in
      the latent-augmented model.
\item \textbf{Oracle comparison.}
      Re-run the CPA pipeline with the true DAGs given (bypassing
      discovery) to isolate classification degradation from discovery
      errors.
\end{enumerate}

\paragraph{Difficulty levels.}
\begin{center}
\small
\begin{tabular}{lcccc}
\toprule
Level   & $p$ & $\rho$ & $K$ & $n$ \\
\midrule
Easy    &  15 & 0.05 &  8 & 1000 \\
Medium  &  15 & 0.10 &  8 &  500 \\
Hard    &  30 & 0.15 &  8 &  500 \\
Extreme &  30 & 0.20 & 12 &  500 \\
\bottomrule
\end{tabular}
\end{center}

\paragraph{Expected outcome.}
Classification F1 should degrade by no more than~$0.15$ relative to the
confounder-free setting at $\rho=0.10$; the oracle baseline should show
$\le 0.05$ degradation, isolating errors to the discovery step.

% ---------------------------------------------------------------------------
\subsection{Semi-Synthetic Benchmarks}\label{sec:semisynthetic}
% ---------------------------------------------------------------------------

Semi-synthetic benchmarks combine a real-world causal graph with controlled
parameter perturbations, providing ecological validity absent from purely
synthetic designs.

\subsubsection{Sachs Protein-Signalling Network}
The Sachs network~\cite{sachs2005causal} models intracellular signalling among
$p=11$ phosphoproteins and phospholipids with $17$ causal edges verified by
domain knowledge.  We construct $K=5$ context variants by (i)~selecting 6
edges as plastic and modifying their linear coefficients by factors drawn from
$\mathrm{Uniform}([0.3,2.0])$, (ii)~removing 2 edges in two of the five contexts
to create disappearing edges, and (iii)~adding 1 emergent edge in the final
context.  Sample sizes follow the original flow-cytometry regime:
$n\in\{500,1000,5000\}$ per context.

\subsubsection{Asia Network}
The Asia (Lauritzen--Spiegelhalter) network~\cite{lauritzen1988local} has
$p=8$ binary variables and $8$ directed edges modelling tuberculosis, lung
cancer, and bronchitis.  We create $K=6$ contexts by varying the conditional
probability tables: 4 edges are plastic (CPT entries perturbed by additive
noise $\sim\mathrm{Uniform}([-0.15,0.15])$, clipped to $[0,1]$), and 1 edge
is removed in the final two contexts.  We sample $n\in\{500,2000\}$ per
context.  Despite the small scale, the discrete variable type tests the
generality of the $\sqrt{\mathrm{JSD}}$ distance beyond Gaussian settings.

\subsubsection{Alarm Network}
The Alarm network~\cite{beinlich1989alarm} has $p=37$ variables and $46$ edges,
modelling patient monitoring in an intensive-care unit.  We generate $K=4$
contexts by: (i)~partitioning edges into invariant (60\%), plastic (30\%), and
disappearing (10\%); (ii)~perturbing CPTs of plastic edges with Dirichlet noise
($\alpha=5$); (iii)~removing disappearing edges in contexts 3--4.  With
$n=1000$ per context, this benchmark tests medium-scale structural change
detection on a realistic topology.

\subsubsection{Insurance Network}
The Insurance network~\cite{binder1997adaptive} has $p=27$ variables and $52$
edges modelling automobile insurance risk.  We use $K=8$ contexts with a
richer perturbation schedule: 50\% invariant, 25\% parametrically plastic,
15\% structurally plastic (parent set changes), 10\% emergent.  This benchmark
stresses the full four-class classification capability with
$n\in\{500,1000\}$ per context.

% ---------------------------------------------------------------------------
\subsection{Real-World Validation}\label{sec:realworld}
% ---------------------------------------------------------------------------

To complement the semi-synthetic benchmarks where ground truth is known by
construction, we include a real-world validation study using the original
multi-condition flow-cytometry data of \citet{sachs2005causal}.

\subsubsection{Sachs Multi-Condition Protocol}
The original Sachs dataset consists of $14$ experimental conditions
(9~stimulatory and inhibitory interventions plus 5 observational
conditions) measuring $p=11$ phosphoproteins.  We treat each condition
as a context ($K=14$), using $n=853$ cells per condition (the minimum
across conditions).  Unlike the semi-synthetic Sachs benchmark, no
ground-truth mechanism labels are available; instead we evaluate against
domain-knowledge expectations:
\begin{enumerate}[nosep]
  \item \textbf{Known invariances:} The \mbox{Raf--Mek--Erk} cascade
        should be classified as parametrically plastic (coefficients
        change with stimulation) but structurally invariant across
        most conditions.
  \item \textbf{Known emergence:} The PIP3 $\to$ PIP2 mechanism is
        expected to be active only under anti-CD3/CD28 stimulation.
  \item \textbf{Tipping points:} The PKC $\to$ Raf activation switch
        should exhibit a tipping point between the PMA-stimulated and
        unstimulated conditions.
\end{enumerate}

\noindent\textbf{Protocol.} We run CPA with default parameters and
report: (i)~plasticity classifications for all 17 known edges,
(ii)~agreement with the domain-knowledge expectations above, and
(iii)~any novel mechanism-change patterns not previously described in
the literature.  Results are evaluated qualitatively (expert assessment)
and quantitatively via concordance with the semi-synthetic ground truth
on matching edges.

\begin{quote}
\textbf{FC14 (Real-World Sachs).}
On the original 14-condition Sachs data, CPA correctly classifies
$\ge 70\%$ of edges whose expected plasticity category is known from
domain knowledge (minimum 5 of 7 testable edges).\\
\emph{Falsified if:} concordance $<50\%$ (fewer than 4 of 7).
\end{quote}

% ---------------------------------------------------------------------------
\subsection{Baselines}\label{sec:baselines}
% ---------------------------------------------------------------------------

We compare CPA against nine baselines spanning na\"ive, principled, and
oracle approaches.

\begin{enumerate}[label=\textbf{BL\arabic*.}, leftmargin=2em]
\item \textbf{IND-PHC} (\emph{Independent Per-Context + Post-Hoc Comparison}).
      Run the PC algorithm independently on each context; compare DAGs via
      raw SHD and Jaccard index; classify an edge as invariant if present in
      every context with the same orientation.  No formal certificates.
      This is the na\"ive lower bound.

\item \textbf{ICP} (\emph{Invariant Causal Prediction})~\cite{peters2016causal}.
      Identifies invariant causal predictors across environments with
      binary invariant/non-invariant output and Type-I error control.  
      Cannot detect emergence, structural plasticity, or tipping points.

\item \textbf{Pooled-PC}.
      Pool all context data into a single dataset; run PC once.  Every
      edge is implicitly ``invariant.''  Tests whether multi-context
      analysis adds value over na\"ive aggregation.

\item \textbf{JCI-PC} (\emph{Joint Causal Inference})~\cite{mooij2020joint}.
      Joint causal inference across multiple datasets with context
      variables.  Infers a single DAG with context-dependent edge
      annotations but produces no plasticity descriptors.

\item \textbf{PELT-on-Correlations}.
      Compute pairwise Pearson correlations for each variable pair across
      ordered contexts; apply PELT~\cite{killick2012optimal} changepoint
      detection directly to the correlation time series.  This is a
      ``causal-structure-unaware'' tipping-point baseline.

\item \textbf{Oracle-DAG + CPA-Classifier} (\emph{new from critique}).
      Supply the \emph{true} ground-truth DAGs to the CPA classification
      pipeline (Algorithms~2--5), bypassing causal discovery entirely.
      This oracle isolates the ceiling performance of the classification
      module and quantifies how much CPA performance is limited by
      DAG estimation errors versus classification design.

\item \textbf{CD-NOD} (\emph{Causal Discovery from Nonstationary/heterogeneous
      Data})~\cite{huang2020causal}.
      Exploits distributional shifts across contexts to identify changing
      causal modules, producing a binary changed/unchanged classification
      per variable.  Does not provide a continuous plasticity descriptor,
      structural vs.\ parametric decomposition, or tipping-point localisation.

\item \textbf{Random QD Search}.
      Replace the curiosity-driven CD-QD search (Algorithm~3) with
      uniform random sampling of the genome space
      $\mathcal{G}=2^{[K]}\times 2^{[n]}\times\Phi$.  Same evaluation
      budget (10\,000 evaluations).

\item \textbf{Standard MAP-Elites}~\cite{mouret2015illuminating}.
      MAP-Elites with isoline variation~\cite{vassiliades2018discovering}
      and random mutation, but \emph{without} the curiosity signal
      ($\mu=0$).  Tests whether curiosity-driven exploration adds value.
\end{enumerate}

Table~\ref{tab:baseline-capabilities} summarises which capabilities each
baseline supports.

\begin{table}[t]
\centering
\caption{Baseline capability matrix.  \cmark = supported;
         \xmark = not supported; {\tiny$\sim$} = partial.}
\label{tab:baseline-capabilities}
\small
\begin{tabular}{lccccc}
\toprule
Baseline & \rotatebox{70}{Parametric} &
           \rotatebox{70}{Structural} &
           \rotatebox{70}{Emergence} &
           \rotatebox{70}{Tipping pts} &
           \rotatebox{70}{Certificates} \\
\midrule
IND-PHC           & {\tiny$\sim$} & {\tiny$\sim$} & \xmark & \xmark & \xmark \\
ICP               & \cmark & \xmark & \xmark & \xmark & {\tiny$\sim$} \\
Pooled-PC         & \xmark & \xmark & \xmark & \xmark & \xmark \\
JCI-PC            & {\tiny$\sim$} & {\tiny$\sim$} & \xmark & \xmark & \xmark \\
PELT-on-Corr.     & \xmark & \xmark & \xmark & {\tiny$\sim$} & \xmark \\
Oracle-DAG        & \cmark & \cmark & \cmark & \cmark & \cmark \\
CD-NOD            & {\tiny$\sim$} & \xmark & \xmark & \xmark & \xmark \\
Random QD         & {\tiny$\sim$} & {\tiny$\sim$} & {\tiny$\sim$} & \xmark & \xmark \\
Std.\ MAP-Elites  & {\tiny$\sim$} & {\tiny$\sim$} & {\tiny$\sim$} & \xmark & \xmark \\
\bottomrule
\end{tabular}
\end{table}

% ---------------------------------------------------------------------------
\subsection{Metrics Suite}\label{sec:metrics}
% ---------------------------------------------------------------------------

We evaluate CPA on a suite of twelve metrics, each tied to at least one
falsifiable claim.

\begin{enumerate}[label=\textbf{M\arabic*.}]
\item \textbf{Macro-Averaged F1 Score.}
      Harmonic mean of precision and recall, macro-averaged across all
      plasticity classes $\{\textsc{invariant},\textsc{parametric},
      \textsc{structural},\textsc{emergent}\}$.
      Applies to Claims~1, 2, 8, 11.

\item \textbf{Mean Absolute Deviation (MAD) of Tipping-Point Localisation.}
      For each true tipping point $\theta^*$, compute
      $|\hat\theta-\theta^*|$ where $\hat\theta$ is the nearest detected
      changepoint.  Average over all tipping-point edges and replications.
      Applies to Claim~3.

\item \textbf{Expected Calibration Error (ECE).}
      For certificates with nominal confidence $c$:
      $\mathrm{ECE}=\sum_{b\in\mathcal{B}}\frac{|B_b|}{N}
      |\text{emp}_b - c_b|$ over buckets
      $\mathcal{B}=\{0.70,0.80,0.90,0.95\}$.
      Applies to Claim~4.

\item \textbf{Archive Coverage Ratio.}
      Fraction of reachable CVT cells occupied after a fixed evaluation
      budget.  Reachable cells are estimated by a very long
      ($100\mathrm{K}$~evaluation) random search.
      Applies to Claim~5.

\item \textbf{Normalised Alignment Score (NAS).}
      $\mathrm{NAS} = 1 - \mathrm{SHD}(\pi(G_a),G_b)/\mathrm{SHD}_{\max}$,
      where $\pi$ is the alignment mapping and $\mathrm{SHD}_{\max}$ is
      the maximum possible disagreement.
      Applies to Claim~6.

\item \textbf{Wall-Clock Time (seconds).}
      Total elapsed time for the full pipeline, measured on a single
      CPU core (Intel i7 $\ge 3.0$\,GHz, 32\,GB RAM).
      Applies to Claim~7.

\item \textbf{AUROC for Descriptor Discriminability.}
      Area under the ROC curve using the $\ell^2$-norm of the plasticity
      descriptor $\|\psi_i\|_2$ as a score for invariant-vs-non-invariant
      binary classification.  95\%~CI via DeLong's test.
      Applies to Claim~9.

\item \textbf{Tipping-Point Detection AUC-ROC.}
      Binary classification AUC: for each context index, is there a true
      tipping point?  Detector output is the maximum mechanism divergence
      at that index, aggregated over all edges.
      Applies to Claim~3.

\item \textbf{Structural Hamming Distance (SHD).}
      Number of edge additions, deletions, and reversals needed to
      transform one DAG into another.  Used as the perturbation measure
      in the SHD experiment (Claim~12).

\item \textbf{Descriptor Perturbation ($\ell^\infty$).}
      $\|\hat\psi-\psi\|_\infty =
      \max(|\hat\psi_S-\psi_S|,|\hat\psi_P-\psi_P|,
      |\hat\psi_E-\psi_E|,|\hat\psi_{CS}-\psi_{CS}|)$.
      Validates the bound in Theorem~8.

\item \textbf{Cohen's $d$ Effect Size.}
      Standardised mean difference between CPA and baseline:
      $d = (\bar x_{\text{CPA}}-\bar x_{\text{BL}})/s_{\text{pooled}}$.
      Quantifies practical significance beyond $p$-values.

\item \textbf{Certificate Issuance Rate.}
      Fraction of truly invariant mechanisms for which a certificate is
      successfully issued (true-positive rate of certificate generation).
      Measures the system's \emph{willingness to commit}, not just
      accuracy.
\end{enumerate}

% ---------------------------------------------------------------------------
\subsection{Falsifiable Claims}\label{sec:claims}
% ---------------------------------------------------------------------------

We pre-register thirteen falsifiable claims.  Each specifies a measurable
criterion, a support condition, and a \emph{falsification} condition whose
satisfaction would refute the claim.

\begin{enumerate}[label=\textbf{FC\arabic*.}]
\item \textbf{Parametric Classification Accuracy.}
      On FSVP ($p\!=\!20$, $K\!=\!10$, $n\!=\!1000$), CPA achieves
      macro-averaged F1 $\ge 0.90$ over 50 replications.\\
      \emph{Falsified if:} upper bound of the 95\%~CI falls below $0.85$.\\
      \emph{Supported if:} lower 95\%~CI $>0.90$; improvement over IND-PHC
      $\ge 0.10$ F1 (paired Wilcoxon, $p<0.01$).

\item \textbf{Structural Emergence Detection.}
      On CSAC ($p\!=\!15$, $K\!=\!8$), CPA detects emergent edges with
      precision $\ge 0.85$ and recall $\ge 0.80$.\\
      \emph{Falsified if:} upper 95\%~CI for precision $<0.80$ or recall $<0.75$.\\
      \emph{Supported if:} lower CIs exceed thresholds; context boundary
      correct within $\pm 1$ for $\ge 75\%$ of emergent edges.

\item \textbf{Tipping-Point Localisation.}
      On TPS ($K\!=\!20$), CPA localises tipping points within $\pm 1$
      context index of truth for $\ge 80\%$ of tipping-point edges.\\
      \emph{Falsified if:} mean MAD $>2.0$ or 80th-pctl MAD $>3.0$.\\
      \emph{Supported if:} mean MAD $\le 1.0$, 80th-pctl MAD $\le 2.0$,
      detection AUC-ROC $\ge 0.90$.

\item \textbf{Certificate Calibration.}
      Among edges assigned confidence $c$, the empirical retention rate
      under stability selection lies in $[c-0.05,\,c+0.05]$ for each
      confidence bucket.\\
      \emph{Falsified if:} deviation $>0.10$ for any bucket over 30 replications.\\
      \emph{Supported if:} ECE $\le 0.04$.

\item \textbf{QD Coverage Superiority.}
      Curiosity-driven CD-QD fills $\ge 30\%$ more CVT cells than random
      search and $\ge 15\%$ more than standard MAP-Elites
      (10\,000 evaluations, CSAC $p\!=\!15$, $K\!=\!8$).\\
      \emph{Falsified if:} improvement over random $<15\%$ at upper 95\%~CI.\\
      \emph{Supported if:} improvements exceed thresholds with paired
      Wilcoxon $p<0.01$, Cohen's $d\ge 0.8$.

\item \textbf{DAG Alignment Quality.}
      CADA achieves NAS $\ge 0.85$ on CSAC, outperforming node-identity
      alignment by $\ge 0.10$ and VF2 by $\ge 0.05$.\\
      \emph{Falsified if:} upper 95\%~CI for NAS $<0.80$ or VF2
      improvement not significant ($p\ge 0.05$).

\item \textbf{Scalability.}
      Full pipeline on ($p\!=\!50$, $K\!=\!10$, $n\!=\!500$) completes in
      $\le 30$\,min wall-clock (single core); ($p\!=\!100$, $K\!=\!5$,
      $n\!=\!500$) in $\le 120$\,min.\\
      \emph{Falsified if:} median exceeds $2\times$ the stated limit.\\
      \emph{Supported if:} medians within limits; time-vs-$p$ fits
      polynomial of degree $\le 3$ ($R^2\ge 0.95$).

\item \textbf{Sachs Semi-Synthetic.}
      On Sachs ($p\!=\!11$, $K\!=\!5$), CPA recovers the correct stability
      classification for $\ge 80\%$ of 17 known edges.\\
      \emph{Falsified if:} accuracy $<70\%$ at upper 95\%~CI.

\item \textbf{Descriptor Discriminability.}
      Plasticity descriptors achieve AUROC $\ge 0.90$ for invariant-vs-plastic
      discrimination on FSVP+CSAC combined.\\
      \emph{Falsified if:} AUROC $<0.85$ via DeLong's test.

\item \textbf{End-to-End Reproducibility.}
      Given fixed seeds, two independent runs produce bit-identical class
      assignments and descriptor values identical to $10^{-8}$ tolerance.\\
      \emph{Falsified if:} any of 100 paired runs (10 benchmarks $\times$
      10 seed pairs) differ by $>10^{-6}$.

\item \textbf{Latent-Confounder Graceful Degradation.}
      On LC ($\rho=0.10$), classification F1 decreases by $\le 0.15$
      compared to the confounder-free setting.\\
      \emph{Falsified if:} F1 drop $>0.25$ at lower 95\%~CI over 30
      replications.\\
      \emph{Supported if:} Oracle-DAG shows $\le 0.05$ drop, confirming
      that degradation is due to discovery, not classification.

\item \textbf{Structural Stability Bound.}
      On the SHD perturbation experiment ($s\in\{1,2,5,10\}$), the
      observed $\|\hat\psi-\psi\|_\infty$ is bounded by the theoretical
      $f(s,n,d_{\max})$ from Theorem~8 for $\ge 95\%$ of trials.\\
      \emph{Falsified if:} $>10\%$ of trials violate the bound for any~$s$.

\item \textbf{CPA vs.\ CD-NOD.}
      On CSAC ($p\!=\!15$, $K\!=\!8$), CPA achieves strictly higher
      macro-averaged F1 than CD-NOD for the binary changed/unchanged
      classification (collapsing CPA categories to binary), and additionally
      provides structural vs.\ parametric decomposition unavailable to CD-NOD.\\
      \emph{Falsified if:} CD-NOD F1 $\ge$ CPA F1 at upper 95\%~CI
      (paired Wilcoxon, $p\ge 0.05$).\\
      \emph{Supported if:} CPA F1 exceeds CD-NOD F1 by $\ge 0.05$ with
      $p<0.01$ and Cohen's $d\ge 0.5$.
\end{enumerate}

% ---------------------------------------------------------------------------
\subsection{Statistical Protocol}\label{sec:statprotocol}
% ---------------------------------------------------------------------------

\paragraph{Replication strategy.}
Each benchmark configuration is run for $R=50$ independent replications
(30 for computationally expensive configurations: $p\ge 50$ or $K\ge 20$).
Each replication uses a distinct random seed drawn from the range
$[1,10^6]$ and recorded in the experiment manifest.

\paragraph{Confidence intervals.}
All point estimates are accompanied by 95\% bootstrap confidence intervals
($B=2000$ resamples, bias-corrected and accelerated).

\paragraph{Hypothesis tests.}
Pairwise comparisons use the paired Wilcoxon signed-rank test at significance
level $\alpha=0.01$ with Bonferroni correction for the 9-baseline family.
For AUROC comparisons we use DeLong's test~\cite{delong1988comparing}.
All tests are two-sided unless directionality is specified in the claim.

\paragraph{Effect sizes.}
We report Cohen's $d$ alongside every significance test.  A result is
considered practically significant only if $|d|\ge 0.5$ (medium effect).

\paragraph{Multiple testing.}
The 13 falsifiable claims constitute a single claim family.  We apply the
Benjamini--Hochberg procedure~\cite{benjamini1995controlling} at FDR level
$q=0.05$ to control the expected proportion of falsely supported claims.

\paragraph{Pre-registration.}
The full analysis plan, including all claim thresholds and test specifications,
is committed to the repository before any data are generated.

% ---------------------------------------------------------------------------
\subsection{Ablation Studies}\label{sec:ablations}
% ---------------------------------------------------------------------------

Six ablation studies isolate the contribution of each major module.

\begin{enumerate}[label=\textbf{ABL\arabic*.}]
\item \textbf{Curiosity Signal ($\mu=0$ vs.\ $\mu=0.5$).}
      Remove the curiosity signal from CD-QD search.
      \emph{Expected:} coverage drops $\ge 10\%$; convergence speed drops
      $\ge 20\%$.  Targets FC5.

\item \textbf{DAG Alignment (CADA vs.\ node-identity).}
      Replace CADA with na\"ive index-based alignment.
      \emph{Expected:} F1 drops $\ge 0.10$ on CSAC where variable
      correspondence is non-trivial.  Targets FC1, FC6.

\item \textbf{Certificate Module Removal.}
      Remove Algorithm~5 (RCG); report bare classifications.
      \emph{Expected:} $\ge 15\%$ of ``invariant'' labels overturned under
      bootstrap.  Targets FC4.

\item \textbf{Descriptor Component Ablation.}
      Remove each of $\psi_S$, $\psi_P$, $\psi_E$, $\psi_{CS}$
      individually.
      \emph{Expected:} removing $\psi_S$ degrades structural-change F1 by
      $\ge 0.20$; removing $\psi_E$ removes all emergence detection.
      Targets FC1, FC9.

\item \textbf{Tipping-Point Detector (PELT vs.\ sliding window).}
      Replace PELT with a simple sliding-window $t$-test
      (window size~3).
      \emph{Expected:} PELT achieves $\ge 0.10$ better AUC-ROC, especially
      for multiple simultaneous tipping points.  Targets FC3.

\item \textbf{Mechanism Distance ($\sqrt{\mathrm{JSD}}$ vs.\ Euclidean).}
      Replace $\sqrt{\mathrm{JSD}}$ with Euclidean distance on regression
      coefficients.
      \emph{Expected:} $\sqrt{\mathrm{JSD}}$ outperforms Euclidean by
      $\ge 0.05$ F1 and $\ge 0.05$ AUC-ROC.  Targets FC1, FC3, FC9.
\end{enumerate}


% ─────────────────────────────────────────────────────────────────────────────
\section{Implementation Architecture}\label{sec:implementation}
% ─────────────────────────────────────────────────────────────────────────────

This section describes the modular software architecture that realises the
theory of Sections~3--5.  The implementation is a pure-Python library
(Python $\ge 3.10$) with dependencies limited to NumPy, SciPy,
scikit-learn, and \texttt{causal-learn}~\cite{zheng2024causal}.

% ---------------------------------------------------------------------------
\subsection{Module Hierarchy}
% ---------------------------------------------------------------------------

The codebase is organised into six top-level packages, each mapped to the
theoretical constructs they implement:

\begin{center}
\small
\begin{tabular}{lll}
\toprule
Package & Theory & Primary algorithms \\
\midrule
\texttt{core.data\_model}         & Def.\ 1--3     & MCCM data structures \\
\texttt{core.mechanism\_distance} & Def.\ 4--5, Thm.\ 1 & $\sqrt{\mathrm{JSD}}$ computation \\
\texttt{alignment.cada}           & Def.\ 6, Thm.\ 1, 7 & Alg.\ 1 (CADA) \\
\texttt{descriptors.plasticity}   & Def.\ 7--8, Thm.\ 2, 4 & Alg.\ 2 (PDC) \\
\texttt{exploration.qd\_search}   & Def.\ 9--12, Thm.\ 6 & Alg.\ 3 (CD-QD) \\
\texttt{detection.tipping\_points}& Def.\ 13, Thm.\ 5   & Alg.\ 4 (TPD) \\
\texttt{certificates.robustness}  & Def.\ 14, Thm.\ 3, 8 & Alg.\ 5 (RCG) \\
\texttt{pipeline.orchestrator}    & --- & Three-phase pipeline \\
\texttt{diagnostics.sensitivity}  & Thm.\ 8 & SHD perturbation testing \\
\texttt{benchmarks.generators}    & --- & Generators 1--4, semi-synthetic \\
\bottomrule
\end{tabular}
\end{center}

Each module exposes a narrow public API (typically a single \texttt{run()}
entry-point) and communicates with other modules exclusively through the
MCCM data structures defined in \texttt{core.data\_model}.

% ---------------------------------------------------------------------------
\subsection{Theorem-to-Module Mapping}
% ---------------------------------------------------------------------------

We make explicit which formal result justifies each computational step:

\begin{itemize}
\item \textbf{Theorem~1} (DAG alignment metric properties) guarantees that the
      alignment cost in \texttt{alignment.cada} satisfies the pseudo-metric
      axioms.  This justifies the use of $d_{\mathrm{DAG}}$ as a
      well-defined distance in downstream CVT tessellation.
\item \textbf{Theorem~2} (classification correctness) provides the
      sample-complexity conditions that \texttt{descriptors.plasticity}
      checks before issuing a classification.  If
      $n_c < n_{\min}(\varepsilon,\alpha)$, the module flags the
      classification as ``provisional.''
\item \textbf{Theorem~3} (certificate soundness) is the logical backbone of
      \texttt{certificates.robustness}: the margin condition
      $\delta_i > \varepsilon_{\text{est}}$ is checked explicitly, and
      certificates are only issued when satisfied.
\item \textbf{Theorem~5} (tipping-point consistency) justifies the PELT
      cost function used in \texttt{detection.tipping\_points}: the
      BIC-penalised cost guarantees consistency as $K\to\infty$ under
      the piecewise-Lipschitz assumption.
\item \textbf{Theorem~6} (atlas completeness) motivates the ergodic mutation
      operator in \texttt{exploration.qd\_search}.  The mutation
      kernel assigns positive probability to every genome in $\mathcal{G}$,
      ensuring eventual coverage of all reachable CVT cells.
\item \textbf{Theorem~7} (NP-hardness / FPT tractability) dictates the
      algorithmic strategy in \texttt{alignment.cada}: for bounded-degree
      DAGs ($d_{\max}\le 6$), the six-phase heuristic provably recovers
      the optimal alignment in polynomial time.
\item \textbf{Theorem~8} (structural stability) is implemented in
      \texttt{diagnostics.sensitivity}, which computes the theoretical
      perturbation bound $f(s,n,d_{\max})$ and compares it with
      empirical descriptor shifts under controlled SHD injection.
\end{itemize}

% ---------------------------------------------------------------------------
\subsection{Data Flow}\label{sec:dataflow}
% ---------------------------------------------------------------------------

The pipeline proceeds in three phases:

\paragraph{Phase~1: Foundation.}
Raw multi-context data $\{D_c\}_{c=1}^K$ enter the pipeline.  An external
causal-discovery algorithm (PC or GES via \texttt{causal-learn}) estimates
per-context DAGs $\{\hat G_c\}$.  CADA aligns all $\binom{K}{2}$ DAG pairs,
populating a pairwise alignment cache $\mathcal{A}$ of size $O(K^2)$.  PDC
then reads $\mathcal{A}$ and the estimated conditional distributions to
compute plasticity descriptors $\{\psi_i\}_{i=1}^{|\bar{V}|}$ and initial
classifications $\{\kappa_i\}$.

\paragraph{Phase~2: Exploration.}
CD-QD search operates on the genome space
$\mathcal{G}=2^{[K]}\times 2^{[n]}\times\Phi$, where each genome encodes a
context subset, a variable subset, and descriptor-computation parameters.
The quality function evaluates classification confidence and robustness;
the behaviour descriptor maps each genome to its 4D plasticity signature.
The archive is tessellated into $L=1000$ CVT cells via Lloyd's algorithm
($50$ iterations).  CD-QD runs for a configurable number of iterations
(default $10\,000$), emitting the populated atlas $\mathcal{A}_T$.

\paragraph{Phase~3: Validation.}
TPD scans the mechanism-divergence sequence
$\{d_\mu(\hat P^{(k)}_i,\hat P^{(k+1)}_i)\}_{k=1}^{K-1}$ for each
variable $i$, applying PELT with BIC penalty and permutation validation
($B_{\text{perm}}=200$, Benjamini--Hochberg FDR at $q=0.05$).
RCG issues robustness certificates by: (i)~stability selection
(50\% subsample $\times$ 100 rounds) for structural invariance;
(ii)~parametric bootstrap ($B=500$) for mechanism-distance UCBs.
The final atlas assembles classifications, certificates, and tipping-point
annotations into a single JSON output.

% ---------------------------------------------------------------------------
\subsection{Computational Budget Allocation}\label{sec:budget}
% ---------------------------------------------------------------------------

For the practical regime ($p=50$, $K=10$, $n=500$ per context), the
approximate wall-clock budget on a single Intel i7 core is:

\begin{center}
\small
\begin{tabular}{lrc}
\toprule
Component & Time & \% \\
\midrule
Phase~1: Causal discovery ($K\times$\,GES) & 8\,min & 27\% \\
Phase~1: CADA alignment ($\binom{K}{2}$ pairs) & 4\,min & 13\% \\
Phase~1: PDC descriptors & 1\,min & 3\% \\
Phase~2: CD-QD search ($10^4$ evals) & 2\,min & 7\% \\
Phase~3: TPD tipping-points & $<\!1$\,min & 2\% \\
Phase~3: RCG certificates (stability sel.\ + bootstrap) & 14\,min$^{\dagger}$ & 47\% \\
Overhead (I/O, CVT, assembly) & $<\!1$\,min & 1\% \\
\midrule
\textbf{Total} & \textbf{$\sim$30\,min} & 100\% \\
\bottomrule
\end{tabular}
\end{center}

Certificate generation dominates the budget because stability selection
requires $100\times K$ DAG re-estimations per mechanism class.  This is
an order-of-magnitude improvement over the full-bootstrap alternative
($500\times K^2$ re-estimations), achieved by adopting stability selection
per the reviewer critique.

\smallskip\noindent${}^{\dagger}$The 14\,min figure assumes sequential
processing of all $n=50$ mechanisms with $R_{\mathrm{stab}}=100$
rounds.  In practice, the structural check is skipped for mechanisms
whose parent sets agree across all contexts (typically $\sim 60\%$),
reducing the effective cost.  The per-mechanism complexity
$O(R_{\mathrm{stab}}\cdot K\cdot T_{\mathrm{DAG}}+K^2\cdot B\cdot
d_{\max})$ yields $\sim 4$\,hours if \emph{all} mechanisms require
stability selection without skipping; the 14\,min estimate reflects the
typical skip rate of $\sim 60\%$.

% ---------------------------------------------------------------------------
\subsection{Parallelisation Strategy}\label{sec:parallel}
% ---------------------------------------------------------------------------

Three levels of parallelism are exploited:

\begin{enumerate}
\item \textbf{Per-context parallelism.}
      Phase~1 discovery runs $K$ independent GES instances concurrently
      (embarrassingly parallel; linear speedup up to $K$ cores).

\item \textbf{Per-pair parallelism.}
      CADA alignment of $\binom{K}{2}$ pairs is independent and
      distributable across cores.  With 8 cores, alignment time is
      reduced by $\sim 7\times$.

\item \textbf{Per-mechanism parallelism.}
      Certificate generation for each mechanism is independent.
      With $|\bar V|=50$ mechanisms and 8 cores, RCG time drops from
      14~min to $\sim 2$~min.
\end{enumerate}

With 8-core parallelism, the total pipeline for ($p\!=\!50$, $K\!=\!10$)
drops from $\sim 30$\,min to $\sim 6$\,min, well within the stated
scalability targets.

% ---------------------------------------------------------------------------
\subsection{Engineering Challenges}\label{sec:challenges}
% ---------------------------------------------------------------------------

Several components present non-trivial engineering challenges:

\paragraph{CADA alignment at scale.}
Although Theorem~7 guarantees FPT tractability for bounded-degree DAGs,
the six-phase heuristic requires careful tuning of the CI-fingerprint
scoring (Phase~3) to avoid $O(p^3)$ behaviour in the Hungarian matching
step.  We mitigate this by pre-filtering candidate correspondences using
Markov-blanket overlap, reducing the bipartite matching to at most
$O(d_{\max}^2\cdot p)$ candidates.

\paragraph{Stability selection at scale.}
Each round of stability selection requires a half-sample DAG re-estimation.
For $p=100$ with GES, this takes $\sim 3$\,s per round.  With
$n_{\text{inv}}\approx 60$ invariant mechanisms and $100$ rounds each,
the sequential cost is $60\times 100\times 3\,\text{s}\approx 5$\,hours.
Per-mechanism parallelism (8 cores) reduces this to $\sim 40$\,min.
Further acceleration via caching of shared sufficient statistics across
subsamples brings the cost to $\sim 15$\,min.

\paragraph{CVT tessellation in 4D.}
Lloyd's algorithm for 1000 cells in $\mathbb{R}^4$ converges slowly
(requiring $\sim 50$ iterations with $10^5$ samples per iteration).
We initialise with $k$-means++ and terminate when cell centroids move
by $<10^{-4}$ in $\ell^2$-norm between iterations.

\paragraph{Numerical stability of $\sqrt{\mathrm{JSD}}$.}
Computing $\sqrt{\mathrm{JSD}}$ between Gaussian conditionals admits
a closed-form expression involving log-determinants of covariance
matrices.  For near-singular covariance (e.g., $n<p$), we add a
Tikhonov regulariser $\lambda I$ with $\lambda=10^{-6}$ and verify
that the resulting distance remains a valid metric up to
$O(\lambda)$~perturbation.

% ---------------------------------------------------------------------------
\subsection{Expanded Implementation Summary}\label{sec:expanded-impl}
% ---------------------------------------------------------------------------

The complete implementation comprises 84{,}338 non-empty lines of Python
across 217 files (125 source modules, 83 test modules, 9 configuration/utility),
with 2{,}198 unit tests passing. Beyond the core packages described above,
thirteen additional packages realise the full algorithmic and evaluation
infrastructure:

\begin{center}
\small
\begin{tabular}{llrc}
\toprule
Package & Purpose & Files & Lines \\
\midrule
\texttt{cpa/inference/}     & Do-calculus, counterfactuals, ID algorithm     & 4 & 3{,}190 \\
\texttt{cpa/mec/}           & CPDAG/PAG, Meek rules, d-sep (Bayes-Ball), MEC enum. & 5 & 2{,}555 \\
\texttt{cpa/ci\_tests/}     & Fisher-z, KCI/HSIC, $\chi^2$/G-test, CMI (KSG), adaptive & 5 & 2{,}147 \\
\texttt{cpa/scores/}        & BIC/eBIC, BGe, BDeu, interventional BIC, score caching  & 5 & 1{,}845 \\
\texttt{cpa/sampling/}      & Order/partition/structure MCMC, parallel tempering, DAG bootstrap & 5 & 2{,}430 \\
\texttt{cpa/operators/}     & MB crossover, edge-preserving crossover, GES local search & 4 & 1{,}839 \\
\texttt{cpa/baselines/}     & IND-PHC, pooled, ICP, CD-NOD, JCI, GES, LSEM-pool       & 7 & 2{,}596 \\
\texttt{cpa/analysis/}      & Convergence, supermartingale stopping, ergodicity, Sobol/Morris & 5 & 1{,}919 \\
\texttt{cpa/data/}          & SB1--SB5 generators, nonlinear SCM, semi-synthetic, perturbation & 4 & 1{,}723 \\
\texttt{cpa/streaming/}     & IPU (Algorithm~6), online alignment, windowed CUSUM      & 4 & 1{,}703 \\
\texttt{cpa/scalability/}   & Parent-set caching, approx.\ BIC, sparse DAG ops         & 4 & 1{,}429 \\
\texttt{cpa/config/}        & Experiment runner, hyperparameter tuning, component registry & 3 & 976 \\
\texttt{cpa/certificates/}  & Fisher info.\ Lipschitz, Boltzmann stability, ACT (Alg.~7) & 3$^{*}$ & --- \\
\bottomrule
\end{tabular}
\end{center}

\noindent${}^{*}$Three new files added to the existing
\texttt{certificates.robustness} package.

\paragraph{Algorithm--module mapping.}
All seven algorithms specified in the theory are implemented:
Algorithms~1--5 in the core packages (Section~\ref{sec:implementation}),
Algorithm~6 (Incremental Plasticity Update) in \texttt{cpa/streaming/},
and Algorithm~7 (Asymptotic Certificate Tightening) in \texttt{cpa/certificates/}.
The \texttt{cpa/baselines/} package implements all seven evaluation baselines
(IND-PHC, pooled, ICP, CD-NOD, JCI, GES, LSEM-pool), and
\texttt{cpa/data/} implements all five synthetic benchmarks SB1--SB5 together
with the semi-synthetic benchmarks (Sachs, ALARM, Insurance).


% ─────────────────────────────────────────────────────────────────────────────
\section{Discussion and Conclusion}\label{sec:discussion}
% ─────────────────────────────────────────────────────────────────────────────

% ---------------------------------------------------------------------------
\subsection{Summary of Contributions}
% ---------------------------------------------------------------------------

This paper introduces the \emph{Causal-Plasticity Atlas} (CPA), a framework
for systematically classifying causal mechanisms as invariant, parametrically
plastic, structurally plastic, or emergent across heterogeneous observational
contexts.  The contributions span three levels:

\begin{enumerate}
\item \textbf{Formal foundations.}
      We define the Multi-Context Causal Model (MCCM) with variable-set
      mismatch as a first-class citizen (Definition~3), introduce the
      $\sqrt{\mathrm{JSD}}$-based DAG alignment distance as a proper
      pseudo-metric (Theorem~1), and provide a four-dimensional plasticity
      descriptor that decomposes mechanism change along orthogonal axes
      (Definition~7).  Robustness certificates with explicit soundness
      guarantees (Theorem~3) ensure that classifications carry calibrated
      confidence.

\item \textbf{Algorithmic pipeline.}
      Five algorithms---CADA alignment, PDC descriptor computation,
      CD-QD quality-diversity search, PELT-based tipping-point detection,
      and stability-selection-based certificate generation---are woven into
      a three-phase pipeline with $O(K^2 p^2 d^2)$ alignment cost and
      sub-quadratic certificate cost via stability selection.

\item \textbf{Rigorous evaluation protocol.}
      Twelve falsifiable claims, four synthetic generators (including
      a latent-confounder generator added in response to critique),
      four semi-synthetic benchmarks, eight baselines (including an
      oracle-DAG ceiling), twelve metrics, and six ablation studies
      provide a comprehensive and pre-registered evaluation framework.
\end{enumerate}

% ---------------------------------------------------------------------------
\subsection{Limitations}
% ---------------------------------------------------------------------------

We acknowledge several important limitations:

\paragraph{Causal sufficiency.}
The framework assumes approximate causal sufficiency (Assumption~A3).
While Generator~4 (LC) demonstrates graceful degradation under moderate
latent confounding ($\rho\le 0.10$), the theoretical guarantees of
Theorems~2--3 no longer hold when $\rho > 0.15$.  In domains where
latent confounders are pervasive (e.g., social science, genomics),
the plasticity classifications should be interpreted as conditional on
the discovered DAGs, not as causal ground truth.

\paragraph{Faithfulness.}
Theorem~2 requires the faithfulness assumption (A2), which can be violated
in practice by near-cancellation of causal paths or deterministic
relationships.  The minimum edge-weight bound of $0.3$ in our synthetic
generators ensures faithfulness by construction, but real-world data may
contain near-unfaithful configurations that degrade classification accuracy.

\paragraph{Linear Gaussian restriction.}
The primary evaluation uses linear Gaussian SEMs.  While $\sqrt{\mathrm{JSD}}$
is defined for arbitrary distributions, the closed-form Gaussian computation
is substantially faster than kernel-based estimation.  Extension to
nonlinear mechanisms requires nonparametric conditional density estimation,
whose sample complexity scales poorly with dimensionality.  A nonlinear
benchmark (Generator~5 in the approach specification) is planned but not
included in the primary evaluation.

\paragraph{Context ordering for tipping points.}
Tipping-point detection (Algorithm~4) requires a natural ordering of
contexts.  When contexts lack a canonical order (e.g., geographic regions
rather than time points), the tipping-point module is inapplicable.
The remaining modules (alignment, classification, certificates) remain
valid for unordered context sets.

\paragraph{Scalability ceiling.}
The practical regime is $p\in[10,100]$ variables and $K\in[3,50]$ contexts.
Beyond $p=200$ the pairwise alignment cache becomes memory-intensive
($O(K^2 p^2)\sim 1$\,GB for $K=20$, $p=200$), and certificate generation
time grows super-linearly.  High-dimensional applications would require
variable screening or hierarchical context clustering, neither of which is
analysed theoretically in this work.

\paragraph{QD archive interpretability.}
The atlas organises mechanisms by their four-dimensional plasticity signature,
but the downstream \emph{use} of the archive by domain scientists remains
under-specified.  We provide coverage and quality metrics but do not evaluate
whether the archive improves scientific decision-making in a user study.

% ---------------------------------------------------------------------------
\subsection{Connection to Transportability}
% ---------------------------------------------------------------------------

The CPA framework can be viewed through the lens of the
transportability calculus of Bareinboim and
Pearl~\cite{bareinboim2016causal}.  In the transportability framework,
a causal effect is \emph{transportable} from source to target population
when the intervening mechanisms are invariant---that is, when the
selection diagram indicates no context-specific perturbation.  CPA's
invariant classification corresponds precisely to transportable mechanisms:
an edge classified \textsc{invariant} with a high-confidence certificate
is one whose conditional distribution is statistically indistinguishable
across all observed contexts, suggesting (but not guaranteeing) that the
mechanism would remain stable under transport to a new, unseen context.

The plastic and emergent categories extend transportability analysis
by characterising \emph{how} a mechanism changes.  A parametrically plastic
mechanism may still be partially transportable if the parameter variation
follows a predictable trend (e.g., linear drift with context), enabling
extrapolation-based transport.  A structurally plastic mechanism requires
per-context causal reasoning and is generally non-transportable without
domain-specific knowledge.  An emergent mechanism signals that the causal
vocabulary itself changes across populations---a scenario that
transportability theory addresses via selection diagrams but that has
received limited computational attention.

We emphasise that CPA operates in a purely observational setting and
does not leverage interventional data.  The transportability framework is
more powerful when interventional experiments are available; CPA provides
a complementary analysis when only multi-context observational data exist.

% ---------------------------------------------------------------------------
\subsection{Future Work}
% ---------------------------------------------------------------------------

Several directions merit investigation:

\begin{enumerate}
\item \textbf{Interventional contexts.}
      Incorporating contexts with known interventions (as in
      ICP~\cite{peters2016causal} and JCI~\cite{mooij2020joint})
      would strengthen both discovery and classification.  The
      MCCM framework naturally accommodates intervention targets
      as context metadata.

\item \textbf{Nonlinear and non-Gaussian mechanisms.}
      Extending the $\sqrt{\mathrm{JSD}}$ computation to nonparametric
      settings via kernel density estimation or normalising flows would
      broaden applicability.  The formal framework is distribution-agnostic;
      only the computational module requires extension.

\item \textbf{Temporal causal dynamics.}
      The current framework treats contexts as exchangeable (except for
      tipping-point detection).  Modelling temporal dependencies between
      contexts---e.g., via causal time-series models or
      dynamic Bayesian networks---would enable tracking of mechanism
      evolution over time.

\item \textbf{Active context selection.}
      Given a partially populated atlas, which new context would maximise
      information gain?  Integrating the CPA with Bayesian experimental
      design could yield an active-learning loop for causal mechanism
      exploration.

\item \textbf{Causal representation learning.}
      When the variable set is not given a priori (e.g., in image or
      text domains), causal representation
      learning~\cite{scholkopf2021toward} could provide the variable
      extraction layer upstream of the CPA pipeline.

\item \textbf{User studies.}
      Evaluating whether the atlas improves domain scientists' ability
      to identify mechanism changes, compared to standard per-context
      causal-discovery output, requires human-subjects experiments.
\end{enumerate}

% ---------------------------------------------------------------------------
\subsection{Broader Impact}
% ---------------------------------------------------------------------------

The ability to systematically map how causal mechanisms change across
contexts has implications across several domains.  In \emph{medicine},
understanding which drug--response pathways are invariant across patient
subgroups versus plastic with respect to genotype or comorbidity can inform
precision-medicine strategies.  In \emph{climate science}, identifying
tipping points in ecosystem causal networks---e.g., the collapse of a
pollinator--plant interaction under temperature stress---can support early
warning systems.  In \emph{economics}, distinguishing invariant structural
relationships (e.g., supply--demand) from context-sensitive ones
(e.g., consumer behaviour under policy regimes) can improve the external
validity of policy recommendations.

We note potential for misuse: plasticity classifications based on noisy
observational data could be over-interpreted as definitive causal knowledge.
The robustness certificates provide a partial safeguard by quantifying
uncertainty, but users must understand that certificates are conditional on
the assumptions (faithfulness, approximate sufficiency) and the quality of
the underlying causal-discovery algorithm.  We recommend that CPA outputs
be treated as hypothesis-generating rather than hypothesis-confirming, and
that classifications with low-confidence certificates be subjected to
domain-expert review before informing consequential decisions.


% ─────────────────────────────────────────────────────────────────────────────
%  Bibliography
% ─────────────────────────────────────────────────────────────────────────────

\begin{thebibliography}{99}

\bibitem{pearl2009causality}
J.~Pearl.
\newblock \emph{Causality: Models, Reasoning, and Inference}.
\newblock Cambridge University Press, 2nd edition, 2009.

\bibitem{spirtes2000causation}
P.~Spirtes, C.~Glymour, and R.~Scheines.
\newblock \emph{Causation, Prediction, and Search}.
\newblock MIT Press, 2nd edition, 2000.

\bibitem{peters2016causal}
J.~Peters, P.~B{\"u}hlmann, and N.~Meinshausen.
\newblock Causal inference by using invariant prediction: identification and
  confidence intervals.
\newblock \emph{Journal of the Royal Statistical Society: Series B},
  78(5):947--1012, 2016.

\bibitem{mooij2020joint}
J.~M.~Mooij, S.~Magliacane, and T.~Claassen.
\newblock Joint causal inference from multiple contexts.
\newblock \emph{Journal of Machine Learning Research}, 21(99):1--108, 2020.

\bibitem{mouret2015illuminating}
J.-B.~Mouret and J.~Clune.
\newblock Illuminating search spaces by mapping elites.
\newblock \emph{arXiv preprint arXiv:1504.04909}, 2015.

\bibitem{killick2012optimal}
R.~Killick, P.~Fearnhead, and I.~A.~Eckley.
\newblock Optimal detection of changepoints with a linear computational cost.
\newblock \emph{Journal of the American Statistical Association},
  107(500):1590--1598, 2012.

\bibitem{bareinboim2016causal}
E.~Bareinboim and J.~Pearl.
\newblock Causal inference and the data-fusion problem.
\newblock \emph{Proceedings of the National Academy of Sciences},
  113(27):7345--7352, 2016.

\bibitem{scholkopf2012causal}
B.~Sch{\"o}lkopf, D.~Janzing, J.~Peters, E.~Sgouritsa, K.~Zhang, and
  J.~Mooij.
\newblock On causal and anticausal learning.
\newblock In \emph{Proceedings of the 29th International Conference on Machine
  Learning (ICML)}, pp.~459--466, 2012.

\bibitem{meinshausen2010stability}
N.~Meinshausen and P.~B{\"u}hlmann.
\newblock Stability selection.
\newblock \emph{Journal of the Royal Statistical Society: Series B},
  72(4):417--473, 2010.

\bibitem{sachs2005causal}
K.~Sachs, O.~Perez, D.~Pe'er, D.~A.~Lauffenburger, and G.~P.~Nolan.
\newblock Causal protein-signaling networks derived from multiparameter
  single-cell data.
\newblock \emph{Science}, 308(5721):523--529, 2005.

\bibitem{lauritzen1988local}
S.~L.~Lauritzen and D.~J.~Spiegelhalter.
\newblock Local computations with probabilities on graphical structures and
  their application to expert systems.
\newblock \emph{Journal of the Royal Statistical Society: Series B},
  50(2):157--194, 1988.

\bibitem{beinlich1989alarm}
I.~A.~Beinlich, H.~J.~Suermondt, R.~M.~Chavez, and G.~F.~Cooper.
\newblock The {ALARM} monitoring system: A case study with two probabilistic
  inference techniques for belief networks.
\newblock In \emph{Proceedings of the Second European Conference on Artificial
  Intelligence in Medicine}, pp.~247--256, 1989.

\bibitem{binder1997adaptive}
J.~Binder, D.~Koller, S.~Russell, and K.~Kanazawa.
\newblock Adaptive probabilistic networks with hidden variables.
\newblock \emph{Machine Learning}, 29(2):213--244, 1997.

\bibitem{endres2003new}
D.~M.~Endres and J.~E.~Schindelin.
\newblock A new metric for probability distributions.
\newblock \emph{IEEE Transactions on Information Theory}, 49(7):1858--1860,
  2003.

\bibitem{friedman1999data}
N.~Friedman, M.~Goldszmidt, and A.~Wyner.
\newblock Data analysis with {B}ayesian networks: A bootstrap approach.
\newblock In \emph{Proceedings of the 15th Conference on Uncertainty in
  Artificial Intelligence (UAI)}, pp.~196--205, 1999.

\bibitem{chickering2002optimal}
D.~M.~Chickering.
\newblock Optimal structure identification with greedy search.
\newblock \emph{Journal of Machine Learning Research}, 3:507--554, 2002.

\bibitem{zheng2018dags}
X.~Zheng, B.~Aragam, P.~Ravikumar, and E.~P.~Xing.
\newblock {DAGs} with {NO TEARS}: Continuous optimization for structure
  learning.
\newblock In \emph{Advances in Neural Information Processing Systems (NeurIPS)},
  pp.~9472--9483, 2018.

\bibitem{hauser2012characterization}
A.~Hauser and P.~B{\"u}hlmann.
\newblock Characterization and greedy learning of interventional {M}arkov
  equivalence classes of directed acyclic graphs.
\newblock \emph{Journal of Machine Learning Research}, 13:2409--2464, 2012.

\bibitem{tian2001causal}
J.~Tian and J.~Pearl.
\newblock Causal discovery from changes.
\newblock In \emph{Proceedings of the 17th Conference on Uncertainty in
  Artificial Intelligence (UAI)}, pp.~512--521, 2001.

\bibitem{zhang2017causal}
K.~Zhang, B.~Huang, J.~Zhang, C.~Glymour, and B.~Sch{\"o}lkopf.
\newblock Causal discovery from nonstationary/heterogeneous data: Skeleton
  estimation and orientation determination.
\newblock In \emph{Proceedings of the 26th International Joint Conference on
  Artificial Intelligence (IJCAI)}, pp.~1347--1353, 2017.

\bibitem{huang2020causal}
B.~Huang, K.~Zhang, J.~Zhang, J.~Ramsey, R.~Sanchez-Romero, C.~Glymour, and
  B.~Sch{\"o}lkopf.
\newblock Causal discovery from heterogeneous/nonstationary data.
\newblock \emph{Journal of Machine Learning Research}, 21(89):1--53, 2020.

\bibitem{rothenhausler2021anchor}
D.~Rothenh{\"a}usler, N.~Meinshausen, P.~B{\"u}hlmann, and J.~Peters.
\newblock Anchor regression: Heterogeneous data meet causality.
\newblock \emph{Journal of the Royal Statistical Society: Series B},
  83(2):215--246, 2021.

\bibitem{heinzedeml2018invariant}
C.~Heinze-Deml, J.~Peters, and N.~Meinshausen.
\newblock Invariant causal prediction for nonlinear models.
\newblock \emph{Journal of Causal Inference}, 6(2):20170016, 2018.

\bibitem{pfister2019invariant}
N.~Pfister, P.~B{\"u}hlmann, and J.~Peters.
\newblock Invariant causal prediction for sequential data.
\newblock \emph{Journal of the American Statistical Association},
  114(527):1264--1276, 2019.

\bibitem{adams2007bayesian}
R.~P.~Adams and D.~J.~C.~MacKay.
\newblock Bayesian online changepoint detection.
\newblock \emph{arXiv preprint arXiv:0710.3742}, 2007.

\bibitem{benjamini1995controlling}
Y.~Benjamini and Y.~Hochberg.
\newblock Controlling the false discovery rate: A practical and powerful
  approach to multiple testing.
\newblock \emph{Journal of the Royal Statistical Society: Series B},
  57(1):289--300, 1995.

\bibitem{delong1988comparing}
E.~R.~DeLong, D.~M.~DeLong, and D.~L.~Clarke-Pearson.
\newblock Comparing the areas under two or more correlated receiver operating
  characteristic curves: A nonparametric approach.
\newblock \emph{Biometrics}, 44(3):837--845, 1988.

\bibitem{vassiliades2018discovering}
V.~Vassiliades, K.~Chatzilygeroudis, and J.-B.~Mouret.
\newblock Using centroidal {V}oronoi tessellations to scale up the
  multidimensional archive of phenotypic elites algorithm.
\newblock \emph{IEEE Transactions on Evolutionary Computation},
  22(4):623--630, 2018.

\bibitem{lehman2011novelty}
J.~Lehman and K.~O.~Stanley.
\newblock Abandoning objectives: Evolution through the search for novelty alone.
\newblock \emph{Evolutionary Computation}, 19(2):189--223, 2011.

\bibitem{scholkopf2021toward}
B.~Sch{\"o}lkopf, F.~Locatello, S.~Bauer, N.~R.~Ke, N.~Kalchbrenner,
  A.~Goyal, and Y.~Bengio.
\newblock Toward causal representation learning.
\newblock \emph{Proceedings of the IEEE}, 109(5):612--634, 2021.

\bibitem{vovk2005algorithmic}
V.~Vovk, A.~Gammerman, and G.~Shafer.
\newblock \emph{Algorithmic Learning in a Random World}.
\newblock Springer, 2005.

\bibitem{schwarz1978estimating}
G.~Schwarz.
\newblock Estimating the dimension of a model.
\newblock \emph{The Annals of Statistics}, 6(2):461--464, 1978.

\bibitem{shah2013variable}
R.~D.~Shah and R.~J.~Samworth.
\newblock Variable selection with error control: Another look at stability
  selection.
\newblock \emph{Journal of the Royal Statistical Society: Series B},
  75(1):55--80, 2013.

\bibitem{gao2010survey}
X.~Gao, B.~Xiao, D.~Tao, and X.~Li.
\newblock A survey of graph edit distance.
\newblock \emph{Pattern Analysis and Applications}, 13(1):113--129, 2010.

\bibitem{lin1991divergence}
J.~Lin.
\newblock Divergence measures based on the {S}hannon entropy.
\newblock \emph{IEEE Transactions on Information Theory}, 37(1):145--151, 1991.

\bibitem{kalisch2007estimating}
M.~Kalisch and P.~B{\"u}hlmann.
\newblock Estimating high-dimensional directed acyclic graphs with the
  {PC}-algorithm.
\newblock \emph{Journal of Machine Learning Research}, 8:613--636, 2007.

\bibitem{lloyd1982least}
S.~P.~Lloyd.
\newblock Least squares quantization in {PCM}.
\newblock \emph{IEEE Transactions on Information Theory}, 28(2):129--137, 1982.

\bibitem{hoeffding1963probability}
W.~Hoeffding.
\newblock Probability inequalities for sums of bounded random variables.
\newblock \emph{Journal of the American Statistical Association},
  58(301):13--30, 1963.

\bibitem{zheng2024causal}
X.~Zheng, C.~Huang, Y.~Chen, and K.~Zhang.
\newblock \texttt{causal-learn}: Causal discovery in {Python}.
\newblock \emph{Journal of Machine Learning Research}, 25(60):1--8, 2024.

\bibitem{kuhn1955hungarian}
H.~W.~Kuhn.
\newblock The {H}ungarian method for the assignment problem.
\newblock \emph{Naval Research Logistics Quarterly}, 2(1--2):83--97, 1955.

\bibitem{maathuis2009estimating}
M.~H.~Maathuis, M.~Kalisch, and P.~B{\"u}hlmann.
\newblock Estimating high-dimensional intervention effects from observational
  data.
\newblock \emph{The Annals of Statistics}, 37(6A):3133--3164, 2009.

\end{thebibliography}


% ─────────────────────────────────────────────────────────────────────────────
%  Appendix
% ─────────────────────────────────────────────────────────────────────────────
\appendix

\section{Proofs of Supporting Lemmas}\label{app:proofs}

This appendix provides detailed proofs for the key lemmas that support the
main theorems in Section~5.

\subsection{Proof of Lemma~1: $\sqrt{\mathrm{JSD}}$ Satisfies the Triangle Inequality}

\begin{lemma}[Triangle inequality for $\sqrt{\mathrm{JSD}}$]
\label{lem:jsd-triangle}
For any three probability distributions $P$, $Q$, $R$ defined on a common
measurable space $(\Omega,\mathcal{F})$,
\[
  \sqrt{\mathrm{JSD}(P\|R)}
  \;\le\;
  \sqrt{\mathrm{JSD}(P\|Q)} + \sqrt{\mathrm{JSD}(Q\|R)}.
\]
\end{lemma}

\begin{proof}
The Jensen--Shannon divergence is defined as
$\mathrm{JSD}(P\|Q) = \tfrac{1}{2}\mathrm{KL}(P\|M) +
\tfrac{1}{2}\mathrm{KL}(Q\|M)$, where $M = \tfrac{1}{2}(P+Q)$.
Endres and Schindelin~\cite{endres2003new} proved that $\sqrt{\mathrm{JSD}}$
is a metric on the space of probability distributions.  Their proof proceeds
by showing that $\mathrm{JSD}$ is the square of a Hilbertian metric:
there exists an embedding $\phi$ of distributions into a Hilbert space
$\mathcal{H}$ such that
$\mathrm{JSD}(P\|Q) = \|\phi(P)-\phi(Q)\|_{\mathcal{H}}^2$.

Given this embedding, the triangle inequality follows from the triangle
inequality in $\mathcal{H}$:
\begin{align*}
\sqrt{\mathrm{JSD}(P\|R)}
&= \|\phi(P)-\phi(R)\|_{\mathcal{H}} \\
&\le \|\phi(P)-\phi(Q)\|_{\mathcal{H}} + \|\phi(Q)-\phi(R)\|_{\mathcal{H}} \\
&= \sqrt{\mathrm{JSD}(P\|Q)} + \sqrt{\mathrm{JSD}(Q\|R)}.
\end{align*}

For the specific construction in our framework, where $P$, $Q$, $R$ are
conditional distributions $P(X_i\mid\mathrm{Pa}_i)$ parametrised by
regression coefficients and residual variances, the JSD has the closed-form
expression (for multivariate Gaussians $P=\mathcal{N}(\mu_1,\Sigma_1)$,
$Q=\mathcal{N}(\mu_2,\Sigma_2)$):
\begin{align*}
\mathrm{JSD}(P\|Q) &= \tfrac{1}{2}\mathrm{KL}(P\|M)
  + \tfrac{1}{2}\mathrm{KL}(Q\|M), \\
M &= \mathcal{N}\!\left(\tfrac{\mu_1+\mu_2}{2},\,
  \tfrac{\Sigma_1+\Sigma_2}{2} + \tfrac{(\mu_1-\mu_2)(\mu_1-\mu_2)^\top}{4}
  \right),
\end{align*}
where the KL divergences are computed via the standard formula for
Gaussians.  The triangle inequality holds regardless of the parametric
family, as it depends only on the Hilbertian embedding property.
\end{proof}

\subsection{Proof of Lemma~2: Concentration of Empirical $\sqrt{\mathrm{JSD}}$}

\begin{lemma}[Concentration bound]\label{lem:jsd-concentration}
Let $\hat P_n$ denote the empirical conditional distribution estimated from
$n$ i.i.d.\ samples, and let $P$ be the true conditional.  Assume
$d_{\max}$ parents and bounded fourth moments.  Then for any $\delta>0$:
\[
  \Pr\!\left[\left|\sqrt{\mathrm{JSD}(\hat P_n\|Q)}
  - \sqrt{\mathrm{JSD}(P\|Q)}\right|
  > \varepsilon\right]
  \;\le\; 2\exp\!\left(-\frac{n\varepsilon^2}
  {C\,d_{\max}\,\sigma_{\max}^4}\right),
\]
where $C$ is a universal constant and $\sigma_{\max}^2$ is the maximum
marginal variance.
\end{lemma}

\begin{proof}
We decompose the proof into three steps.

\textbf{Step 1: Stability of JSD under parameter perturbation.}
For Gaussian conditionals $P=\mathcal{N}(\beta^\top x,\sigma^2)$, the JSD
between $P$ and $Q=\mathcal{N}(\beta'^\top x,{\sigma'}^2)$ depends
smoothly on $(\beta,\sigma^2,\beta',{\sigma'}^2)$.  Specifically,
$\mathrm{JSD}$ is Lipschitz in the parameters with constant
$L = O(d_{\max}\,\sigma_{\max}^{-2})$:
\[
  |\mathrm{JSD}(P_1\|Q) - \mathrm{JSD}(P_2\|Q)|
  \le L\,\|(\beta_1,\sigma_1^2)-(\beta_2,\sigma_2^2)\|_2.
\]

\textbf{Step 2: Concentration of regression estimates.}
Standard results for OLS regression give that with probability $\ge 1-\delta$:
\[
  \|\hat\beta - \beta\|_2 \le \sigma\sqrt{\frac{d_{\max}+\log(1/\delta)}{n}},
  \qquad
  |\hat\sigma^2 - \sigma^2| \le \sigma^2\sqrt{\frac{2\log(1/\delta)}{n}}.
\]

\textbf{Step 3: Combining via the square-root.}
Since $|\sqrt{a}-\sqrt{b}| \le |a-b|/\sqrt{\min(a,b)}$ for $a,b>0$,
and JSD is lower-bounded by $\Omega(\|\beta-\beta'\|_2^2/\sigma_{\max}^2)$
for well-separated mechanisms, we obtain the stated bound by combining
Steps~1--2 and applying the union bound over the $d_{\max}+1$ estimated
parameters.
\end{proof}

\subsection{Proof of Lemma~3: Alignment Cost Consistency}

\begin{lemma}[Alignment consistency]\label{lem:alignment-consistency}
Let $\hat\pi_n$ denote the alignment computed by CADA from empirical
mechanism distances, and let $\pi^*$ denote the optimal alignment under
true mechanism distances.  Then
$\Pr[\hat\pi_n = \pi^*] \ge 1 - |V|^2\exp(-cn/d_{\max})$
for a constant $c>0$ depending on the minimum mechanism-distance gap
$\Delta_{\min} = \min_{\pi\ne\pi^*}(d_{\mathrm{DAG}}^{\pi}-d_{\mathrm{DAG}}^{\pi^*})$.
\end{lemma}

\begin{proof}
By Lemma~\ref{lem:jsd-concentration}, each pairwise mechanism distance
is estimated to within $\varepsilon$ with probability
$\ge 1-2\exp(-cn\varepsilon^2/d_{\max})$.  There are at most $|V|^2$
pairwise distances.  By a union bound, with probability
$\ge 1-2|V|^2\exp(-cn\varepsilon^2/d_{\max})$, \emph{all} pairwise
distances are within $\varepsilon$ of their true values.

When $\varepsilon < \Delta_{\min}/(2|V|)$, the total alignment cost
perturbation is at most $|V|\varepsilon < \Delta_{\min}/2$, which is
insufficient to change the optimal alignment.  Setting
$\varepsilon = \Delta_{\min}/(2|V|)$ yields the result.
\end{proof}


\section{Algorithm Parameter Sensitivity Analysis}\label{app:sensitivity}

This appendix analyses the sensitivity of CPA to its key hyperparameters.

\subsection{CADA Alignment Weights}

The CADA alignment cost function uses three weights:
$\lambda_{\text{struct}}$ (structural disagreement),
$\lambda_{\text{mech}}$ (mechanism distance), and
$\lambda_{\text{miss}}$ (missing-variable penalty).

\paragraph{Sensitivity protocol.}
We fix CSAC Medium ($p=15$, $K=8$, $n=500$) and vary each weight
independently while holding the other two constant:
\begin{itemize}
  \item $\lambda_{\text{struct}} \in \{0.1, 0.3, 0.5, 0.7, 1.0\}$
        (default: $0.5$),
  \item $\lambda_{\text{mech}} \in \{0.1, 0.3, 0.5, 0.7, 1.0\}$
        (default: $0.3$),
  \item $\lambda_{\text{miss}} \in \{0.5, 1.0, 2.0, 5.0, 10.0\}$
        (default: $2.0$).
\end{itemize}

\paragraph{Expected behaviour.}
The alignment score should be relatively insensitive to moderate perturbations
of the weights (NAS variation $< 0.05$ within a factor of $2\times$ of the
default) but degrade significantly when $\lambda_{\text{mech}} \to 0$
(ignoring mechanism distances reduces alignment to pure graph matching)
or $\lambda_{\text{miss}} \to 0$ (removing the missing-variable penalty
allows degenerate alignments that map all variables to $\bot$).

\subsection{Plasticity Classification Thresholds}

The plasticity spectrum (Definition~8) uses four thresholds:
$\tau_S$ (structural plasticity), $\tau_P$ (parametric plasticity),
$\tau_E$ (emergence), and $\tau_{CS}$ (context sensitivity).

\paragraph{Sensitivity protocol.}
On FSVP Medium, vary each threshold over a grid of 10 values spanning
$[0.5\tau_{\text{default}},\,2.0\tau_{\text{default}}]$ and record
the macro-averaged F1 score.

\paragraph{Expected behaviour.}
The F1 score should exhibit a clear peak near the default threshold value,
with gradual degradation as thresholds move away from the optimum.
If the peak is broad (F1 variation $<0.03$ within $\pm 20\%$ of the default),
the classification is robust to threshold choice.  If the peak is sharp,
adaptive threshold selection (e.g., via cross-validation over context
subsets) would be warranted.

\subsection{QD Search Parameters}

Key parameters of Algorithm~3 (CD-QD) include:
\begin{itemize}
  \item Number of CVT cells $L \in \{100, 500, 1000, 2000, 5000\}$
        (default: $1000$),
  \item Curiosity weight $\mu \in \{0, 0.1, 0.3, 0.5, 0.7, 1.0\}$
        (default: $0.5$),
  \item Population size $|\text{pop}| \in \{50, 100, 200, 500\}$
        (default: $100$),
  \item Mutation rate $p_m \in \{0.05, 0.10, 0.20, 0.30, 0.50\}$
        (default: $0.10$).
\end{itemize}

\paragraph{Expected behaviour.}
Coverage should increase with $L$ up to a saturation point (diminishing
returns beyond $L \approx 2000$).  The curiosity weight $\mu$ should show
a non-monotonic relationship with archive quality: $\mu=0$ (no curiosity)
under-explores; $\mu=1$ (pure curiosity) sacrifices quality for coverage.
The default $\mu=0.5$ is expected to balance exploration and exploitation.

\subsection{Certificate Generation Parameters}

\begin{itemize}
  \item Stability selection rounds $R_{\text{ss}} \in \{50, 100, 200, 500\}$
        (default: $100$),
  \item Subsample fraction $f_{\text{ss}} \in \{0.3, 0.5, 0.7\}$
        (default: $0.5$),
  \item Bootstrap replicates for parametric certificates
        $B_{\text{param}} \in \{100, 200, 500, 1000\}$
        (default: $500$),
  \item Confidence level $\alpha \in \{0.01, 0.05, 0.10\}$
        (default: $0.05$).
\end{itemize}

\paragraph{Expected behaviour.}
Certificate calibration (ECE) should improve with $R_{\text{ss}}$ and
$B_{\text{param}}$, with diminishing returns.  The subsample fraction
$f_{\text{ss}}=0.5$ is recommended by Meinshausen and
B\"uhlmann~\cite{meinshausen2010stability}; deviations should not
significantly affect calibration unless the sample size is very small
($n<200$).


\section{Extended Complexity Analysis}\label{app:complexity}

\subsection{Phase~1: Foundation}

\paragraph{Causal discovery.}
For each of $K$ contexts, GES runs in
$O(p^2\cdot n + p^{d_{\max}+2})$~\cite{chickering2002optimal}, where
the first term is for computing sufficient statistics and the second is
for greedy equivalence search with bounded indegree.  Total:
$O(K(p^2 n + p^{d_{\max}+2}))$.

\paragraph{CADA alignment.}
For each of $\binom{K}{2}$ pairs:
\begin{enumerate}
  \item Anchor propagation: $O(p)$ (linear scan of shared names).
  \item MB candidate generation: $O(p \cdot d_{\max}^2)$
        (for each unanchored variable, generate candidates based on
        Markov-blanket overlap).
  \item CI-fingerprint scoring: $O(u \cdot d_{\max} \cdot n)$
        (for $u$ unanchored variables, compute conditional independence
        fingerprints, each requiring $O(n)$ for partial correlation).
  \item Hungarian matching: $O(u^3)$, where $u \le p$ is the number
        of unanchored variables.
  \item Edge classification: $O(p \cdot d_{\max})$.
  \item Score computation: $O(p)$.
\end{enumerate}
Total per pair: $O(p\cdot d_{\max}^2 + u^3 + u\cdot d_{\max}\cdot n)$.
Total for all pairs: $O(K^2(p\cdot d_{\max}^2 + u^3 + p\cdot d_{\max}\cdot n))$.

In the practical regime ($p=50$, $d_{\max}=4$, $K=10$, $n=500$, $u\le 10$):
$\binom{10}{2}=45$ pairs, each requiring
$O(50\cdot 16 + 1000 + 50\cdot 4\cdot 500) = O(10^5)$ operations.
Total: $\approx 4.5\times 10^6$ operations ($< 1$\,second).

\paragraph{PDC descriptors.}
For each of $|\bar V|$ variables, compute $\binom{K}{2}$ pairwise
$\sqrt{\mathrm{JSD}}$ values (each $O(d_{\max}^2)$ for Gaussians):
$O(|\bar V|\cdot K^2\cdot d_{\max}^2)$.

\subsection{Phase~2: Exploration}

Each CD-QD evaluation requires: (i)~extracting a context/variable subset
($O(K\cdot p)$), (ii)~computing descriptors on the subset ($O(K^2 d_{\max}^2)$),
(iii)~evaluating quality ($O(p)$).
Total for $T=10^4$ evaluations:
$O(T\cdot(K\cdot p + K^2 d_{\max}^2))$.

CVT tessellation via Lloyd's algorithm:
$O(I\cdot S\cdot L\cdot D)$ where $I=50$ iterations,
$S=10^5$ samples, $L=1000$ cells, $D=4$ dimensions.
Total: $O(2\times 10^{10})$ comparisons, but each is $O(D)=O(1)$, so
$\approx 10$\,seconds.

\subsection{Phase~3: Validation}

\paragraph{TPD.}
For each variable, PELT on a sequence of length $K-1$:
$O(K^2)$ with pruning~\cite{killick2012optimal}.
Permutation validation: $B_{\text{perm}}\cdot K^2$ per variable.
Total: $O(|\bar V|\cdot B_{\text{perm}}\cdot K^2)$.

\paragraph{RCG.}
Structural certificates via stability selection:
$R_{\text{ss}}$ rounds, each requiring a half-sample DAG estimation
($O(p^2 n/2 + p^{d_{\max}+2})$) per context.
Total: $O(n_{\text{inv}}\cdot R_{\text{ss}}\cdot K\cdot(p^2 n + p^{d_{\max}+2}))$.

Parametric certificates via bootstrap:
$B_{\text{param}}$ resamples, each requiring $O(d_{\max}^2\cdot n)$
per mechanism.
Total: $O(n_{\text{inv}}\cdot K^2\cdot B_{\text{param}}\cdot d_{\max}^2\cdot n)$.

\subsection{Overall Scaling}

\begin{center}
\small
\begin{tabular}{lccc}
\toprule
Component & Scaling & $p\!=\!50,K\!=\!10$ & $p\!=\!100,K\!=\!10$ \\
\midrule
Discovery & $O(K p^2 n)$ & $1.3\times 10^8$ & $5\times 10^8$ \\
Alignment & $O(K^2 p d n)$ & $4.5\times 10^6$ & $2\times 10^7$ \\
Descriptors & $O(p K^2 d^2)$ & $4\times 10^4$ & $1.6\times 10^5$ \\
QD search & $O(T K p)$ & $5\times 10^6$ & $10^7$ \\
Tipping pts & $O(p B K^2)$ & $5\times 10^6$ & $2\times 10^7$ \\
Certificates & $O(n_{\text{inv}} R K p^2 n)$ & $7.5\times 10^{10}$ & $6\times 10^{11}$ \\
\bottomrule
\end{tabular}
\end{center}

The certificate generation step dominates, as shown in
Section~\ref{sec:budget}.  With stability selection (100 rounds rather than
500 bootstrap DAG re-estimations per $K^2$ pairs), the cost is reduced by
a factor of $\sim 25K$ relative to full bootstrap.


\section{Additional Benchmark Specifications}\label{app:benchmarks}

\subsection{Semi-Synthetic Perturbation Protocol}

For all semi-synthetic benchmarks (Sachs, Asia, Alarm, Insurance), we use
the following standardised perturbation protocol to generate
multi-context variants:

\begin{enumerate}
\item \textbf{Edge classification.}
      Partition the edge set $E$ of the reference network into:
      \begin{itemize}
        \item Invariant ($\lfloor f_I |E| \rfloor$ edges,
              $f_I \in \{0.5, 0.6, 0.7\}$): parameters unchanged.
        \item Parametrically plastic ($\lfloor f_P |E| \rfloor$ edges,
              $f_P = 1 - f_I - f_D - f_E$): for continuous networks,
              multiply regression coefficients by a factor
              $\gamma_k \sim \mathrm{Uniform}([0.3, 2.0])$ per context;
              for discrete networks, perturb CPT entries by additive noise
              $\eta \sim \mathrm{Uniform}([-\Delta, \Delta])$, clipped to
              $[0,1]$ and renormalised, with $\Delta \in \{0.05, 0.10, 0.15\}$.
        \item Disappearing ($\lfloor f_D |E| \rfloor$ edges,
              $f_D \in \{0.05, 0.10, 0.15\}$): removed in contexts
              $k > k^*$, with $k^* \sim \mathrm{Uniform}(\{2,\dots,K-1\})$.
        \item Emergent ($\lfloor f_E |E| \rfloor$ edges,
              $f_E \in \{0.05, 0.10\}$): new edges not in $E$, added in
              contexts $k > k^*$ while preserving acyclicity.
      \end{itemize}

\item \textbf{Sample generation.}
      For each context, sample $n$ observations from the perturbed
      network using ancestral sampling (for discrete networks) or the
      linear SEM (for Gaussian networks).

\item \textbf{Ground-truth labels.}
      Record per-edge plasticity labels based on the perturbation
      assignment.  For disappearing and emergent edges, record the
      transition context index.
\end{enumerate}

\subsection{Benchmark Network Statistics}

\begin{center}
\small
\begin{tabular}{lcccccc}
\toprule
Network & $p$ & $|E|$ & Max indegree & Avg.\ degree & Type & $K$ \\
\midrule
Sachs     &  11 &  17 & 3 & 1.55 & Continuous &  5 \\
Asia      &   8 &   8 & 2 & 1.00 & Discrete   &  6 \\
Alarm     &  37 &  46 & 4 & 1.24 & Discrete   &  4 \\
Insurance &  27 &  52 & 3 & 1.93 & Discrete   &  8 \\
\bottomrule
\end{tabular}
\end{center}

\subsection{SHD Perturbation Experiment Details}

To validate Theorem~8 (structural stability), we inject controlled
errors into estimated DAGs and measure the resulting perturbation of
plasticity descriptors.

\paragraph{Perturbation procedure.}
Given a true DAG $G$ with $|E|$ edges:
\begin{enumerate}
  \item For each $s \in \{1, 2, 5, 10\}$, generate 50 perturbed DAGs
        $\tilde G^{(j)}_s$ by applying $s$ random operations, each
        chosen uniformly from $\{\text{add edge}, \text{delete edge},
        \text{reverse edge}\}$ while maintaining acyclicity.
  \item Run the full CPA classification pipeline on $\tilde G^{(j)}_s$.
  \item Compute the descriptor perturbation
        $\Delta_j = \|\tilde\psi^{(j)} - \psi^*\|_\infty$.
  \item Compare $\Delta_j$ against the theoretical bound
        $f(s, n, d_{\max}) = C' \cdot s \cdot d_{\max} / \sqrt{n}$
        from Theorem~8.
\end{enumerate}

\paragraph{Expected outcome.}
For $\ge 95\%$ of trials, $\Delta_j \le f(s, n, d_{\max})$.
The empirical scaling of $\Delta_j$ with $s$ should be approximately
linear, confirming the theoretical prediction.

\subsection{Computational Cost of Benchmark Suite}

\begin{center}
\small
\begin{tabular}{lcccc}
\toprule
Benchmark & Replications & Time/rep & Total & Cores \\
\midrule
FSVP (4 levels) & $4 \times 50$ & 1--20\,min & 10\,h & 8 \\
CSAC (4 levels) & $4 \times 50$ & 2--30\,min & 15\,h & 8 \\
TPS (4 levels)  & $4 \times 50$ & 5--45\,min & 25\,h & 8 \\
LC (4 levels)   & $4 \times 30$ & 5--30\,min & 10\,h & 8 \\
Semi-synthetic  & $4 \times 50$ & 1--10\,min &  5\,h & 8 \\
SHD perturb.    & $4 \times 50$ & 2--15\,min &  8\,h & 8 \\
Ablations (6)   & $6 \times 30$ & 5--30\,min & 20\,h & 8 \\
\midrule
\textbf{Total}  & & & \textbf{$\sim$93\,h} & 8 \\
\bottomrule
\end{tabular}
\end{center}

The full benchmark suite requires approximately 93 CPU-hours
($\sim 12$\,hours wall-clock with 8-core parallelism), fitting comfortably
within a single-machine execution constraint.

\end{document}